\documentclass[11pt]{article}
\usepackage[margin=1in]{geometry}
\usepackage[T1]{fontenc}
\usepackage[utf8]{inputenc}
\usepackage[french]{babel}
\usepackage{amsmath,amssymb,amsthm}
\usepackage[colorlinks=true,linkcolor=blue,urlcolor=blue]{hyperref}
\newtheorem{problem}{Problème}
\newenvironment{refs}{\par\small\noindent\emph{Références :}\begin{itemize}\setlength\itemsep{0pt}}{\end{itemize}\normalsize}
\title{Hors de la Caverne \\ \Large La salle des collégiens}
\author{}
\date{}
\begin{document}
\maketitle
\noindent\emph{Des problèmes, pas de réponses.} Chaque problème ci-dessous est posé
sans solution. Les références indiquent où trouver les connaissances nécessaires.
\medskip


\section*{Algèbre abstraite}

\begin{problem}[{Trois entiers consécutifs --- Collège}]
\textbf{ALG-44.} \textbf{Problème.} Des entiers consécutifs sont des entiers qui se suivent, comme 7, 8, 9. Dans
tout ce problème, il s'agit de démontrer, pas de vérifier sur des exemples.
\begin{enumerate}
\item[(a)] Démontrez que la somme de trois entiers consécutifs est toujours divisible par 3. On appellera
$ n $ le premier des trois, et l'on prendra soin de dire pourquoi le raisonnement vaut aussi
lorsque $ n $ est négatif.
\item[(b)] Démontrez que la somme de cinq entiers consécutifs est toujours divisible par 5, puis que la
somme de quatre entiers consécutifs n'est jamais divisible par 4.
\item[(c)] Déterminez, en le démontrant, tous les entiers $ k \ge 2 $ tels que la somme de $ k $
entiers consécutifs quelconques soit toujours divisible par $ k $.
\end{enumerate}

\end{problem}

\noindent Le passage de \og{} ça marche sur mes exemples \fg{} à \og{} ça marche toujours, et voici
pourquoi \fg{} est exactement ce que la lettre $ n $ rend possible : une seule ligne de calcul
littéral remplace une infinité de vérifications. C'est le geste fondateur de l'algèbre, et il est
plus ancien qu'on ne croit --- al-Khwârizmî, à Bagdad vers 820, écrit encore ses inconnues en toutes
lettres dans le traité dont le titre a donné le mot \textit{algèbre}, et il faudra attendre Viète à
la fin du XVIe siècle pour que les lettres deviennent le langage ordinaire du calcul.
La partie (c) montre qu'une belle règle a souvent une frontière nette, et que la trouver est plus
instructif que la règle elle-même.

\begin{refs}
\item Contexte : Calcul algébrique --- Wikipédia (\url{https://fr.wikipedia.org/wiki/Calcul_alg%C3%A9brique})
\item Contexte : Parité (arithmétique) --- Wikipédia (\url{https://fr.wikipedia.org/wiki/Parit%C3%A9_(arithm%C3%A9tique)})
\item Contexte : Al-Khwârizmî --- Wikipédia (\url{https://fr.wikipedia.org/wiki/Al-Khw%C3%A2rizm%C3%AE})
\end{refs}

\bigskip

\begin{problem}[{\og{} Pense à un nombre\dots{} \fg{} --- Collège}]
\textbf{ALG-45.} \textbf{Problème.} Voici un tour de magie numérique. Pensez à un nombre ; ajoutez 5 ; multipliez
le résultat par 2 ; retirez 4 ; divisez par 2 ; enfin, retirez le nombre auquel vous aviez pensé.
Le magicien annonce alors que le résultat est 3, quel que soit le nombre choisi au départ.
\begin{enumerate}
\item[(a)] Vérifiez le tour sur trois nombres de départ différents, dont un nombre négatif et un nombre
décimal.
\item[(b)] En appelant $ x $ le nombre pensé, écrivez l'expression obtenue après chaque étape, puis
démontrez que le résultat final ne dépend pas de $ x $. Précisez à quel moment exactement le
$ x $ disparaît, et par quelle propriété du calcul.
\item[(c)] Inventez un tour du même genre dont la réponse annoncée est toujours 7, et démontrez qu'il
fonctionne. Inventez-en ensuite un dont la réponse est le double du nombre pensé, et démontrez-le
aussi.
\end{enumerate}

\end{problem}

\noindent Ces tours circulent dans les cours de récréation depuis des siècles, et ils sont
de l'algèbre déguisée : le magicien ne devine rien, il a construit une expression qui se simplifie
en une constante. Martin Gardner, qui a tenu la chronique de récréations mathématiques du
\textit{Scientific American} de 1956 à 1981, en a recueilli et analysé des centaines, précisément
parce qu'expliquer le truc et démontrer une identité sont le même travail. La partie (c) est la
plus formatrice : construire un tour, c'est résoudre le problème à l'envers, et l'on ne peut le
faire qu'en comprenant le mécanisme.

\begin{refs}
\item Contexte : Calcul algébrique --- Wikipédia (\url{https://fr.wikipedia.org/wiki/Calcul_alg%C3%A9brique})
\item Contexte : Martin Gardner --- Wikipédia (\url{https://fr.wikipedia.org/wiki/Martin_Gardner})
\end{refs}

\bigskip

\begin{problem}[{Les identités remarquables démontrées par des aires --- Collège}]
\textbf{ALG-46.} \textbf{Problème.} Dans tout ce problème, $ a $ et $ b $ désignent des longueurs, donc des
nombres strictement positifs, et chaque égalité doit être établie par un découpage de figure et un
calcul d'aires --- pas par un développement algébrique.
\begin{enumerate}
\item[(a)] Un carré de côté $ a+b $ est découpé en quatre morceaux par deux traits. Démontrez, en
calculant l'aire totale de deux façons, que
\[ (a+b)^2 = a^2 + 2ab + b^2 . \]
\item[(b)] On suppose $ 0 < b < a $. Démontrez de même que $ (a-b)^2 = a^2 - 2ab + b^2 $ en
retirant deux bandes à un carré de côté $ a $, et expliquez pourquoi le petit carré de côté
$ b $ est d'abord retiré deux fois, puis rendu une fois.
\item[(c)] Démontrez que $ a^2 - b^2 = (a-b)(a+b) $ en découpant en deux morceaux la figure obtenue en
ôtant d'un carré de côté $ a $ un carré de côté $ b $ placé dans un coin, puis en recollant ces
deux morceaux pour former un rectangle. Précisez les dimensions du rectangle obtenu.
\end{enumerate}

\end{problem}

\noindent Le livre II des \textit{Éléments} d'Euclide, écrit vers 300 av. J.-C., ne contient
aucune formule : ses propositions 4 et 5 sont exactement les identités ci-dessus, énoncées comme
des théorèmes sur des rectangles et des carrés, et démontrées par découpage. Les historiens
appellent cela l'algèbre géométrique des Grecs --- ils n'avaient pas de notation littérale, mais ils
avaient les figures, et les figures suffisaient. La méthode de complétion du carré d'al-Khwârizmî,
neuf siècles plus tard, est encore un dessin ; comprendre ces identités comme des aires plutôt que
comme des recettes à retenir change durablement la façon dont on les manipule.

\begin{refs}
\item Contexte : Identité remarquable --- Wikipédia (\url{https://fr.wikipedia.org/wiki/Identit%C3%A9_remarquable})
\item Contexte : Éléments (Euclide) --- Wikipédia (\url{https://fr.wikipedia.org/wiki/%C3%89l%C3%A9ments_(Euclide)})
\end{refs}

\bigskip

\begin{problem}[{Une hausse puis une baisse --- Collège, short}]
\textbf{ALG-47.} \textbf{Problème.} Un prix $ p $ subit une hausse de 20 \%, puis le nouveau prix
subit une baisse de 20 \% ; démontrez que l'on ne retrouve jamais le prix de départ, sauf si
$ p = 0 $, et déterminez la variation totale exprimée en pourcentage. Déterminez ensuite, en le
justifiant, le pourcentage de baisse qui annulerait exactement une hausse de 20 \%.
\end{problem}

\noindent C'est l'un des rares résultats de collège que l'on peut utiliser toute sa vie :
les pourcentages successifs se multiplient et ne s'additionnent pas, ce qui explique pourquoi une
remise sur une remise n'est pas la somme des deux, et pourquoi un placement qui perd 50 \% puis
gagne 50 \% a bel et bien perdu. La démonstration tient en une ligne dès qu'on écrit la hausse comme
une multiplication par un coefficient plutôt que comme une addition --- c'est le changement de point
de vue qui fait tout le travail.

\begin{refs}
\item Contexte : Pourcentage --- Wikipédia (\url{https://fr.wikipedia.org/wiki/Pourcentage})
\item Contexte : Proportionnalité --- Wikipédia (\url{https://fr.wikipedia.org/wiki/Proportionnalit%C3%A9})
\end{refs}

\bigskip

\begin{problem}[{Le carré des nombres qui finissent par 5 --- Collège, short}]
\textbf{ALG-48.} \textbf{Problème.} On observe que $ 25^2 = 625 $, $ 35^2 = 1225 $ et
$ 45^2 = 2025 $ : chaque résultat se termine par 25, et ce qui précède ces deux chiffres est le
produit du chiffre des dizaines par l'entier suivant. Démontrez cette règle pour tout entier de la
forme $ 10a + 5 $ avec $ a $ entier naturel, et déterminez si elle vaut encore quand $ a $ a
lui-même plusieurs chiffres, comme dans $ 115^2 $.
\end{problem}

\noindent Les calculateurs mentaux professionnels du XIXe siècle vivaient de
règles de ce genre, et les manuels de calcul rapide en sont pleins ; ce qui les distingue d'un
simple truc, c'est qu'une identité remarquable les démontre en deux lignes. Une fois la
démonstration écrite, la règle cesse d'être à retenir : on la reconstruit. C'est la différence
entre savoir un tour et savoir pourquoi il ne peut pas échouer.

\begin{refs}
\item Contexte : Identité remarquable --- Wikipédia (\url{https://fr.wikipedia.org/wiki/Identit%C3%A9_remarquable})
\item Contexte : Calcul mental --- Wikipédia (\url{https://fr.wikipedia.org/wiki/Calcul_mental})
\end{refs}

\bigskip

\begin{problem}[{Le partage, posé en équation --- Collège}]
\textbf{ALG-49.} \textbf{Problème.} Poser une équation, c'est nommer d'une lettre la quantité cherchée, puis
traduire chaque phrase de l'énoncé en une écriture mathématique.
\begin{enumerate}
\item[(a)] Trois personnes se partagent 120 euros : la deuxième reçoit le double de la première, et la
troisième reçoit 15 euros de plus que la deuxième. Déterminez les trois parts en posant une
équation d'inconnue la part de la première, et démontrez que la solution est unique.
\item[(b)] Problème 40 du papyrus Rhind, copié en Égypte vers 1650 av. J.-C. : partager 100 pains entre
cinq hommes de telle sorte que les cinq parts se suivent avec un écart constant, et que le septième
de la somme des trois plus grandes parts soit égal à la somme des deux plus petites. Traduisez
l'énoncé en calcul littéral et résolvez-le en justifiant chaque étape.
\item[(c)] Construisez un énoncé de partage du même type qui n'admet aucune solution acceptable, et
démontrez qu'il n'en admet aucune. Expliquez ensuite ce que le passage par une inconnue apporte de
plus qu'un tâtonnement, même bien organisé.
\end{enumerate}

\end{problem}

\noindent Le papyrus Rhind, copié par le scribe Ahmès d'après un original plus ancien
encore et conservé aujourd'hui au British Museum, est l'un des plus vieux textes mathématiques que
nous ayons ; ses problèmes de partage de pains et de bière sont des problèmes d'équation avant que
le mot n'existe. Les Égyptiens les résolvaient par la méthode dite de fausse position --- on essaie
une valeur commode, on regarde de combien on se trompe, puis on corrige proportionnellement --- et
comparer cette méthode à la vôtre est instructif : elle fonctionne exactement parce que la relation
est linéaire.

\begin{refs}
\item Contexte : Papyrus Rhind --- Wikipédia (\url{https://fr.wikipedia.org/wiki/Papyrus_Rhind})
\item Contexte : Équation du premier degré --- Wikipédia (\url{https://fr.wikipedia.org/wiki/%C3%89quation_du_premier_degr%C3%A9})
\end{refs}

\bigskip

\begin{problem}[{L'épitaphe de Diophante --- Collège}]
\textbf{ALG-50.} \textbf{Problème.} Une épigramme de l'\textit{Anthologie grecque} décrit ainsi la vie de Diophante :
son enfance occupa le sixième de sa vie ; un douzième de plus, et sa joue se couvrit de barbe ;
après un septième de plus il se maria ; cinq ans plus tard naquit un fils, qui vécut la moitié de
la vie de son père ; le père survécut quatre ans à son enfant, puis mourut.
\begin{enumerate}
\item[(a)] En appelant $ x $ l'âge de Diophante à sa mort, traduisez chacune des six indications de
l'épigramme en une expression littérale.
\item[(b)] Écrivez l'équation que traduit la phrase \og{} ces durées mises bout à bout composent toute sa
vie \fg{}, puis résolvez-la en justifiant chaque étape.
\item[(c)] Vérifiez que toutes les durées obtenues sont bien des nombres entiers d'années, et expliquez
pourquoi l'auteur de l'épigramme ne pouvait pas choisir les fractions au hasard. Remplacez ensuite
le sixième par un cinquième et déterminez, en le justifiant, si l'énigme modifiée admet encore une
solution entière.
\end{enumerate}

\end{problem}

\noindent Diophante d'Alexandrie, auteur des \textit{Arithmétiques}, a vécu au
IIIe siècle de notre ère --- les dates ne sont pas sûres, et c'est en partie pourquoi
cette épigramme, transmise par l'\textit{Anthologie grecque} et probablement composée bien après sa
mort, a tant circulé. Le charme du problème tient à ce que la poésie et l'algèbre y coïncident : la
narration impose sa contrainte à chaque vers, et l'équation est le poème traduit. C'est aussi
l'ancêtre direct de tous les problèmes d'âges, et Diophante a donné son nom aux équations dont on
ne cherche que les solutions entières.

\begin{refs}
\item Contexte : Diophante d'Alexandrie --- Wikipédia (\url{https://fr.wikipedia.org/wiki/Diophante_d%27Alexandrie})
\item Contexte : Anthologie grecque --- Wikipédia (\url{https://fr.wikipedia.org/wiki/Anthologie_grecque})
\end{refs}

\bigskip

\begin{problem}[{Les poules et les lapins --- Collège}]
\textbf{ALG-51.} \textbf{Problème.} Dans une basse-cour on ne compte que des poules, à deux pattes, et des lapins,
à quatre pattes. On dénombre 35 têtes et 94 pattes.
\begin{enumerate}
\item[(a)] Déterminez le nombre de poules et le nombre de lapins par un tâtonnement organisé dans un
tableau, et démontrez que votre tableau ne peut pas avoir laissé échapper une autre solution.
\item[(b)] Reprenez la question en n'introduisant qu'une seule inconnue et une seule équation du premier
degré, et démontrez que la solution est unique.
\item[(c)] On compte maintenant 35 têtes et 95 pattes. Démontrez qu'aucune basse-cour ne peut donner ce
résultat. Déterminez ensuite, en le justifiant, tous les nombres de pattes possibles pour 35
têtes.
\end{enumerate}

\end{problem}

\noindent Le problème apparaît sous cette forme exacte --- 35 têtes, 94 pattes --- dans le
\textit{Sunzi Suanjing}, manuel chinois du IIIe au Ve siècle, avec des faisans
et des lapins ; il a voyagé jusqu'au Japon, où on l'appelle le calcul des grues et des tortues, et
il figure encore aujourd'hui dans les manuels du monde entier. Sa longévité vient de ce qu'il se
laisse attaquer de trois façons également honnêtes : par tableau, par une inconnue, ou par
l'astuce dite \og{} faites lever les lapins sur leurs pattes arrière \fg{}. La partie (c) est la plus
mathématique des trois, car elle demande de démontrer une impossibilité, ce qu'aucun tâtonnement
ne saura jamais faire.

\begin{refs}
\item Contexte : Sunzi Suanjing --- Wikipédia (\url{https://fr.wikipedia.org/wiki/Sunzi_Suanjing})
\item Contexte : Système d'équations linéaires --- Wikipédia (\url{https://fr.wikipedia.org/wiki/Syst%C3%A8me_d%27%C3%A9quations_lin%C3%A9aires})
\end{refs}

\bigskip

\begin{problem}[{La somme de Gauss --- Collège}]
\textbf{ALG-52.} \textbf{Problème.} On s'intéresse aux sommes d'entiers qui se suivent, et à la façon de les
calculer sans les additionner un par un.
\begin{enumerate}
\item[(a)] Calculez $ 1 + 2 + 3 + \cdots + 100 $ en appariant le premier terme avec le dernier, le
deuxième avec l'avant-dernier, et ainsi de suite. Démontrez que tous les couples ainsi formés ont
la même somme, et justifiez le nombre de couples obtenus.
\item[(b)] Démontrez que pour tout entier $ n \ge 1 $,
\[ 1 + 2 + 3 + \cdots + n = \frac{n(n+1)}{2} , \]
en écrivant la somme deux fois, une fois à l'endroit et une fois à l'envers, puis en ajoutant les
deux lignes colonne par colonne. Expliquez pourquoi cette démonstration évite d'avoir à distinguer
le cas où $ n $ est pair de celui où $ n $ est impair.
\item[(c)] Démontrez que la somme des $ n $ premiers nombres impairs vaut $ n^2 $, et illustrez le
résultat par un dessin de carrés emboîtés en expliquant ce que représente chaque équerre
ajoutée.
\end{enumerate}

\end{problem}

\noindent L'anecdote est célèbre : le jeune Gauss, à l'école primaire de Brunswick, aurait
mis quelques secondes à additionner une longue suite d'entiers que le maître croyait bonne à
occuper la classe une heure. Elle nous vient du mémorial que son ami Sartorius von Waltershausen
publie en 1856, l'année suivant la mort de Gauss ; la version précise \og{} de 1 à 100 \fg{} est un
embellissement plus tardif, et l'anecdote a probablement été polie par le temps. Le procédé, lui,
est bien plus vieux que Gauss --- on le trouve chez les pythagoriciens sous forme de cailloux
disposés en triangle --- et la partie (c) donne l'une des plus belles démonstrations sans mots de
toutes les mathématiques.

\begin{refs}
\item Contexte : Suite arithmétique --- Wikipédia (\url{https://fr.wikipedia.org/wiki/Suite_arithm%C3%A9tique})
\item Contexte : Carl Friedrich Gauss --- Wikipédia (\url{https://fr.wikipedia.org/wiki/Carl_Friedrich_Gauss})
\item Contexte : Nombre triangulaire --- Wikipédia (\url{https://fr.wikipedia.org/wiki/Nombre_triangulaire})
\end{refs}

\bigskip

\begin{problem}[{Échelles, agrandissements et vitesses --- Collège}]
\textbf{ALG-53.} \textbf{Problème.} Trois situations où l'on croit reconnaître de la proportionnalité, et où il faut
regarder de près.
\begin{enumerate}
\item[(a)] Sur une carte à l'échelle $ 1 : 25\,000 $, deux villages sont distants de 8,4 cm. Déterminez
la distance réelle en kilomètres et justifiez le calcul par un tableau de proportionnalité, en
disant ce que représente le coefficient employé.
\item[(b)] Une maquette de bateau est construite à l'échelle $ 1 : 50 $. Démontrez que l'aire de sa
voilure est 2 500 fois plus petite que celle du bateau réel, et son volume 125 000 fois plus petit.
Expliquez pourquoi les trois facteurs de réduction sont différents alors que l'échelle est
unique.
\item[(c)] Un cycliste parcourt 27 km en 1 h 12 min. Un second roule 20 \% plus vite que le premier.
Déterminez le temps mis par le second, en le justifiant, puis démontrez que le temps de parcours
n'est pas proportionnel à la vitesse mais lui est inversement proportionnel.
\end{enumerate}

\end{problem}

\noindent La partie (b) est la raison pour laquelle on ne peut pas fabriquer un insecte
géant ni un éléphant miniature : quand une longueur est multipliée par $ k $, les aires le sont
par $ k^2 $ et les volumes par $ k^3 $, si bien que la masse croît beaucoup plus vite que la
section des os qui doit la porter. Galilée l'expose dès 1638 dans les \textit{Discours concernant deux
sciences nouvelles}, et J. B. S. Haldane en tire en 1926 un essai resté fameux, \textit{On Being the
Right Size}. La partie (c) traque l'erreur la plus répandue en proportionnalité : croire que
deux grandeurs liées varient forcément dans le même sens et dans le même rapport.

\begin{refs}
\item Contexte : Proportionnalité --- Wikipédia (\url{https://fr.wikipedia.org/wiki/Proportionnalit%C3%A9})
\item Contexte : Échelle (proportion) --- Wikipédia (\url{https://fr.wikipedia.org/wiki/%C3%89chelle_(proportion)})
\end{refs}

\bigskip

\begin{problem}[{Multiplier autour d'un carré --- Collège, short}]
\textbf{ALG-54.} \textbf{Problème.} Pour calculer $ 19 \times 21 $ de tête on remarque que les
deux facteurs encadrent 20 ; démontrez que $ (n-1)(n+1) = n^2 - 1 $ pour tout nombre $ n $, puis
établissez la règle plus générale donnant $ (n-k)(n+k) $. Servez-vous-en pour justifier un calcul
mental de $ 97 \times 103 $, puis déterminez si tout produit de deux entiers peut se ramener de
cette façon à une différence de deux carrés.
\end{problem}

\noindent Cette identité est la plus utile des trois identités remarquables, et elle est
aussi la clef d'une méthode de factorisation que Fermat employait au XVIIe siècle pour
casser les grands nombres : chercher à écrire $ N $ comme une différence de deux carrés, ce qui
le factorise aussitôt. La dernière question du problème est donc historiquement chargée --- la
réponse dépend de la parité, et elle explique pourquoi la méthode de Fermat s'applique bien aux
nombres impairs.

\begin{refs}
\item Contexte : Identité remarquable --- Wikipédia (\url{https://fr.wikipedia.org/wiki/Identit%C3%A9_remarquable})
\item Contexte : Factorisation --- Wikipédia (\url{https://fr.wikipedia.org/wiki/Factorisation})
\end{refs}

\bigskip

\begin{problem}[{Le nombre retourné --- Collège, short}]
\textbf{ALG-55.} \textbf{Problème.} Prenez un nombre entier de deux chiffres, retournez-le, et
soustrayez le plus petit des deux nombres du plus grand ; démontrez que le résultat est toujours un
multiple de 9, et déterminez exactement quels multiples de 9 peuvent être obtenus. Reprenez ensuite
la question avec un nombre de trois chiffres et son retourné.
\end{problem}

\noindent L'écriture décimale est une convention, et l'algèbre sert ici à la rendre
visible : dès qu'on écrit un nombre de deux chiffres sous la forme $ 10a + b $, le résultat
devient une évidence de calcul là où il paraissait un tour de passe-passe. Le même déguisement
explique le critère de divisibilité par 9, la preuve par neuf des comptables médiévaux, et une
bonne moitié des tours de cartes numériques. Il vaut la peine d'essayer aussi en base autre que
dix : le 9 y est remplacé par un autre nombre, et savoir lequel, c'est avoir compris la
démonstration.

\begin{refs}
\item Contexte : Critère de divisibilité --- Wikipédia (\url{https://fr.wikipedia.org/wiki/Crit%C3%A8re_de_divisibilit%C3%A9})
\item Contexte : Système décimal --- Wikipédia (\url{https://fr.wikipedia.org/wiki/Syst%C3%A8me_d%C3%A9cimal})
\end{refs}

\bigskip


\section*{Théorie des nombres}

\begin{problem}[{Les nombres 9, 99, 999 et la règle de trois --- Collège}]
\textbf{NT-54.} \textbf{Problème.} Tout entier à quatre chiffres s'écrit $ n = 1000a + 100b + 10c + d $, où
$ a, b, c, d $ sont ses chiffres.
\begin{enumerate}
\item[(a)] Démontrez que $ 9 $, $ 99 $ et $ 999 $ sont tous les trois divisibles par $ 3 $ et par
$ 9 $.
\item[(b)] Vérifiez l'égalité
\[ n = (999a + 99b + 9c) + (a + b + c + d) , \]
puis démontrez que $ n $ est divisible par $ 3 $ exactement lorsque $ a + b + c + d $ l'est.
Démontrez de même le critère par $ 9 $.
\item[(c)] Expliquez comment votre argument s'adapte à un nombre ayant un nombre quelconque de chiffres, et
dites clairement quelle propriété des nombres $ 9, 99, 999, 9999, \ldots $ le fait fonctionner.
\item[(d)] Le même raisonnement donnerait-il un critère de divisibilité par $ 7 $ fondé sur la somme des
chiffres ? Exhibez deux entiers ayant la même somme de chiffres dont l'un est divisible par $ 7 $ et
l'autre non, puis dites précisément à quel endroit l'argument de (b) cesse de fonctionner.
\end{enumerate}

\end{problem}

\noindent La règle est enseignée à tous les écoliers et sa raison n'est presque jamais
montrée, alors qu'elle tient en une ligne dès qu'on regarde $ 999 $ au lieu de $ 1000 $. Le
procédé est très ancien : les arithméticiens indiens puis les mathématiciens de langue arabe s'en
servaient comme contrôle de comptable sous le nom de \og{} preuve par neuf \fg{}, et il voyage vers l'Europe
avec la numération décimale de position, celle-là même qu'al-Khwârizmî expose au IXe siècle
et que Fibonacci fait connaître en 1202. La partie (d) est là pour une raison de fond : un critère
n'est pas un fait sur les nombres, c'est un fait sur \textit{l'écriture} des nombres en base dix, et il
s'effondre dès que le diviseur ne divise pas $ 10 - 1 $. Le problème NT-37, plus bas, redémontre la
même chose dans la langue des congruences --- comparez les deux rédactions.

\begin{refs}
\item Contexte : Critère de divisibilité --- Wikipédia (\url{https://fr.wikipedia.org/wiki/Crit%C3%A8re_de_divisibilit%C3%A9})
\item Contexte : Système décimal --- Wikipédia (\url{https://fr.wikipedia.org/wiki/Syst%C3%A8me_d%C3%A9cimal})
\item Contexte : Al-Khwârizmî --- Wikipédia (\url{https://fr.wikipedia.org/wiki/Al-Khw%C3%A2rizm%C3%AE})
\end{refs}

\bigskip

\begin{problem}[{Le crible d'Ératosthène jusqu'à 100 --- Collège}]
\textbf{NT-55.} \textbf{Problème.} Écrivez les entiers de $ 2 $ à $ 100 $. Barrez les multiples de $ 2 $ autres
que $ 2 $, puis les multiples de $ 3 $ autres que $ 3 $, et ainsi de suite ; ce qui reste est la
liste des nombres premiers inférieurs à $ 100 $.
\begin{enumerate}
\item[(a)] Menez le crible à son terme et donnez la liste obtenue.
\item[(b)] Démontrez que tout entier $ n \ge 2 $ qui n'est pas premier possède un diviseur premier
$ p $ vérifiant $ p \times p \le n $. En déduire, avec démonstration, qu'il suffit de barrer les
multiples de $ 2, 3, 5 $ et $ 7 $ pour que tout ce qui survit jusqu'à $ 100 $ soit premier.
\item[(c)] Quand on barre les multiples de $ p $, on peut commencer directement à $ p \times p $.
Démontrez que rien n'est perdu, c'est-à-dire que tout multiple de $ p $ strictement compris entre
$ p $ et $ p \times p $ a déjà été barré à une étape précédente.
\item[(d)] Dans votre liste, $ 3, 5, 7 $ sont trois nombres premiers en progression de deux en deux.
Démontrez qu'il n'y en a pas d'autre : si $ p, p+2, p+4 $ sont tous premiers, alors $ p = 3 $.
(Regardez les restes de $ p $ dans la division par $ 3 $.)
\end{enumerate}

\end{problem}

\noindent Ératosthène de Cyrène, qui dirigeait la bibliothèque d'Alexandrie au
IIIe siècle av. J.-C. et qui mesura la circonférence de la Terre à partir de deux ombres,
a laissé son nom à ce procédé de tri --- le mot \og{} crible \fg{} dit bien ce qu'il fait : il laisse tomber les
composés et retient les premiers. Ce qui mérite l'attention n'est pas la mécanique, que n'importe
quel élève exécute, mais les deux justifications demandées en (b) et (c) : elles expliquent pourquoi
on peut s'arrêter, et c'est exactement ce qui sépare une recette d'un théorème. La question (d) est un
premier contact avec une idée qui traverse toute la théorie des nombres : parmi trois entiers espacés
de deux, l'un est toujours multiple de $ 3 $, qu'on le veuille ou non. Les couples $ (p, p+2) $,
eux, semblent se poursuivre indéfiniment --- c'est la conjecture des nombres premiers jumeaux, toujours
ouverte aujourd'hui.

\begin{refs}
\item Contexte : Crible d'Ératosthène --- Wikipédia (\url{https://fr.wikipedia.org/wiki/Crible_d%27%C3%89ratosth%C3%A8ne})
\item Contexte : Nombre premier --- Wikipédia (\url{https://fr.wikipedia.org/wiki/Nombre_premier})
\item Contexte : Nombres premiers jumeaux --- Wikipédia (\url{https://fr.wikipedia.org/wiki/Nombres_premiers_jumeaux})
\end{refs}

\bigskip

\begin{problem}[{Pair, impair --- Collège, short}]
\textbf{NT-56.} \textbf{Problème.} Démontrez que la somme de deux nombres impairs est toujours
paire, et que leur produit est toujours impair. Votre rédaction doit utiliser l'écriture
$ 2k + 1 $ d'un nombre impair et ne s'appuyer sur aucun exemple.
\end{problem}

\noindent C'est probablement la première vraie démonstration qu'on puisse écrire en entier,
et elle contient déjà tout : on remplace \og{} j'ai essayé beaucoup de nombres \fg{} par une lettre qui les
désigne tous à la fois. La contrainte de ne citer aucun exemple n'est pas une coquetterie --- vérifier
mille cas ne démontre rien, et le jour où l'on rencontre un énoncé faux à partir du millionième
entier, on comprend pourquoi.

\begin{refs}
\item Contexte : Parité (arithmétique) --- Wikipédia (\url{https://fr.wikipedia.org/wiki/Parit%C3%A9_(arithm%C3%A9tique)})
\item Voir aussi : L'art de la démonstration (\url{proofs.html})
\end{refs}

\bigskip

\begin{problem}[{La somme du petit Gauss --- Collège}]
\textbf{NT-57.} \textbf{Problème.} On veut calculer $ 1 + 2 + 3 + \cdots + n $ sans poser l'addition.
\begin{enumerate}
\item[(a)] Pour $ n = 100 $, groupez les termes deux par deux en partant des extrémités : $ 1 + 100 $,
$ 2 + 99 $, et ainsi de suite. Expliquez pourquoi toutes ces sommes sont égales, combien il y a de
groupes, et concluez.
\item[(b)] Démontrez la formule générale
\[ 1 + 2 + \cdots + n = \frac{n(n+1)}{2} \]
en écrivant la somme deux fois, la seconde à l'envers, et en additionnant les deux lignes terme à
terme. Traitez aussi le cas où $ n $ est impair, où le groupement de (a) laisse un terme seul.
\item[(c)] Démontrez que $ \frac{n(n+1)}{2} $ est toujours un entier, c'est-à-dire que le produit de deux
entiers consécutifs est toujours pair.
\item[(d)] Démontrez que la somme des $ n $ premiers nombres impairs vaut $ n \times n $. Une figure
faite de carrés emboîtés suffit, à condition que vous expliquiez précisément ce que chaque équerre de
la figure représente.
\end{enumerate}

\end{problem}

\noindent L'anecdote est célèbre et sa forme exacte est incertaine : un maître d'école de
Brunswick aurait donné à sa classe l'addition des cent premiers entiers pour avoir la paix, et Carl
Friedrich Gauss, âgé de sept ou huit ans, aurait posé son ardoise presque aussitôt. Ce qui est sûr,
c'est que le procédé lui est bien antérieur --- les nombres de la forme $ n(n+1)/2 $ sont les
\textit{nombres triangulaires} des Grecs, comptés en disposant des cailloux en triangle, et Nicomaque de
Gérase les décrit vers l'an 100. La partie (d) montre l'autre manière de démontrer en arithmétique :
non par le calcul mais par la figure, ce qu'on appelle une preuve sans mots. Les deux voies sont
également rigoureuses --- à condition qu'on sache dire ce que la figure démontre.

\begin{refs}
\item Contexte : Nombre triangulaire --- Wikipédia (\url{https://fr.wikipedia.org/wiki/Nombre_triangulaire})
\item Contexte : Carl Friedrich Gauss --- Wikipédia (\url{https://fr.wikipedia.org/wiki/Carl_Friedrich_Gauss})
\item Contexte : Preuve sans mots --- Wikipédia (\url{https://fr.wikipedia.org/wiki/Preuve_sans_mots})
\end{refs}

\bigskip

\begin{problem}[{L'algorithme d'Euclide par soustractions --- Collège}]
\textbf{NT-58.} \textbf{Problème.} Soient $ a $ et $ b $ deux entiers strictement positifs avec
$ a \ge b $. On admet que $ a $ et $ b $ possèdent un plus grand diviseur commun, noté
$ \mathrm{PGCD}(a, b) $.
\begin{enumerate}
\item[(a)] Démontrez que tout diviseur commun de $ a $ et $ b $ est aussi un diviseur commun de $ b $
et $ a - b $, et réciproquement. En déduire l'égalité
\[ \mathrm{PGCD}(a, b) = \mathrm{PGCD}(b,\ a - b) . \]
\item[(b)] Appliquez cette égalité à répétition, en remplaçant chaque fois le plus grand des deux nombres
par leur différence, pour calculer $ \mathrm{PGCD}(1071, 462) $. Écrivez toutes les étapes.
\item[(c)] Expliquez pourquoi ce procédé s'arrête toujours, quels que soient les entiers de départ.
\item[(d)] Rendez la fraction $ \dfrac{1071}{462} $ irréductible, puis démontrez que la fraction obtenue ne
peut plus être simplifiée.
\end{enumerate}

\end{problem}

\noindent C'est le plus ancien algorithme non trivial encore en usage quotidien : il occupe
les propositions 1 et 2 du livre VII des \textit{Éléments} d'Euclide, vers 300 av. J.-C., où il est
énoncé pour des longueurs et non pour des nombres --- on retranche la plus courte de la plus longue
jusqu'à ce que les deux coïncident, et l'on obtient ainsi leur commune mesure. La version par
soustractions demandée ici est exactement celle d'Euclide ; la version par restes, plus rapide, est
la même idée abrégée, et le problème NT-03 en tire l'identité de Bézout puis toute l'arithmétique
élémentaire. Notez que la question (c) est la vraie : un procédé qui ne s'arrêterait pas ne
calculerait rien.

\begin{refs}
\item Contexte : Algorithme d'Euclide --- Wikipédia (\url{https://fr.wikipedia.org/wiki/Algorithme_d%27Euclide})
\item Contexte : Plus grand commun diviseur --- Wikipédia (\url{https://fr.wikipedia.org/wiki/Plus_grand_commun_diviseur})
\item Contexte : Nombres premiers entre eux --- Wikipédia (\url{https://fr.wikipedia.org/wiki/Nombres_premiers_entre_eux})
\end{refs}

\bigskip

\begin{problem}[{Deux phares --- Collège, short}]
\textbf{NT-59.} \textbf{Problème.} Un phare s'allume toutes les $ 12 $ secondes, un autre toutes
les $ 18 $ secondes ; à l'instant $ 0 $, ils s'allument ensemble. Démontrez qu'ils se rallument
ensemble exactement toutes les $ 36 $ secondes --- c'est-à-dire que $ 36 $ convient, et qu'aucune
durée strictement plus courte ne convient.
\end{problem}

\noindent Le nombre $ 36 $ est le plus petit commun multiple de $ 12 $ et de $ 18 $, et
l'énoncé demande les deux moitiés qu'on oublie toujours : qu'il marche, et qu'il soit le plus petit.
La seconde moitié se démontre en divisant un éventuel instant commun par $ 36 $ et en regardant le
reste --- un raisonnement qu'on retrouvera identique, mot pour mot, pour l'ordre d'un élément dans un
groupe.

\begin{refs}
\item Contexte : Plus petit commun multiple --- Wikipédia (\url{https://fr.wikipedia.org/wiki/Plus_petit_commun_multiple})
\end{refs}

\bigskip

\begin{problem}[{Le dernier chiffre des puissances de 2 --- Collège, short}]
\textbf{NT-60.} \textbf{Problème.} Démontrez que le dernier chiffre de $ 2^n $, pour
$ n \ge 1 $, ne prend que quatre valeurs et qu'il se répète avec une période de quatre --- en
justifiant que le dernier chiffre de $ 2^{n+1} $ ne dépend que du dernier chiffre de $ 2^n $. Déterminez
alors le dernier chiffre de $ 2^{2026} $ sans calculer ce nombre.
\end{problem}

\noindent Le dernier chiffre d'un nombre est son reste dans la division par $ 10 $, et
l'observation décisive est que ce reste se propage : pour connaître le chiffre suivant, il suffit du
chiffre courant, jamais du nombre entier. Une quantité qui ne dépend que d'elle-même à l'étape
précédente finit forcément par se répéter, puisqu'il n'y a que dix chiffres possibles --- c'est le
premier exemple de périodicité de toute l'arithmétique, et le problème NT-36 en donne la version
lycée.

\begin{refs}
\item Contexte : Puissance d'un nombre --- Wikipédia (\url{https://fr.wikipedia.org/wiki/Puissance_d%27un_nombre})
\item Contexte : Système décimal --- Wikipédia (\url{https://fr.wikipedia.org/wiki/Syst%C3%A8me_d%C3%A9cimal})
\end{refs}

\bigskip

\begin{problem}[{6 et 28, les nombres parfaits --- Collège}]
\textbf{NT-61.} \textbf{Problème.} Un entier strictement positif est dit \textit{parfait} lorsqu'il est égal à la somme
de ses diviseurs autres que lui-même. Ainsi $ 6 = 1 + 2 + 3 $.
\begin{enumerate}
\item[(a)] Dressez la liste complète des diviseurs de $ 28 $, justifiez que votre liste ne peut en oublier
aucun, et vérifiez que $ 28 $ est parfait.
\item[(b)] Démontrez de même que $ 496 = 2^4 \times 31 $ est parfait. (Commencez par démontrer que ses seuls
diviseurs sont les $ 2^k $ pour $ 0 \le k \le 4 $ et leurs produits par $ 31 $.)
\item[(c)] Démontrez qu'un nombre premier n'est jamais parfait. Démontrez également qu'aucune puissance de
$ 2 $ n'est parfaite.
\item[(d)] Les quatre nombres parfaits connus des Grecs sont $ 6, 28, 496, 8128 $. Écrivez chacun sous la
forme $ 2^{k-1} \times (2^k - 1) $, constatez que le second facteur est premier à chaque fois, et
énoncez la règle de fabrication qu'Euclide en a tirée. Vous n'avez pas à la démontrer : dites
seulement, avec précision, ce qu'elle affirme et ce qu'elle n'affirme pas.
\end{enumerate}

\end{problem}

\noindent La proposition 36 du livre IX des \textit{Éléments} donne la recette : si
$ 2^k - 1 $ est premier, alors $ 2^{k-1}(2^k - 1) $ est parfait. Euler démontrera vingt siècles
plus tard que \textit{tout} nombre parfait pair est de cette forme --- la réciproque exacte, et l'un des
plus beaux allers-retours de l'histoire des mathématiques. Ces nombres ont porté une lourde charge
symbolique : saint Augustin écrit que Dieu créa le monde en six jours parce que six est parfait, et
Nicomaque de Gérase les rangeait déjà parmi les nombres \og{} ni excessifs ni défectueux \fg{}. Ce que la
question (d) doit vous faire sentir, c'est la place d'un trou : personne n'a jamais trouvé de nombre
parfait \textit{impair}, et personne n'a jamais démontré qu'il n'en existe pas. La question est ouverte
depuis plus de deux mille ans.

\begin{refs}
\item Contexte : Nombre parfait --- Wikipédia (\url{https://fr.wikipedia.org/wiki/Nombre_parfait})
\item Contexte : Nombre de Mersenne premier --- Wikipédia (\url{https://fr.wikipedia.org/wiki/Nombre_de_Mersenne_premier})
\item Contexte : Nicomaque de Gérase --- Wikipédia (\url{https://fr.wikipedia.org/wiki/Nicomaque_de_G%C3%A9rase})
\item Voir aussi : Problèmes ouverts (\url{open-problems.html})
\end{refs}

\bigskip

\begin{problem}[{Autant de nombres pairs que d'entiers --- Collège}]
\textbf{NT-62.} \textbf{Problème.} L'intuition dit qu'il y a \og{} deux fois moins \fg{} de nombres pairs que d'entiers,
puisqu'un entier sur deux est pair. On va voir que cette phrase n'a pas de sens tant qu'on n'a pas
dit ce que compter veut dire.
\begin{enumerate}
\item[(a)] À chaque entier $ n \ge 1 $ on associe le nombre pair $ 2n $. Démontrez que deux entiers
différents reçoivent des nombres pairs différents, et que tout nombre pair strictement positif est
atteint par exactement un entier.
\item[(b)] Un hôtel possède une infinité de chambres numérotées $ 1, 2, 3, \ldots $, toutes occupées. Un
voyageur se présente. Expliquez comment le loger sans mettre personne dehors, et démontrez que dans
votre solution aucun client ne se retrouve sans chambre et qu'aucune chambre n'accueille deux
clients.
\item[(c)] Reprenez la question (b) avec une infinité de nouveaux voyageurs arrivant en même temps, et
justifiez votre réponse avec la même rigueur.
\item[(d)] Expliquez avec vos propres mots pourquoi la phrase \og{} il y a moins de nombres pairs que d'entiers \fg{}
et la phrase \og{} il y en a autant \fg{} peuvent être vraies toutes les deux, selon ce que l'on décide
d'appeler \og{} autant \fg{}.
\end{enumerate}

\end{problem}

\noindent Galilée avait déjà buté sur ce paradoxe en 1638, en remarquant qu'il y a autant de
carrés parfaits que d'entiers alors que les carrés se raréfient : il en concluait qu'il ne faut pas
appliquer aux collections infinies les mots \og{} plus \fg{}, \og{} moins \fg{} et \og{} autant \fg{}. Georg Cantor a tranché
autrement à partir de 1874 : il a gardé les mots et changé leur définition --- deux collections ont
autant d'éléments lorsqu'on peut les apparier une à une, sans reste ni répétition. C'est exactement ce
que vous démontrez en (a). L'hôtel est une image due à David Hilbert, qui s'en servait vers 1924 dans
ses cours pour faire sentir que l'infini n'obéit pas aux habitudes du fini. Rien ici ne dépasse le
programme du collège, et pourtant vous êtes en train de manipuler l'objet qui a divisé les
mathématiciens pendant trente ans.

\begin{refs}
\item Contexte : Hôtel de Hilbert --- Wikipédia (\url{https://fr.wikipedia.org/wiki/H%C3%B4tel_de_Hilbert})
\item Contexte : Ensemble dénombrable --- Wikipédia (\url{https://fr.wikipedia.org/wiki/Ensemble_d%C3%A9nombrable})
\item Contexte : Georg Cantor --- Wikipédia (\url{https://fr.wikipedia.org/wiki/Georg_Cantor})
\item Voir aussi : Logique, théorie des ensembles et fondements (\url{pure-foundations.html})
\end{refs}

\bigskip

\begin{problem}[{Les fractions des scribes égyptiens --- Collège}]
\textbf{NT-63.} \textbf{Problème.} Les scribes égyptiens n'écrivaient que des fractions de numérateur $ 1 $, et
exprimaient toute autre fraction comme une somme de telles fractions, toutes différentes. Par exemple
$ \frac{3}{4} = \frac{1}{2} + \frac{1}{4} $.
\begin{enumerate}
\item[(a)] Écrivez $ \frac{2}{5} $ et $ \frac{3}{7} $ comme sommes de fractions de numérateur $ 1 $
deux à deux distinctes, et vérifiez vos égalités par le calcul.
\item[(b)] Démontrez que pour tout entier $ n \ge 1 $,
\[ \frac{1}{n} = \frac{1}{n+1} + \frac{1}{n(n+1)} . \]
\item[(c)] Déduisez de (b) qu'une même fraction admet une infinité d'écritures égyptiennes différentes, en
expliquant précisément comment on passe de l'une à la suivante.
\item[(d)] Un scribe doit partager $ 5 $ pains entre $ 8 $ personnes, chaque personne devant recevoir la
même chose et aucun pain ne devant être coupé en plus de morceaux que nécessaire. Proposez un partage,
et démontrez que chacun reçoit bien $ \frac{5}{8} $ de pain.
\end{enumerate}

\end{problem}

\noindent Le papyrus Rhind, copié vers 1550 av. J.-C. par le scribe Ahmès d'après un original
plus ancien de deux siècles, s'ouvre sur une table donnant $ 2/n $ comme somme de fractions
unitaires pour tous les $ n $ impairs de $ 3 $ à $ 101 $ : c'est le plus vieux document
mathématique substantiel qui nous soit parvenu, et il est essentiellement fait de ce problème. On
ignore encore pourquoi les Égyptiens s'astreignaient à cette contrainte ; l'explication la plus
souvent avancée est celle du partage effectif --- couper des pains en parts égales visibles est plus
simple avec des moitiés et des quarts qu'avec des cinq-huitièmes. L'identité de la question (b) est
aussi la clé de l'algorithme glouton que Fibonacci décrit en 1202 : prendre
à chaque étape la plus grande fraction unitaire qui tient, et recommencer. Qu'il se termine toujours
est un joli théorème ; qu'il donne la décomposition la plus courte est faux.

\begin{refs}
\item Contexte : Fraction égyptienne --- Wikipédia (\url{https://fr.wikipedia.org/wiki/Fraction_%C3%A9gyptienne})
\item Contexte : Papyrus Rhind --- Wikipédia (\url{https://fr.wikipedia.org/wiki/Papyrus_Rhind})
\item Contexte : Ahmès --- Wikipédia (\url{https://fr.wikipedia.org/wiki/Ahm%C3%A8s})
\end{refs}

\bigskip

\begin{problem}[{Les lapins de Fibonacci --- Collège}]
\textbf{NT-64.} \textbf{Problème.} On pose $ F_1 = 1 $, $ F_2 = 1 $, et chaque terme suivant est la somme des
deux précédents : $ F_{n+2} = F_{n+1} + F_n $.
\begin{enumerate}
\item[(a)] Calculez $ F_1 $ jusqu'à $ F_{12} $. Expliquez ensuite en quoi cette règle traduit exactement
l'énoncé de Fibonacci : un couple de lapins devient fécond au bout d'un mois et engendre ensuite un
couple par mois, aucun ne meurt.
\item[(b)] Écrivez la suite des parités (pair ou impair) des douze premiers termes. Démontrez que $ F_n $
est pair exactement lorsque $ n $ est un multiple de $ 3 $ --- il suffit de raisonner sur les
parités, jamais sur les valeurs.
\item[(c)] Vérifiez sur plusieurs valeurs de $ n $, puis démontrez, l'égalité
\[ F_1 + F_2 + \cdots + F_n = F_{n+2} - 1 . \]
(Remarquez que $ F_k = F_{k+2} - F_{k+1} $ et additionnez ces égalités : presque tout se
simplifie.)
\item[(d)] Calculez les quotients $ F_{n+1} / F_n $ pour $ n $ allant de $ 1 $ à $ 11 $. Décrivez ce
que vous observez et expliquez pourquoi cette observation, si nette soit-elle, ne constitue pas une
démonstration.
\end{enumerate}

\end{problem}

\noindent Leonardo de Pise, dit Fibonacci, pose le problème des lapins dans le \textit{Liber
abaci} de 1202 --- un livre écrit d'abord pour convaincre les marchands italiens d'abandonner les
chiffres romains, et dont la postérité a surtout retenu un seul exercice. La suite lui est
pourtant bien antérieure : les prosodistes indiens Pingala, Virahanka et Hemachandra la connaissaient
en comptant les manières de composer un vers avec des syllabes brèves et longues, plusieurs siècles
plus tôt. Le quotient de la question (d) s'approche du nombre d'or $ (1 + \sqrt{5})/2 $, ce que vous
verrez démontré plus tard ; la vraie leçon de cette question est ailleurs, dans l'écart entre
\og{} je vois que \fg{} et \og{} je démontre que \fg{}. Aucun élève n'y échappe et aucun mathématicien non plus.

\begin{refs}
\item Contexte : Suite de Fibonacci --- Wikipédia (\url{https://fr.wikipedia.org/wiki/Suite_de_Fibonacci})
\item Contexte : Liber abaci --- Wikipédia (\url{https://fr.wikipedia.org/wiki/Liber_abaci})
\item Contexte : Nombre d'or --- Wikipédia (\url{https://fr.wikipedia.org/wiki/Nombre_d%27or})
\end{refs}

\bigskip

\begin{problem}[{La divisibilité par 11 --- Collège, short}]
\textbf{NT-65.} \textbf{Problème.} Démontrez qu'un entier à quatre chiffres
$ n = 1000a + 100b + 10c + d $ est divisible par $ 11 $ exactement lorsque $ d - c + b - a $
l'est. (Partez de $ 11 \times 91 = 1001 $, $ 11 \times 9 = 99 $ et $ 11 \times 1 = 11 $.)
\end{problem}

\noindent Le critère par $ 3 $ marchait parce que $ 9, 99, 999 $ sont des multiples de
$ 3 $ ; celui-ci marche parce que $ 11, 1001, 100001 $ le sont de $ 11 $, tandis que
$ 99, 9999 $ le sont aussi --- d'où l'alternance des signes, qui est la petite surprise de l'affaire.
C'est l'occasion de voir qu'un critère de divisibilité n'est jamais gratuit : il dépend de la base
d'écriture, et il faut à chaque fois trouver quels nombres formés de $ 9 $ et de $ 1 $ le diviseur
accepte de diviser.

\begin{refs}
\item Contexte : Critère de divisibilité --- Wikipédia (\url{https://fr.wikipedia.org/wiki/Crit%C3%A8re_de_divisibilit%C3%A9})
\end{refs}

\bigskip


\section*{Géométrie}

\begin{problem}[{Pourquoi les angles d'un triangle font 180$^\circ$ --- Collège}]
\textbf{GEO-41.} \textbf{Problème.} Soit $ ABC $ un triangle. Tracez la droite passant par $ A $ et parallèle à
$ (BC) $.
\begin{enumerate}
\item[(a)] En utilisant les angles alternes-internes formés par cette parallèle avec $ (AB) $ puis avec
$ (AC) $, démontrez que
\[ \widehat{A} + \widehat{B} + \widehat{C} = 180^\circ . \]
Nommez explicitement, à chaque étape, la propriété du cours que vous employez.
\item[(b)] Déduisez-en que la somme des angles d'un quadrilatère non croisé vaut $ 360^\circ $, puis que
celle d'un polygone convexe à $ n $ côtés vaut $ (n-2) \times 180^\circ $. Justifiez le découpage
que vous utilisez.
\item[(c)] Sur une sphère, on peut tracer un triangle dont les trois angles sont droits : partez du pôle
Nord, descendez jusqu'à l'équateur, parcourez un quart d'équateur, remontez. Sa somme d'angles vaut
$ 270^\circ $. Expliquez précisément quelle étape de votre démonstration de (a) devient fausse dans
ce cas.
\end{enumerate}

\end{problem}

\noindent C'est la proposition 32 du livre I des \textit{Éléments} d'Euclide, et la question (c)
en désigne le point sensible : la démonstration ne tient que parce qu'on a pu tracer \textit{une seule}
parallèle à $ (BC) $ passant par $ A $. Cet énoncé est le cinquième postulat d'Euclide, qu'il
avait lui-même isolé comme suspect et que vingt siècles de mathématiciens ont tenté en vain de
déduire des quatre autres. Gauss, Bolyai et Lobatchevski ont fini par comprendre, dans les années
1820-1830, qu'on ne le pouvait pas : il existe des géométries cohérentes où il est faux, et où la
somme des angles d'un triangle n'est jamais $ 180^\circ $. Autrement dit, ce résultat de sixième est
équivalent au postulat le plus discuté de l'histoire des mathématiques --- le problème GEO-06, plus bas,
reprend l'affaire.

\begin{refs}
\item Contexte : Somme des angles d'un triangle --- Wikipédia (\url{https://fr.wikipedia.org/wiki/Somme_des_angles_d%27un_triangle})
\item Contexte : Angles alternes-internes --- Wikipédia (\url{https://fr.wikipedia.org/wiki/Angles_alternes-internes})
\item Contexte : Géométrie sphérique --- Wikipédia (\url{https://fr.wikipedia.org/wiki/G%C3%A9om%C3%A9trie_sph%C3%A9rique})
\end{refs}

\bigskip

\begin{problem}[{Pythagore par découpage --- Collège}]
\textbf{GEO-42.} \textbf{Problème.} Soit un triangle rectangle de côtés de l'angle droit $ a $ et $ b $, et
d'hypoténuse $ c $. Prenez deux carrés identiques de côté $ a + b $. Dans le premier, placez
quatre copies du triangle dans les quatre coins de façon à laisser au centre un quadrilatère de côté
$ c $ ; dans le second, placez les quatre mêmes copies de façon à laisser deux carrés vides, de
côtés $ a $ et $ b $.
\begin{enumerate}
\item[(a)] Démontrez que le quadrilatère central de la première figure est bien un carré. C'est le point que
tout le monde saute : il faut prouver que ses angles sont droits, en utilisant que les angles aigus du
triangle rectangle ont pour somme $ 90^\circ $.
\item[(b)] En comparant les aires des deux grands carrés, démontrez que
\[ a^2 + b^2 = c^2 . \]
\item[(c)] Le triangle de côtés $ 3, 4, 5 $ est rectangle. Cela découle-t-il du théorème que vous venez de
démontrer, ou d'un autre énoncé ? Formulez précisément l'énoncé nécessaire et dites en quoi il diffère
du précédent.
\item[(d)] Un camarade affirme : \og{} j'ai vérifié sur vingt triangles rectangles dessinés et mesurés, donc le
théorème est vrai. \fg{} Expliquez pourquoi ce n'est pas une démonstration, et ce que ses mesures
établissent réellement.
\end{enumerate}

\end{problem}

\noindent Le résultat était connu et utilisé bien avant Pythagore : la tablette babylonienne
Plimpton 322, vers 1800 av. J.-C., aligne des triplets d'entiers vérifiant la relation, et les cordes
à nœuds des arpenteurs égyptiens s'en servaient pour tracer des angles droits. La démonstration par
découpage demandée ici est essentiellement celle qui accompagne le \textit{Zhoubi Suanjing} chinois, où
la figure porte le nom de \textit{xian tu}, et on la retrouve chez Bhāskara II au XIIe siècle
en Inde, accompagnée, dit la tradition, du seul mot \og{} Regarde ! \fg{}. C'est le théorème qui compte le
plus de démonstrations différentes --- plusieurs centaines recensées, dont une due à un futur président
des États-Unis, James Garfield, en 1876. La question (c) sépare deux énoncés qu'on confond presque
toujours : le théorème et sa réciproque, et seule la réciproque permet de \textit{conclure} qu'un
triangle est rectangle.

\begin{refs}
\item Contexte : Théorème de Pythagore --- Wikipédia (\url{https://fr.wikipedia.org/wiki/Th%C3%A9or%C3%A8me_de_Pythagore})
\item Contexte : Plimpton 322 --- Wikipédia (\url{https://fr.wikipedia.org/wiki/Plimpton_322})
\item Contexte : Zhoubi Suanjing --- Wikipédia (\url{https://fr.wikipedia.org/wiki/Zhoubi_Suanjing})
\item Contexte : Preuve sans mots --- Wikipédia (\url{https://fr.wikipedia.org/wiki/Preuve_sans_mots})
\end{refs}

\bigskip

\begin{problem}[{La médiatrice, ensemble des points équidistants --- Collège}]
\textbf{GEO-43.} \textbf{Problème.} Soient $ A $ et $ B $ deux points distincts, $ I $ le milieu de
$ [AB] $, et $ (d) $ la perpendiculaire à $ (AB) $ passant par $ I $.
\begin{enumerate}
\item[(a)] Démontrez que tout point $ M $ de $ (d) $ vérifie $ MA = MB $.
\item[(b)] Démontrez la réciproque : tout point $ M $ du plan vérifiant $ MA = MB $ appartient à
$ (d) $. C'est cette moitié-là qui justifie qu'on parle de \textit{l'}ensemble des points
équidistants et non seulement d'une famille de points équidistants.
\item[(c)] Soient $ A $, $ B $, $ C $ trois points non alignés. En appliquant (a) et (b) aux
médiatrices de $ [AB] $ et de $ [BC] $, démontrez que les trois médiatrices du triangle
$ ABC $ se coupent en un même point, et qu'il existe un cercle passant par les trois sommets.
\item[(d)] Trois villages veulent creuser un puits à égale distance des trois. Démontrez qu'un tel puits
existe toujours, sauf dans un cas que vous préciserez, et expliquez géométriquement ce qui se passe
dans ce cas.
\end{enumerate}

\end{problem}

\noindent La double implication de (a) et (b) est le premier exemple, dans une scolarité, d'un
\textit{ensemble caractérisé par une propriété} --- un lieu géométrique. Le mot est ancien : les Grecs
appelaient \textit{topos} une figure définie par une condition, et la plus grande partie de la géométrie
classique consiste à reconnaître qu'une condition apparemment compliquée décrit en réalité une droite
ou un cercle. Euclide construit le cercle circonscrit à la proposition 5 du livre IV. La question (d)
est là pour l'exception : les trois médiatrices d'un triangle aplati sont parallèles, et le puits
part à l'infini --- un cas dégénéré qu'une bonne rédaction n'a pas le droit de passer sous silence.

\begin{refs}
\item Contexte : Médiatrice --- Wikipédia (\url{https://fr.wikipedia.org/wiki/M%C3%A9diatrice})
\item Contexte : Cercle circonscrit --- Wikipédia (\url{https://fr.wikipedia.org/wiki/Cercle_circonscrit})
\end{refs}

\bigskip

\begin{problem}[{Encadrer $\pi$, puis l'aire du disque --- Collège}]
\textbf{GEO-44.} \textbf{Problème.} On considère un cercle de rayon $ 1 $, donc de périmètre $ 2\pi $.
\begin{enumerate}
\item[(a)] Inscrivez dans ce cercle un carré, et circonscrivez-lui un autre carré. Calculez les deux
périmètres et démontrez l'encadrement $ 2\sqrt{2} \le \pi \le 4 $, en énonçant clairement
l'hypothèse géométrique que vous admettez sur les longueurs.
\item[(b)] Remplacez le carré inscrit par un hexagone régulier. Démontrez que son côté vaut exactement
$ 1 $ --- décomposez l'hexagone en six triangles et montrez qu'ils sont équilatéraux --- puis
déduisez-en que $ \pi \ge 3 $.
\item[(c)] Découpez maintenant le disque en un grand nombre de secteurs égaux très fins, et rangez-les tête
bêche pour former une figure proche d'un rectangle. Expliquez pourquoi les dimensions de ce
\og{} rectangle \fg{} sont proches de $ \pi $ et de $ 1 $, et pourquoi cela rend plausible que l'aire du
disque vaille $ \pi $.
\item[(d)] Dites précisément ce que l'argument de (c) ne démontre pas. Quel mot, absent du programme du
collège, faudrait-il définir pour qu'il devienne une démonstration ?
\end{enumerate}

\end{problem}

\noindent Archimède écrit vers 250 av. J.-C. un petit traité, \textit{De la mesure du cercle},
qui fait exactement cela : en partant de l'hexagone et en doublant quatre fois le nombre de côtés
jusqu'au polygone à 96 côtés, inscrit puis circonscrit, il obtient $ 223/71 < \pi < 22/7 $.
C'est le premier encadrement de l'histoire assorti d'une démonstration, et le procédé --- la méthode
d'exhaustion héritée d'Eudoxe --- est l'ancêtre direct du calcul intégral. Archimède démontre aussi, dans
le même traité, que l'aire du disque est celle d'un triangle de base le périmètre et de hauteur le
rayon, ce qui est précisément l'énoncé que le découpage de (c) fait deviner. La question (d) doit vous
mener au mot manquant : \textit{limite}. Vous le rencontrerez au lycée, et la page
Analyse réelle et théorie de la mesure (\url{pure-real-analysis.html}) montre ce qu'il a fallu de travail pour le
définir proprement --- deux mille ans après Archimède.

\begin{refs}
\item Contexte : Pi --- Wikipédia (\url{https://fr.wikipedia.org/wiki/Pi})
\item Contexte : Méthode d'exhaustion --- Wikipédia (\url{https://fr.wikipedia.org/wiki/M%C3%A9thode_d%27exhaustion})
\item Contexte : Archimède --- Wikipédia (\url{https://fr.wikipedia.org/wiki/Archim%C3%A8de})
\item Voir aussi : Analyse réelle et théorie de la mesure (\url{pure-real-analysis.html})
\end{refs}

\bigskip

\begin{problem}[{Thalès, l'ombre et la pyramide --- Collège}]
\textbf{GEO-45.} \textbf{Problème.} Un piquet vertical de $ 1{,}50 $ m projette une ombre de $ 2{,}40 $ m. Au même
instant, l'ombre d'un arbre voisin mesure $ 18 $ m.
\begin{enumerate}
\item[(a)] Faites une figure faisant apparaître la configuration de Thalès, et calculez la hauteur de
l'arbre. Justifiez chaque étape par le théorème, énoncé en entier.
\item[(b)] Cette méthode repose sur une hypothèse physique que la figure ne montre pas. Laquelle ? Expliquez
pourquoi la conclusion serait fausse si les rayons du Soleil n'étaient pas sensiblement parallèles, et
si le piquet et l'arbre n'étaient pas tous deux verticaux.
\item[(c)] La tradition rapporte que Thalès mesura ainsi la hauteur de la grande pyramide de Khéops. Le
problème est en réalité plus difficile que celui de l'arbre : le pied de la pyramide n'est pas
accessible. Expliquez quelle longueur supplémentaire il faut mesurer et pourquoi.
\item[(d)] Construisez une figure où deux triangles partagent un sommet et où les longueurs sont dans le même
rapport, mais où la conclusion du théorème est fausse parce que les droites ne sont pas parallèles.
Dites ce que cela prouve sur l'usage du théorème.
\end{enumerate}

\end{problem}

\noindent Thalès de Milet, au VIe siècle av. J.-C., est le premier personnage à qui
la tradition grecque attribue des démonstrations plutôt que des recettes --- le passage, décisif, de
\og{} cela marche \fg{} à \og{} voici pourquoi \fg{}. L'histoire de la pyramide nous vient de Diogène Laërce et de
Plutarque, huit siècles plus tard, et deux versions circulent : l'une où Thalès attend l'heure à
laquelle son ombre égale sa taille, l'autre où il utilise le rapport général, celle du présent
problème. On ne peut pas trancher, et il faut le dire ainsi. Ce qui n'est pas légendaire, c'est le
contenu du théorème : il exprime que la géométrie euclidienne est invariante par changement d'échelle,
propriété que la géométrie sphérique de GEO-41 ne possède pas --- sur une sphère, agrandir une figure la
déforme.

\begin{refs}
\item Contexte : Théorème de Thalès --- Wikipédia (\url{https://fr.wikipedia.org/wiki/Th%C3%A9or%C3%A8me_de_Thal%C3%A8s})
\item Contexte : Thalès --- Wikipédia (\url{https://fr.wikipedia.org/wiki/Thal%C3%A8s})
\item Contexte : Pyramide de Khéops --- Wikipédia (\url{https://fr.wikipedia.org/wiki/Pyramide_de_Kh%C3%A9ops})
\end{refs}

\bigskip

\begin{problem}[{Le patron d'un cube et les faces d'un dé --- Collège}]
\textbf{GEO-46.} \textbf{Problème.} Un patron de cube est une figure faite de six carrés reliés par des côtés, qui se
plie en un cube sans recouvrement.
\begin{enumerate}
\item[(a)] Trouvez le plus grand nombre possible de patrons de cube, deux patrons superposables par rotation
ou retournement étant considérés comme identiques. Il y en a exactement onze : décrivez la méthode qui
vous garantit de n'en oublier aucun.
\item[(b)] Démontrez que deux carrés d'un patron qui partagent un côté ne deviennent jamais des faces
opposées du cube.
\item[(c)] Sur un dé usuel, deux faces opposées portent des nombres dont la somme vaut $ 7 $. Prenez un
patron en croix, numérotez-en une face $ 1 $ et une face voisine $ 2 $ : montrez qu'il reste
exactement deux façons de compléter le dé, et que ces deux dés ne sont pas superposables.
\item[(d)] On double l'arête d'un cube. Démontrez que son aire est multipliée par $ 4 $ et son volume par
$ 8 $, et expliquez pourquoi les exposants $ 2 $ et $ 3 $ apparaissent ici.
\end{enumerate}

\end{problem}

\noindent Le dénombrement des onze patrons est l'un des rares problèmes de collège où l'on
doit démontrer qu'une liste est \textit{complète} --- non pas trouver des exemples, mais fermer la porte.
La bonne méthode consiste à classer les patrons par la longueur de leur plus longue bande de carrés
alignés, ce qui transforme une recherche tâtonnante en un raisonnement fini. Les deux dés de la
question (c) sont symétriques l'un de l'autre sans être superposables : les dés occidentaux et
certains dés asiatiques anciens n'ont pas la même chiralité, et un joueur attentif peut le voir. La
question (d) contient en germe toute la théorie des dimensions --- pourquoi une fourmi peut tomber sans
dommage et pas un éléphant --- et c'est le même exposant $ 3 $ qui explique que Kepler, en empilant
des boulets de canon, ait posé un problème resté ouvert jusqu'en 1998.

\begin{refs}
\item Contexte : Patron (géométrie) --- Wikipédia (\url{https://fr.wikipedia.org/wiki/Patron_(g%C3%A9om%C3%A9trie)})
\item Contexte : Cube --- Wikipédia (\url{https://fr.wikipedia.org/wiki/Cube})
\item Contexte : Dé --- Wikipédia (\url{https://fr.wikipedia.org/wiki/D%C3%A9})
\end{refs}

\bigskip

\begin{problem}[{Le plus court chemin jusqu'à la rivière --- Collège}]
\textbf{GEO-47.} \textbf{Problème.} Une droite $ (d) $ représente la berge d'une rivière. Deux points $ A $ et
$ B $ sont situés du même côté de $ (d) $. Un promeneur part de $ A $, doit toucher la berge en
un point $ M $, puis rejoindre $ B $ ; il cherche à rendre $ AM + MB $ le plus petit possible.
\begin{enumerate}
\item[(a)] Soit $ A' $ le symétrique de $ A $ par rapport à $ (d) $. Démontrez que pour tout point
$ M $ de $ (d) $, on a $ AM = A'M $, donc $ AM + MB = A'M + MB $.
\item[(b)] Démontrez, à l'aide de l'inégalité triangulaire, que $ A'M + MB $ est minimal lorsque $ M $
est le point d'intersection de $ [A'B] $ avec $ (d) $, et que ce point existe et est unique.
\item[(c)] Démontrez qu'en ce point l'angle entre $ [AM] $ et $ (d) $ est égal à l'angle entre
$ [MB] $ et $ (d) $.
\item[(d)] Reprenez le problème lorsque $ A $ et $ B $ sont de part et d'autre de $ (d) $. Que devient
la solution, et pourquoi le détour par le symétrique n'est-il alors plus nécessaire ?
\end{enumerate}

\end{problem}

\noindent Héron d'Alexandrie, au Ier siècle, pose et résout ce problème dans sa
\textit{Catoptrique}, et en tire une conclusion audacieuse : puisque la lumière réfléchie suit un chemin
dont l'angle d'incidence égale l'angle de réflexion, c'est que \textit{la lumière prend le chemin le plus
court}. C'est le premier principe de minimalité de la physique. Fermat le reprend en 1662, corrige
\og{} le plus court \fg{} en \og{} le plus rapide \fg{} et en déduit la loi de la réfraction ; de là descend toute la
mécanique de Maupertuis, d'Euler et de Lagrange, et l'idée que les lois de la nature s'écrivent comme
des minimisations. Tout cela tient dans la symétrie axiale de la question (a) : un pliage de feuille
qui remplace un problème d'optimisation par une ligne droite.

\begin{refs}
\item Contexte : Symétrie axiale --- Wikipédia (\url{https://fr.wikipedia.org/wiki/Sym%C3%A9trie_axiale})
\item Contexte : Inégalité triangulaire --- Wikipédia (\url{https://fr.wikipedia.org/wiki/In%C3%A9galit%C3%A9_triangulaire})
\item Contexte : Héron d'Alexandrie --- Wikipédia (\url{https://fr.wikipedia.org/wiki/H%C3%A9ron_d%27Alexandrie})
\item Contexte : Principe de Fermat --- Wikipédia (\url{https://fr.wikipedia.org/wiki/Principe_de_Fermat})
\end{refs}

\bigskip

\begin{problem}[{Des triangles de même aire --- Collège, short}]
\textbf{GEO-48.} \textbf{Problème.} Soient $ A $ et $ B $ deux points d'une droite $ (d) $, et
$ (d') $ une droite parallèle à $ (d) $. Démontrez que tous les triangles $ ABM $, lorsque
$ M $ parcourt $ (d') $, ont la même aire, et expliquez pourquoi cette aire ne dépend pas de la
position de $ M $ alors que la forme du triangle, elle, change du tout au tout.
\end{problem}

\noindent C'est le lemme du \og{} cisaillement \fg{} : on peut faire glisser le sommet d'un triangle
le long d'une parallèle à la base sans rien changer à l'aire, et cette seule remarque démontre en trois
lignes la moitié des formules d'aire du collège. Euclide s'en sert aux propositions 37 et 38 du livre
I, juste avant d'attaquer Pythagore, et la démonstration célèbre de la proposition 47 n'est presque
rien d'autre que deux cisaillements bien choisis.

\begin{refs}
\item Contexte : Aire d'un triangle --- Wikipédia (\url{https://fr.wikipedia.org/wiki/Aire_d%27un_triangle})
\end{refs}

\bigskip

\begin{problem}[{Quels polygones pavent le plan --- Collège}]
\textbf{GEO-49.} \textbf{Problème.} On veut recouvrir le plan, sans trou ni chevauchement, avec des copies d'un même
polygone régulier, les carreaux se joignant côté contre côté et sommet contre sommet.
\begin{enumerate}
\item[(a)] En utilisant le résultat de GEO-41, démontrez que l'angle intérieur d'un polygone régulier à
$ n $ côtés vaut
\[ \frac{(n-2) \times 180^\circ}{n} . \]
\item[(b)] Autour de chaque sommet du pavage, les angles des carreaux doivent totaliser exactement
$ 360^\circ $. Démontrez que l'angle intérieur doit donc diviser $ 360^\circ $, et déduisez-en que
seuls le triangle équilatéral, le carré et l'hexagone régulier conviennent.
\item[(c)] Traitez explicitement le cas du pentagone régulier : calculez son angle et montrez où le calcul
échoue. Dessinez trois pentagones autour d'un point et mesurez le trou qui reste.
\item[(d)] Parmi les trois pavages possibles, l'hexagone est celui qui, à aire donnée, utilise la plus petite
longueur de cloisons. Cherchez ce que les abeilles en font, et énoncez en une phrase le théorème du
nid d'abeilles en précisant qui l'a démontré et quand.
\end{enumerate}

\end{problem}

\noindent La restriction aux trois pavages réguliers est déjà chez Pappus d'Alexandrie
(IVe siècle), qui remarque au passage que les abeilles ont \og{} choisi \fg{} l'hexagone et y voit
un signe de sagesse géométrique. L'intuition qu'il exprime --- l'hexagone minimise la cloison --- est
restée une conjecture pendant seize siècles, jusqu'à ce que Thomas Hales la démontre en 1999 : c'est
le théorème du nid d'abeilles. Levez ensuite la contrainte de régularité et le sujet explose : les
pavages du plan par des figures quelconques sont classés en dix-sept groupes de symétrie, dont les
mosaïques de l'Alhambra de Grenade illustrent une bonne part six siècles avant qu'on ne sache les
compter, et la question de savoir quels pentagones convexes pavent le plan n'a été close qu'en 2017,
par Michaël Rao.

\begin{refs}
\item Contexte : Pavage du plan --- Wikipédia (\url{https://fr.wikipedia.org/wiki/Pavage_du_plan})
\item Contexte : Polygone régulier --- Wikipédia (\url{https://fr.wikipedia.org/wiki/Polygone_r%C3%A9gulier})
\item Contexte : Pavage hexagonal --- Wikipédia (\url{https://fr.wikipedia.org/wiki/Pavage_hexagonal})
\item Contexte : Théorème du nid d'abeilles --- Wikipédia (en anglais) (\url{https://en.wikipedia.org/wiki/Honeycomb_theorem})
\end{refs}

\bigskip

\begin{problem}[{Deux miroirs parallèles --- Collège, short}]
\textbf{GEO-50.} \textbf{Problème.} Soient $ (d_1) $ et $ (d_2) $ deux droites parallèles
distinctes. Démontrez que la composée des deux symétries axiales d'axes $ (d_1) $ puis $ (d_2) $
est une translation, et déterminez sa direction, son sens et sa longueur en fonction de la distance
entre les deux droites.
\end{problem}

\noindent C'est le premier théorème de composition de transformations qu'on puisse démontrer
sans machinerie, et il ouvre une porte : les symétries axiales engendrent toutes les isométries du
plan, si bien que translations et rotations ne sont pas des objets nouveaux mais des produits de
réflexions. Remplacez \og{} parallèles \fg{} par \og{} sécantes \fg{} et vous obtiendrez une rotation d'angle double
--- la même démonstration, une figure différente --- et vous serez à la porte de la théorie des groupes,
qui classe les frises et les pavages de GEO-49.

\begin{refs}
\item Contexte : Symétrie axiale --- Wikipédia (\url{https://fr.wikipedia.org/wiki/Sym%C3%A9trie_axiale})
\item Contexte : Groupe de frise --- Wikipédia (\url{https://fr.wikipedia.org/wiki/Groupe_de_frise})
\end{refs}

\bigskip

\begin{problem}[{Voir n'est pas démontrer --- Collège, short}]
\textbf{GEO-51.} \textbf{Problème.} Tracez très soigneusement, sur cinq triangles quelconques et
franchement différents, le point de concours des médianes $ G $, celui des médiatrices $ O $ et
celui des hauteurs $ H $ ; constatez à la règle que ces trois points paraissent toujours alignés.
Rédigez ensuite, en une demi-page, pourquoi ces cinq figures --- et cent de plus --- ne démontrent rien,
et ce qu'il faudrait produire à la place.
\end{problem}

\noindent L'alignement est réel : c'est la droite d'Euler, publiée en 1765, deux mille ans
après Euclide --- rappel que la géométrie du triangle n'était pas close. Sa démonstration est au-dessus
du collège et vous attend en GEO-03, plus bas sur cette page. Ce qui est à votre portée aujourd'hui,
et qui compte davantage, c'est de savoir nommer l'écart entre une régularité observée et un théorème :
la conjecture de Pólya, vérifiée sur des centaines de millions de cas, s'est révélée fausse ; celle de
Mertens aussi. Une figure vous dit où chercher, elle ne vous dit jamais que vous avez trouvé.

\begin{refs}
\item Contexte : Droite d'Euler --- Wikipédia (\url{https://fr.wikipedia.org/wiki/Droite_d%27Euler})
\item Voir aussi : L'art de la démonstration (\url{proofs.html})
\end{refs}

\bigskip

\begin{problem}[{Reconnaître deux droites parallèles --- Collège, short}]
\textbf{GEO-52.} \textbf{Problème.} Deux droites $ (d) $ et $ (d') $ sont coupées par une
sécante $ (s) $, et deux angles alternes-internes ainsi formés sont égaux. Démontrez que $ (d) $
et $ (d') $ sont parallèles, en nommant précisément l'axiome ou le théorème du cours qui autorise
chacun de vos pas.
\end{problem}

\noindent C'est la proposition 27 du livre I d'Euclide, et sa place dans les \textit{Éléments}
est remarquable : elle est démontrée \textit{sans} le cinquième postulat, alors que sa réciproque --- si
les droites sont parallèles alors les angles alternes-internes sont égaux, proposition 29 --- en a
besoin. Les deux énoncés se ressemblent comme deux gouttes d'eau et n'ont pas du tout le même statut
logique ; savoir dans quel sens on raisonne est ici la seule chose qui compte.

\begin{refs}
\item Contexte : Angles alternes-internes --- Wikipédia (\url{https://fr.wikipedia.org/wiki/Angles_alternes-internes})
\item Contexte : Axiome des parallèles --- Wikipédia (\url{https://fr.wikipedia.org/wiki/Axiome_des_parall%C3%A8les})
\end{refs}

\bigskip


\section*{Combinatoire et théorie des graphes}

\begin{problem}[{Les poignées de main d'une classe --- Collège}]
\textbf{CMB-45.} \textbf{Problème.} Dans une salle, chaque personne présente serre la main de chaque autre personne
exactement une fois.
\begin{enumerate}
\item[(a)] Six personnes sont présentes. Dressez la liste complète des poignées de main et dénombrez-les,
puis expliquez pourquoi votre liste n'en oublie aucune et n'en compte aucune deux fois.
\item[(b)] Démontrez que pour $ n $ personnes le nombre de poignées de main vaut $ \frac{n(n-1)}{2} $.
On pourra compter d'abord les poignées de main \og{} vues \fg{} par chacun, puis corriger le double
comptage --- et il faudra justifier que la correction est bien une division par 2.
\item[(c)] Un soir, 66 poignées de main ont été échangées et chacun a bien serré la main de tous les
autres. Déterminez, en le justifiant, combien de personnes étaient présentes, et expliquez pourquoi
il ne peut pas y avoir deux réponses.
\end{enumerate}

\end{problem}

\noindent C'est le premier dénombrement que l'on rencontre en général, et il cache déjà la
technique la plus rentable de toute la combinatoire : compter la même collection de deux manières,
puis identifier les deux résultats. Le nombre obtenu est aussi le nombre de segments joignant
$ n $ points, le nombre d'arêtes du graphe complet à $ n $ sommets et le $ (n-1) $-ième nombre
triangulaire --- trois questions qui n'ont l'air de rien avoir en commun jusqu'à ce qu'on les compte.
Euler, en 1736, comptait des ponts ; la même façon de raisonner sert ici à compter des mains
serrées.

\begin{refs}
\item Contexte : Graphe complet --- Wikipédia (\url{https://fr.wikipedia.org/wiki/Graphe_complet})
\item Contexte : Nombre triangulaire --- Wikipédia (\url{https://fr.wikipedia.org/wiki/Nombre_triangulaire})
\end{refs}

\bigskip

\begin{problem}[{Les diagonales d'un polygone --- Collège}]
\textbf{CMB-46.} \textbf{Problème.} Dans un polygone convexe, une diagonale est un segment qui joint deux sommets
non voisins.
\begin{enumerate}
\item[(a)] Tracez toutes les diagonales d'un pentagone convexe, puis celles d'un hexagone convexe, et
dénombrez-les. Expliquez la méthode qui vous garantit de n'en oublier aucune.
\item[(b)] Démontrez qu'un polygone convexe à $ n $ côtés possède exactement
$ \frac{n(n-3)}{2} $ diagonales. Justifiez séparément d'où vient le $ -3 $ et d'où vient la
division par 2.
\item[(c)] Déterminez, en le démontrant, s'il existe un polygone convexe ayant exactement 100 diagonales.
Un décompte pour quelques valeurs de $ n $ ne suffit pas : il faut expliquer pourquoi les valeurs
non testées sont écartées elles aussi.
\end{enumerate}

\end{problem}

\noindent La question a l'air d'un exercice de dessin et elle est en réalité le même
problème que celui des poignées de main, à trois exceptions près par sommet : chaque sommet est
relié à tous les autres sauf lui-même et ses deux voisins. Reconnaître qu'un problème est un
ancien problème déguisé est une compétence à part entière, et c'est précisément ce que fait la
combinatoire quand elle appelle \og{} bijection \fg{} le fait de traduire une question en une autre. La
partie (c) apprend autre chose encore : pour démontrer qu'une chose n'existe pas, il ne suffit
jamais d'avoir cherché sans trouver.

\begin{refs}
\item Contexte : Diagonale --- Wikipédia (\url{https://fr.wikipedia.org/wiki/Diagonale})
\item Contexte : Polygone --- Wikipédia (\url{https://fr.wikipedia.org/wiki/Polygone})
\end{refs}

\bigskip

\begin{problem}[{Chemins sur un quadrillage --- Collège}]
\textbf{CMB-47.} \textbf{Problème.} Une fourmi part du coin inférieur gauche d'un quadrillage et se rend au coin
supérieur droit en se déplaçant uniquement le long des traits, d'un carreau vers la droite ou d'un
carreau vers le haut.
\begin{enumerate}
\item[(a)] Sur un quadrillage carré de 4 carreaux de côté, qui compte donc 25 nœuds, inscrivez à chaque
nœud le nombre de chemins qui y mènent depuis le départ, en additionnant à chaque fois le nombre
écrit à gauche et le nombre écrit en dessous. Démontrez que cette règle d'addition est légitime,
c'est-à-dire que tout chemin arrivant en un nœud y arrive nécessairement par l'une des deux
directions et jamais par les deux.
\item[(b)] Comparez le tableau de nombres obtenu au triangle de Pascal et expliquez précisément pourquoi
les deux constructions produisent les mêmes nombres.
\item[(c)] On interdit maintenant le passage par le nœud central du quadrillage, celui qui se trouve à
2 carreaux de chacun des quatre bords. Dénombrez les chemins restants et justifiez votre méthode :
si vous soustrayez, démontrez que vous soustrayez exactement les chemins interdits, chacun une
seule fois.
\end{enumerate}

\end{problem}

\noindent Le triangle apparaît dans la prosodie indienne de Pingala (vers 200 av. J.-C.),
chez le mathématicien persan al-Karaji vers l'an mille et chez le Chinois Yang Hui au
XIIIe siècle, bien avant le traité de Pascal de 1654 ; il porte donc, comme souvent en
mathématiques, le nom du dernier arrivé. Le dénombrement des chemins en est la lecture la plus
parlante : chaque case du triangle compte les façons d'y descendre depuis le sommet. La partie (c)
introduit le raisonnement par soustraction, qui deviendra plus tard le principe d'inclusion et
d'exclusion.

\begin{refs}
\item Contexte : Triangle de Pascal --- Wikipédia (\url{https://fr.wikipedia.org/wiki/Triangle_de_Pascal})
\item Contexte : Dénombrement --- Wikipédia (\url{https://fr.wikipedia.org/wiki/D%C3%A9nombrement})
\end{refs}

\bigskip

\begin{problem}[{L'arbre des menus --- Collège, short}]
\textbf{CMB-48.} \textbf{Problème.} Un restaurant propose 3 entrées, 4 plats et 2 desserts, et un
menu se compose d'une entrée, d'un plat et d'un dessert. Dessinez l'arbre des choix, dénombrez les
menus possibles, puis démontrez la règle générale qui donne ce nombre sans qu'il soit nécessaire de
dessiner l'arbre.
\end{problem}

\noindent On appelle cette règle le principe multiplicatif, et l'arbre en est la
démonstration dessinée : chaque niveau de branches multiplie le nombre de feuilles. Elle paraît
évidente jusqu'au jour où les choix cessent d'être indépendants --- un plat interdit avec une entrée,
par exemple --- et où l'arbre, lui, continue de dire la vérité que la multiplication ne dit plus.

\begin{refs}
\item Contexte : Dénombrement --- Wikipédia (\url{https://fr.wikipedia.org/wiki/D%C3%A9nombrement})
\item Contexte : Produit cartésien --- Wikipédia (\url{https://fr.wikipedia.org/wiki/Produit_cart%C3%A9sien})
\end{refs}

\bigskip

\begin{problem}[{Des chaussettes dans le noir --- Collège, short}]
\textbf{CMB-49.} \textbf{Problème.} Un tiroir contient 10 chaussettes noires, 10 bleues et 10
rouges, toutes mêlées, et l'on y pioche des chaussettes sans allumer la lumière. Déterminez, en le
démontrant, le plus petit nombre de chaussettes à tirer pour être certain d'avoir une paire de la
même couleur, puis le plus petit nombre garantissant une paire noire.
\end{problem}

\noindent C'est le principe des tiroirs sous sa forme la plus domestique : Dirichlet
l'énonçait dans les années 1830 sous le nom de \textit{Schubfachprinzip} pour démontrer des résultats
d'approximation des irrationnels. Une démonstration complète comporte toujours deux moitiés qu'on
oublie souvent de distinguer : montrer que le nombre annoncé suffit, et exhiber un tirage
malchanceux montrant qu'un de moins ne suffirait pas.

\begin{refs}
\item Contexte : Principe des tiroirs --- Wikipédia (\url{https://fr.wikipedia.org/wiki/Principe_des_tiroirs})
\end{refs}

\bigskip

\begin{problem}[{Les sept ponts de Königsberg --- Collège}]
\textbf{CMB-50.} \textbf{Problème.} La ville de Königsberg était bâtie sur deux rives et deux îles de la rivière
Pregel, reliées par sept ponts ; ses habitants se demandaient s'il existait une promenade
empruntant chaque pont une fois et une seule. La rive nord touche 3 ponts, la rive sud 3 ponts, la
grande île 5 ponts et la petite île 3 ponts.
\begin{enumerate}
\item[(a)] Redessinez la ville en remplaçant chaque morceau de terre par un point et chaque pont par un
trait joignant deux points. Expliquez pourquoi la forme des rives, les distances et la longueur des
ponts ne changent rien à la question.
\item[(b)] Démontrez qu'une telle promenade est impossible. On raisonnera sur un morceau de terre qui
n'est ni le départ ni l'arrivée : à chaque fois que la promenade y entre, elle doit en ressortir
par un autre pont.
\item[(c)] Déterminez, en le justifiant, le nombre minimal de ponts qu'il faudrait ajouter pour qu'une
telle promenade devienne possible, et dites si elle pourrait alors revenir à son point de
départ.
\end{enumerate}

\end{problem}

\noindent Leonhard Euler règle la question en 1736 dans un mémoire pour
l'Académie de Saint-Pétersbourg, et il commence par écarter la géométrie : ni les longueurs ni les
angles n'interviennent, seule compte la façon dont les choses sont reliées. Il crée du même coup
deux disciplines, la théorie des graphes et la topologie, à partir d'une promenade dominicale. La
ville s'appelle aujourd'hui Kaliningrad et deux des sept ponts ont disparu pendant la Seconde
Guerre mondiale ; la question, elle, se repose à l'identique sur la carte d'aujourd'hui.

\begin{refs}
\item Contexte : Problème des sept ponts de Königsberg --- Wikipédia (\url{https://fr.wikipedia.org/wiki/Probl%C3%A8me_des_sept_ponts_de_K%C3%B6nigsberg})
\item Contexte : Leonhard Euler --- Wikipédia (\url{https://fr.wikipedia.org/wiki/Leonhard_Euler})
\item Contexte : Théorie des graphes --- Wikipédia (\url{https://fr.wikipedia.org/wiki/Th%C3%A9orie_des_graphes})
\end{refs}

\bigskip

\begin{problem}[{Combien de matchs dans le tournoi ? --- Collège}]
\textbf{CMB-51.} \textbf{Problème.} Deux formats de compétition se disputent le monde du sport : l'élimination
directe, où le perdant d'un match quitte le tournoi, et le championnat toutes rondes, où chaque
équipe rencontre chaque autre.
\begin{enumerate}
\item[(a)] Un tournoi à élimination directe réunit 64 équipes et se termine par un unique vainqueur.
Déterminez le nombre total de matchs joués et démontrez-le sans additionner les tours un par un.
On pourra compter ce que chaque match produit d'irréversible.
\item[(b)] Démontrez que votre raisonnement vaut encore pour $ n $ équipes lorsque $ n $ n'est pas une
puissance de 2 et que certaines équipes sont exemptées au premier tour.
\item[(c)] Dans un championnat où chaque équipe rencontre chaque autre deux fois, à l'aller et au retour,
182 matchs ont été joués. Déterminez le nombre d'équipes engagées, en justifiant que c'est le
seul possible.
\end{enumerate}

\end{problem}

\noindent Les deux formats donnent des décomptes de natures opposées : l'un croît
proportionnellement au nombre d'équipes, l'autre bien plus vite, et c'est la raison pratique pour
laquelle les grandes compétitions mélangent poules et tableau final. L'argument attendu en (a) est
un modèle du genre --- on ne compte pas les matchs, on compte les éliminations, et l'on remarque que
la correspondance entre les deux est exacte. C'est ce qu'on appelle une démonstration par bijection,
et elle vaut aussi bien pour 64 équipes que pour 4 096.

\begin{refs}
\item Contexte : Tournoi toutes rondes --- Wikipédia (\url{https://fr.wikipedia.org/wiki/Tournoi_toutes_rondes})
\item Contexte : Tournoi à élimination directe --- Wikipédia (en anglais) (\url{https://en.wikipedia.org/wiki/Single-elimination_tournament})
\end{refs}

\bigskip

\begin{problem}[{Colorier une carte avec peu de couleurs --- Collège}]
\textbf{CMB-52.} \textbf{Problème.} Colorier une carte, c'est attribuer une couleur à chaque région de sorte que
deux régions partageant une frontière --- un morceau de ligne, pas seulement un point --- reçoivent
des couleurs différentes.
\begin{enumerate}
\item[(a)] Dessinez une carte de quatre régions qui exige réellement quatre couleurs, et démontrez que
trois ne suffisent pas pour votre carte.
\item[(b)] On trace sur une feuille plusieurs droites qui la traversent entièrement ; elles la découpent
en régions. Démontrez que deux couleurs suffisent toujours à colorier cette carte, en examinant ce
qui se passe lorsqu'on ajoute les droites une par une.
\item[(c)] Le théorème des quatre couleurs affirme que quatre couleurs suffisent pour toute carte tracée
sur une feuille. Construisez une carte de six régions qui se colorie avec trois couleurs et une
autre carte de six régions qui en exige quatre, et démontrez ces deux affirmations.
\end{enumerate}

\end{problem}

\noindent La question est posée en 1852 par Francis Guthrie, qui coloriait les comtés
d'Angleterre, transmise par son frère à De Morgan, et elle résiste ensuite plus d'un siècle : la
démonstration d'Alfred Kempe, publiée en 1879 et acclamée, s'effondre en 1890 sous l'examen de
Percy Heawood, qui sauve tout de même le résultat pour cinq couleurs. Kenneth Appel et Wolfgang
Haken concluent en 1976 à l'aide d'un ordinateur, en vérifiant à la machine près de 1 500
configurations --- première grande démonstration qu'aucun être humain ne peut lire en entier, et
première dispute sérieuse sur ce qu'est une démonstration. Les parties (a) et (b), elles, se
démontrent au crayon.

\begin{refs}
\item Contexte : Théorème des quatre couleurs --- Wikipédia (\url{https://fr.wikipedia.org/wiki/Th%C3%A9or%C3%A8me_des_quatre_couleurs})
\item Contexte : Coloration de graphe --- Wikipédia (\url{https://fr.wikipedia.org/wiki/Coloration_de_graphe})
\end{refs}

\bigskip

\begin{problem}[{Combien de carrés sur un damier ? --- Collège}]
\textbf{CMB-53.} \textbf{Problème.} On considère un damier $ 8 \times 8 $ et l'on ne s'intéresse qu'aux carrés
et rectangles dont les côtés suivent les lignes de la grille.
\begin{enumerate}
\item[(a)] Dénombrez les carrés de côté 1, puis ceux de côté 2, puis ceux de côté 3. Justifiez chaque
décompte en décrivant précisément ce qui détermine un carré une fois sa taille fixée.
\item[(b)] Démontrez que le nombre total de carrés de toutes tailles d'un damier $ n \times n $ vaut
\[ 1^2 + 2^2 + 3^2 + \cdots + n^2 . \]
\item[(c)] Dénombrez maintenant tous les rectangles du damier $ 8 \times 8 $, carrés compris, et
démontrez que votre décompte est complet et sans répétition.
\end{enumerate}

\end{problem}

\noindent Le damier est le laboratoire préféré des dénombreurs parce qu'il est fini,
familier, et qu'il punit immédiatement les décomptes bâclés : la réponse \og{} soixante-quatre \fg{} à la
question du titre est fausse et tout le monde la donne d'abord. La somme des carrés de la partie
(b) était connue d'Archimède, qui s'en servait dans \textit{Sur les spirales} pour évaluer une aire, et les
nombres qu'elle produit s'appellent nombres pyramidaux carrés --- ils comptent les boulets empilés en
pyramide à base carrée. La partie (c) demande le geste décisif de la combinatoire : décrire un
objet par les choix qui le déterminent entièrement.

\begin{refs}
\item Contexte : Nombre pyramidal carré --- Wikipédia (\url{https://fr.wikipedia.org/wiki/Nombre_pyramidal_carr%C3%A9})
\item Contexte : Échiquier --- Wikipédia (\url{https://fr.wikipedia.org/wiki/%C3%89chiquier})
\end{refs}

\bigskip

\begin{problem}[{Le damier et les dominos --- Collège}]
\textbf{CMB-54.} \textbf{Problème.} Un domino recouvre exactement deux cases voisines d'un damier, côte à côte ou
l'une au-dessus de l'autre. Recouvrir un damier signifie le couvrir entièrement, sans chevauchement
ni débordement.
\begin{enumerate}
\item[(a)] Démontrez qu'un damier $ 5 \times 5 $ ne peut pas être recouvert par des dominos.
\item[(b)] On colorie le damier $ 8 \times 8 $ en noir et blanc comme un échiquier. Démontrez que tout
domino posé couvre exactement une case noire et une case blanche, puis qu'après avoir retiré deux
cases de la même couleur le reste ne peut plus être recouvert.
\item[(c)] Une case a été retirée en haut à gauche et une autre en bas à droite. Expliquez, à partir de
(b), pourquoi 31 dominos ne peuvent pas recouvrir les 62 cases restantes, et dites ce que devient
la conclusion si les deux cases retirées sont de couleurs différentes.
\end{enumerate}

\end{problem}

\noindent Popularisé par le philosophe Max Black en 1946 puis par les chroniques de Martin
Gardner dans \textit{Scientific American}, ce petit problème est le modèle de l'argument par
invariant : au lieu d'essayer les pavages un à un --- il y en a beaucoup trop --- on repère une
quantité que chaque coup légal laisse inchangée, et l'on constate qu'elle interdit l'arrivée.
L'idée sert bien au-delà du damier : c'est elle qui décide du taquin, des jeux de jetons et de
quantité de problèmes d'olympiades. La question réciproque --- quels retraits de cases laissent un
pavage possible --- est plus fine et se traite au lycée.

\begin{refs}
\item Contexte : Problème du damier mutilé --- Wikipédia (en anglais) (\url{https://en.wikipedia.org/wiki/Mutilated_chessboard_problem})
\item Contexte : Damier --- Wikipédia (\url{https://fr.wikipedia.org/wiki/Damier})
\item Contexte : Martin Gardner --- Wikipédia (\url{https://fr.wikipedia.org/wiki/Martin_Gardner})
\end{refs}

\bigskip

\begin{problem}[{Le motif d'allumettes --- Collège, short}]
\textbf{CMB-55.} \textbf{Problème.} On aligne des carrés en les collant côte à côte : un carré
demande 4 allumettes, deux carrés en demandent 7, trois carrés en demandent 10. Déterminez, en le
démontrant, le nombre d'allumettes nécessaires pour $ n $ carrés alignés, puis le nombre de carrés
que l'on peut construire avec exactement 2026 allumettes.
\end{problem}

\noindent Poursuivre un motif est facile ; démontrer que la règle devinée est la bonne est
un autre métier, et c'est celui-là que le problème demande. La bonne démonstration ne recalcule pas
les premiers cas : elle explique ce que coûte le passage d'un carré au suivant, et pourquoi ce coût
ne change jamais. C'est le premier contact avec les suites arithmétiques, et déjà avec l'idée que
l'on démontre en regardant la marche, pas les marches.

\begin{refs}
\item Contexte : Suite arithmétique --- Wikipédia (\url{https://fr.wikipedia.org/wiki/Suite_arithm%C3%A9tique})
\item Contexte : Suite (mathématiques) --- Wikipédia (\url{https://fr.wikipedia.org/wiki/Suite_(math%C3%A9matiques)})
\end{refs}

\bigskip

\begin{problem}[{Deux langues et un diagramme --- Collège, short}]
\textbf{CMB-56.} \textbf{Problème.} Dans une classe de 30 élèves, 18 étudient l'anglais, 15
étudient l'espagnol et 7 étudient les deux langues. Déterminez combien d'élèves n'étudient aucune
de ces deux langues, en justifiant votre décompte par un diagramme et par une phrase, puis démontrez
la règle générale que vous avez employée.
\end{problem}

\noindent Additionner 18 et 15 compte deux fois les élèves qui font les deux langues :
c'est l'erreur que le principe d'inclusion et d'exclusion existe pour réparer. Les diagrammes qui
portent le nom de John Venn datent de 1880, mais l'idée de corriger un décompte en retranchant les
recouvrements est bien plus ancienne --- on la trouve chez de Moivre au XVIIIe siècle. Avec
trois langues au lieu de deux, la règle se complique d'un terme supplémentaire, et il vaut la peine
de deviner lequel avant de le lire.

\begin{refs}
\item Contexte : Principe d'inclusion-exclusion --- Wikipédia (\url{https://fr.wikipedia.org/wiki/Principe_d%27inclusion-exclusion})
\item Contexte : Diagramme de Venn --- Wikipédia (\url{https://fr.wikipedia.org/wiki/Diagramme_de_Venn})
\end{refs}

\bigskip


\section*{Probabilités}

\begin{problem}[{Sept contre douze --- Collège}]
\textbf{PRB-45.} \textbf{Problème.} On lance deux dés équilibrés à six faces, l'un rouge et l'un vert, et on
additionne les deux résultats. Un joueur affirme : \og{} il y a onze sommes possibles, de 2 à 12, donc
chacune a une chance sur onze de sortir \fg{}.
\begin{enumerate}
\item[(a)] Expliquez pourquoi ce sont les $ 36 $ couples (dé rouge, dé vert) qui forment les cas
également possibles, et non les onze sommes. Votre argument doit dire précisément ce qui distingue
les deux façons de compter.
\item[(b)] Dressez le tableau des 36 issues, puis comparez la probabilité d'obtenir 7 et celle d'obtenir
12. Justifiez l'écart que vous trouvez en décrivant, en français et sans formule, ce que la somme 7
a de particulier parmi toutes les sommes.
\item[(c)] Le joueur se défend : \og{} j'ai lancé les dés 36 fois et je n'ai obtenu 7 que quatre fois ; votre
calcul est donc faux \fg{}. Rédigez une réponse qui dit ce qu'une probabilité promet et ce qu'elle ne
promet pas.
\end{enumerate}

\end{problem}

\noindent Le grand-duc de Toscane posa à Galilée, vers 1620, la même question pour trois
dés : les joueurs savaient d'expérience que 10 sortait plus souvent que 9, alors que les deux
sommes s'écrivent chacune de six façons. Galilée trancha en quelques lignes, en comptant les issues
ordonnées plutôt que les descriptions d'issues --- c'est le premier raisonnement de probabilité qui
mérite ce nom, et c'est exactement le geste demandé à la question (a). La version à trois dés est
PRB-01, sur cette même page.

\begin{refs}
\item Contexte : Probabilité --- Wikipédia (\url{https://fr.wikipedia.org/wiki/Probabilit%C3%A9})
\item Contexte : Dé --- Wikipédia (\url{https://fr.wikipedia.org/wiki/D%C3%A9})
\item Contexte : Galilée --- Wikipédia (\url{https://fr.wikipedia.org/wiki/Galil%C3%A9e_(savant)})
\end{refs}

\bigskip

\begin{problem}[{La suite qui n'a pas l'air au hasard --- Collège}]
\textbf{PRB-46.} \textbf{Problème.} Deux élèves devaient produire une suite de 30 résultats de pile ou face. L'un a
réellement lancé une pièce et noté ce qu'il obtenait ; l'autre a inventé sa suite de tête. Voici
les deux copies, où P désigne pile et F face.

\noindent Suite 1 : P F P F F P F P P F P F F P F P F P P F P F P F F P F P F P 
Suite 2 : P P F F F F P F P P P F F P F F F F F P P F P P P P F P F F
\begin{enumerate}
\item[(a)] Pour chacune des deux suites, calculez la fréquence de \og{} pile \fg{} et relevez la plus longue série
de résultats identiques consécutifs. Présentez vos résultats dans un tableau.
\item[(b)] Dites laquelle des deux suites vous croyez inventée et défendez votre choix. On attend un
argument appuyé sur les nombres de la question (a), pas une impression visuelle.
\item[(c)] Justifiez qu'une suite précise de 30 lancers a exactement la même probabilité qu'une autre
suite précise de 30 lancers --- y compris celle qui donne trente fois \og{} pile \fg{}. Expliquez ensuite,
en un paragraphe rédigé, pourquoi cela ne contredit pas votre réponse à la question (b).
\end{enumerate}

\end{problem}

\noindent Cette expérience est une démonstration de cours devenue classique --- souvent
attribuée au mathématicien américain Ted Hill : le professeur sort de la salle, une moitié de la
classe lance vraiment une pièce cent fois, l'autre moitié invente sa suite, et il retrouve les
tricheurs en quelques secondes. Les êtres humains alternent trop et produisent presque toujours
des séries trop courtes, parce que nous confondons \og{} au hasard \fg{} et \og{} bien mélangé \fg{}. La question
(c) est le cœur de l'affaire : le hasard ne rend pas certaines suites impossibles, il rend rares
certaines \textit{propriétés} des suites --- la version chiffrée de cet argument est PRB-05, dans la
bande Lycée.

\begin{refs}
\item Contexte : Pile ou face --- Wikipédia (\url{https://fr.wikipedia.org/wiki/Pile_ou_face})
\item Contexte : Hasard --- Wikipédia (\url{https://fr.wikipedia.org/wiki/Hasard})
\item Contexte : Ted Hill --- Wikipédia (en anglais) (\url{https://en.wikipedia.org/wiki/Ted_Hill_(mathematician)})
\end{refs}

\bigskip

\begin{problem}[{Deux chaussettes dans le noir --- Collège, short}]
\textbf{PRB-47.} \textbf{Problème.} Un tiroir contient six chaussettes noires et quatre
chaussettes bleues mélangées, et vous en tirez deux au hasard dans l'obscurité, sans remettre la
première avant de tirer la seconde. Déterminez la probabilité d'obtenir deux chaussettes de la même
couleur en justifiant chaque branche de votre arbre, puis expliquez précisément ce qui changerait
si vous remettiez la première chaussette dans le tiroir.
\end{problem}

\noindent Le tiroir de chaussettes est un exercice de manuel depuis au moins les années
1950, et il circule sous deux formes que tout oppose. La forme \og{} combien faut-il en tirer pour être
\textit{certain} d'avoir une paire ? \fg{} ne demande aucune probabilité : elle relève du principe des
tiroirs et se règle par un décompte. La forme posée ici demande un arbre et donne un nombre entre 0
et 1. Savoir laquelle des deux questions on est en train de résoudre est déjà une compétence de
mathématicien.

\begin{refs}
\item Contexte : Problème d'urne --- Wikipédia (\url{https://fr.wikipedia.org/wiki/Probl%C3%A8me_d%27urne})
\item Contexte : Arbre de probabilité --- Wikipédia (\url{https://fr.wikipedia.org/wiki/Arbre_de_probabilit%C3%A9})
\item Contexte : Principe des tiroirs --- Wikipédia (\url{https://fr.wikipedia.org/wiki/Principe_des_tiroirs})
\end{refs}

\bigskip

\begin{problem}[{La roue du forain --- Collège}]
\textbf{PRB-48.} \textbf{Problème.} À la fête foraine, une roue est partagée en huit secteurs de même taille,
numérotés de 1 à 8 ; le forain jure qu'elle est parfaitement équilibrée. Vous restez une heure à
noter les résultats. Sur 400 tours, les secteurs 1 à 8 sont sortis respectivement 62, 54, 47, 51,
58, 49, 66 et 13 fois.
\begin{enumerate}
\item[(a)] Calculez la fréquence de sortie de chaque secteur et l'effectif que l'on attendrait pour
chacun si la roue était équilibrée. Rassemblez le tout dans un tableau et commentez ligne par
ligne.
\item[(b)] Le forain répond : \og{} le hasard fait ce qu'il veut, vos écarts ne prouvent rien \fg{}. Il a raison
sur le principe. Construisez malgré tout l'argument le plus solide que vous puissiez écrire au
collège pour soutenir que la roue est truquée --- puis dites honnêtement ce qui manque à votre
argument pour être une preuve.
\item[(c)] Le secteur 8 est celui qui fait gagner le gros lot. Expliquez pourquoi un forain malhonnête a
intérêt à truquer précisément ce secteur-là, et pourquoi un joueur qui ne regarde qu'une dizaine de
tours n'a pratiquement aucune chance de s'en apercevoir.
\end{enumerate}

\end{problem}

\noindent Cette méthode a réellement fait fortune. En 1873, l'ingénieur anglais Joseph
Jagger paya des employés pour relever pendant des semaines les numéros sortis aux roulettes du
casino de Monte-Carlo ; l'une des roues était légèrement déséquilibrée, certains numéros sortaient
trop souvent, et Jagger repartit avec une somme considérable. Les casinos rééquilibrent depuis
leurs roues et changent leur emplacement, précisément pour empêcher ce genre de relevé. L'idée que
la fréquence observée finit par ressembler à la probabilité porte un nom --- la loi des grands
nombres --- et se démontre au lycée ; ce qui vous manque ici, et que la question (b) vous demande de
nommer, est justement l'outil qui dit à partir de quel écart on peut cesser de croire au
forain.

\begin{refs}
\item Contexte : Fréquence (statistiques) --- Wikipédia (\url{https://fr.wikipedia.org/wiki/Fr%C3%A9quence_(statistiques)})
\item Contexte : Loi des grands nombres --- Wikipédia (\url{https://fr.wikipedia.org/wiki/Loi_des_grands_nombres})
\item Contexte : Joseph Jagger --- Wikipédia (en anglais) (\url{https://en.wikipedia.org/wiki/Joseph_Jagger})
\end{refs}

\bigskip

\begin{problem}[{Le jeu est-il équitable ? --- Collège}]
\textbf{PRB-49.} \textbf{Problème.} Alice et Bruno lancent ensemble deux dés équilibrés à six faces. Dans chacune
des règles ci-dessous, il n'y a ni match nul ni relance.
\begin{enumerate}
\item[(a)] Règle 1 : Alice gagne si la somme des deux dés est paire, Bruno si elle est impaire. Le jeu
est-il équitable ? Justifiez à partir du tableau des 36 issues.
\item[(b)] Règle 2 : Alice gagne si le produit des deux dés est pair, Bruno s'il est impair. Justifiez
que ce jeu n'est pas équitable, et dites de combien exactement il penche.
\item[(c)] Pour compenser la règle 2, Bruno propose de recevoir $ k $ euros chaque fois qu'il gagne et
de payer 1 euro chaque fois qu'il perd. Déterminez, en raisonnant sur un grand nombre de parties,
la valeur de $ k $ qui rend ce jeu équitable, et justifiez-la. Expliquez ensuite pourquoi
\og{} équitable \fg{} ne veut pas dire la même chose à la question (a) et à la question (c).
\end{enumerate}

\end{problem}

\noindent La notion de jeu équitable est le berceau de toute la théorie : en 1654, le
chevalier de Méré --- de son vrai nom Antoine Gombaud --- apporta à Blaise Pascal des questions de ce
type, et la correspondance entre Pascal et Fermat qui s'ensuivit fonda la discipline. Huygens en
tira en 1657 le premier traité imprimé, où il définit la valeur d'un jeu comme la mise qu'un joueur
devrait accepter de payer pour y entrer : c'est l'ancêtre direct de votre question (c). L'outil qui
la formalise s'appelle l'espérance et se rencontre au lycée ; au collège, le raisonnement \og{} sur un
grand nombre de parties \fg{} suffit, et il a l'avantage de dire ce que l'espérance signifie avant de
dire comment on la calcule.

\begin{refs}
\item Contexte : Antoine Gombaud, chevalier de Méré --- Wikipédia (\url{https://fr.wikipedia.org/wiki/Antoine_Gombaud})
\item Contexte : Christian Huygens --- Wikipédia (\url{https://fr.wikipedia.org/wiki/Christian_Huygens})
\item Contexte : Espérance mathématique --- Wikipédia (\url{https://fr.wikipedia.org/wiki/Esp%C3%A9rance_math%C3%A9matique})
\end{refs}

\bigskip

\begin{problem}[{Un anniversaire commun dans votre classe --- Collège}]
\textbf{PRB-50.} \textbf{Problème.} Un élève d'une classe de 30 affirme : \og{} nous sommes 30 et l'année compte 365
jours, donc il n'y a presque aucune chance que deux d'entre nous soient nés le même jour \fg{}.
\begin{enumerate}
\item[(a)] Comptez le nombre de paires d'élèves que l'on peut former dans cette classe. Expliquez
pourquoi c'est ce nombre-là, et non 30, qu'il faut comparer à 365, et pourquoi le raisonnement de
l'élève passe à côté du problème.
\item[(b)] Distinguez soigneusement deux questions : \og{} un autre élève est-il né le même jour que
\textit{moi} ? \fg{} et \og{} deux élèves quelconques de la classe sont-ils nés le même jour ? \fg{}. Justifiez
que la seconde a beaucoup plus de chances d'être vraie que la première, sans calculer aucune des
deux probabilités.
\item[(c)] Menez l'expérience. Relevez les dates de naissance de votre classe et de trois classes
voisines, comptez combien de ces quatre classes contiennent au moins un anniversaire commun, et
confrontez le résultat à votre raisonnement de la question (a). Discutez enfin ce que quatre
classes permettent --- et ne permettent pas --- de conclure.
\end{enumerate}

\end{problem}

\noindent Le paradoxe des anniversaires fut posé sous cette forme par Richard von Mises en
1939 et rendu célèbre par le manuel de William Feller. Il n'a rien d'un paradoxe logique : c'est
une contradiction entre un calcul juste et une intuition fausse, et l'intuition se trompe parce
qu'elle compte les personnes là où il faut compter les paires. Le calcul exact, avec le seuil
fameux de 23 personnes, est PRB-03 dans la bande Lycée ; la question (a) ci-dessus en contient déjà
toute l'idée, et la question (c) vous fait faire ce que von Mises n'a jamais pris la peine de
faire.

\begin{refs}
\item Contexte : Paradoxe des anniversaires --- Wikipédia (\url{https://fr.wikipedia.org/wiki/Paradoxe_des_anniversaires})
\item Contexte : Probabilité --- Wikipédia (\url{https://fr.wikipedia.org/wiki/Probabilit%C3%A9})
\end{refs}

\bigskip

\begin{problem}[{Deux enfants, dont une fille --- Collège}]
\textbf{PRB-51.} \textbf{Problème.} On admet qu'à chaque naissance, garçon et fille sont également probables, et
on considère une famille de deux enfants.
\begin{enumerate}
\item[(a)] Dressez l'arbre des quatre familles possibles, en distinguant l'aîné du cadet, et justifiez
que ces quatre issues sont également probables.
\item[(b)] On vous dit seulement : \og{} cette famille a au moins une fille \fg{}. Déterminez la probabilité que
les deux enfants soient des filles, en justifiant à partir de votre arbre quelles branches vous
gardez et lesquelles vous éliminez.
\item[(c)] Situation différente : vous croisez dans la rue un enfant de cette famille, choisi au hasard
parmi les deux, et c'est une fille. Reprenez l'arbre --- il faut maintenant y faire figurer le choix
de l'enfant que vous rencontrez --- et déterminez à nouveau la probabilité que les deux soient des
filles.
\item[(d)] Les deux réponses ne sont pas les mêmes. Expliquez, dans un paragraphe rédigé, ce qui a changé
entre (b) et (c), et pourquoi une énigme de probabilité n'est pas encore un problème tant qu'on n'a
pas dit \textit{comment} l'information a été obtenue.
\end{enumerate}

\end{problem}

\noindent Martin Gardner posa l'énigme dans sa chronique de \textit{Scientific American} en
1959, publia une réponse, puis fit machine arrière : la question, telle qu'il l'avait écrite, ne
déterminait pas la réponse. Ce n'est pas une faiblesse de l'énoncé, c'est le contenu même du
problème --- en probabilité, ce qui compte n'est pas le fait que vous apprenez mais la procédure qui
vous l'a fait apprendre. La version calculée avec des probabilités conditionnelles est PRB-23,
dans la bande Lycée ; ici, l'arbre suffit, et il rend la morale plus visible encore.

\begin{refs}
\item Contexte : Paradoxe des deux enfants --- Wikipédia (\url{https://fr.wikipedia.org/wiki/Paradoxe_des_deux_enfants})
\item Contexte : Arbre de probabilité --- Wikipédia (\url{https://fr.wikipedia.org/wiki/Arbre_de_probabilit%C3%A9})
\item Contexte : Martin Gardner --- Wikipédia (\url{https://fr.wikipedia.org/wiki/Martin_Gardner})
\end{refs}

\bigskip

\begin{problem}[{Trois dés qui tournent en rond --- Collège}]
\textbf{PRB-52.} \textbf{Problème.} Trois dés équilibrés à six faces portent les valeurs
$ A = (2,2,4,4,9,9) $, $ B = (1,1,6,6,8,8) $ et $ C = (3,3,5,5,7,7) $. Deux joueurs choisissent
chacun un dé, les lancent, et celui qui obtient le plus grand nombre gagne la partie.
\begin{enumerate}
\item[(a)] Pour chacun des trois duels --- A contre B, B contre C, puis C contre A --- dressez le tableau des
36 issues et déterminez, en justifiant, quel dé a le plus de chances de l'emporter.
\item[(b)] Comparez vos trois réponses. Énoncez avec précision l'idée reçue qu'elles mettent en défaut,
et expliquez en français pourquoi cette idée reçue semblait pourtant évidente.
\item[(c)] Vous proposez ce jeu à un camarade en le laissant choisir son dé le premier. Expliquez
pourquoi vous avez alors l'avantage, et dites ce qui change si c'est vous qui devez choisir en
premier.
\end{enumerate}

\end{problem}

\noindent Ces dés sont dus au statisticien Bradley Efron et furent rendus célèbres par
Martin Gardner dans \textit{Scientific American} en 1970. Leur intérêt n'est pas l'arithmétique mais
ce qu'ils détruisent : presque tout le monde suppose en silence que \og{} a plus de chances de gagner
contre \fg{} se transmet de proche en proche, comme \og{} est plus grand que \fg{} se transmet chez les
nombres. Cette hypothèse tacite sous-tend les classements sportifs, les comparatifs de produits et
certains modes de scrutin --- et voici un contre-exemple qui tient dans trois tableaux. La version à
quatre dés est PRB-24, dans la bande Lycée.

\begin{refs}
\item Contexte : Dés non transitifs --- Wikipédia (\url{https://fr.wikipedia.org/wiki/D%C3%A9s_non_transitifs})
\item Contexte : Martin Gardner --- Wikipédia (\url{https://fr.wikipedia.org/wiki/Martin_Gardner})
\end{refs}

\bigskip

\begin{problem}[{\og{} Deux issues, donc une chance sur deux \fg{} --- Collège, short}]
\textbf{PRB-53.} \textbf{Problème.} Un élève affirme : \og{} quand je joue au loto, soit je gagne,
soit je perds ; il n'y a que deux issues, donc j'ai une chance sur deux de gagner \fg{}. Débusquez
l'erreur en énonçant précisément l'hypothèse qui manque, puis construisez deux expériences n'ayant
chacune que deux issues, l'une où \og{} une chance sur deux \fg{} est juste et l'autre où c'est faux, et
justifiez les deux cas.
\end{problem}

\noindent L'erreur a une lettre de noblesse. Dans l'article \og{} Croix ou pile \fg{} de
l'\textit{Encyclopédie} (1754), d'Alembert soutient qu'en deux lancers d'une pièce la probabilité
d'obtenir au moins une fois croix vaut $ 2/3 $, parce qu'il ne voit que trois cas : croix au
premier coup, pile puis croix, pile puis pile. Il compte des descriptions d'issues au lieu de
compter des issues également probables --- la faute même de l'élève ci-dessus, commise par l'un des
plus grands mathématiciens du siècle. L'hypothèse d'équiprobabilité n'est jamais gratuite : elle se
justifie ou elle se démontre, jamais elle ne se suppose au vu du nombre de cases.

\begin{refs}
\item Contexte : Équiprobabilité --- Wikipédia (\url{https://fr.wikipedia.org/wiki/%C3%89quiprobabilit%C3%A9})
\item Contexte : Jean Le Rond d'Alembert --- Wikipédia (\url{https://fr.wikipedia.org/wiki/Jean_Le_Rond_d%27Alembert})
\item Contexte : Loterie --- Wikipédia (\url{https://fr.wikipedia.org/wiki/Loterie})
\end{refs}

\bigskip

\begin{problem}[{Deux urnes et un dé --- Collège}]
\textbf{PRB-54.} \textbf{Problème.} L'urne n$^\circ$ 1 contient 3 boules rouges et 2 boules vertes ; l'urne n$^\circ$ 2 contient
1 boule rouge et 4 boules vertes. On lance un dé équilibré : si l'on obtient 5 ou 6, on tire une
boule dans l'urne n$^\circ$ 1, sinon on tire une boule dans l'urne n$^\circ$ 2.
\begin{enumerate}
\item[(a)] Dressez l'arbre de cette expérience en portant sur chaque branche la probabilité
correspondante, et justifiez que la somme des probabilités des branches issues d'un même nœud vaut
toujours 1.
\item[(b)] Déterminez la probabilité de tirer une boule rouge, en justifiant chaque étape de votre
calcul sur l'arbre.
\item[(c)] Un camarade objecte : \og{} il y a en tout 4 boules rouges parmi les 10 boules des deux urnes,
donc la probabilité de tirer une rouge vaut $ 4/10 $ \fg{}. Expliquez pourquoi ce raisonnement est
faux ici, puis décrivez précisément une règle de tirage --- différente de celle de l'énoncé --- pour
laquelle il serait juste.
\end{enumerate}

\end{problem}

\noindent Les urnes remplies de boules de couleur sont le laboratoire de la théorie des
probabilités depuis Jacques Bernoulli, dont l'\textit{Ars conjectandi} (publié en 1713) en fait
l'instrument de base du raisonnement sur le hasard. Leur intérêt est d'être totalement
transparentes : on voit ce qu'il y a dedans, il ne reste donc rien à deviner et tout à démontrer.
La question (c) contient la seule difficulté véritable --- une probabilité n'est une fraction de cas
favorables que lorsque les cas sont également probables, et deux urnes tirées avec des chances
inégales ne se laissent pas verser dans un même sac.

\begin{refs}
\item Contexte : Problème d'urne --- Wikipédia (\url{https://fr.wikipedia.org/wiki/Probl%C3%A8me_d%27urne})
\item Contexte : Arbre de probabilité --- Wikipédia (\url{https://fr.wikipedia.org/wiki/Arbre_de_probabilit%C3%A9})
\end{refs}

\bigskip

\begin{problem}[{Au moins un six --- Collège, short}]
\textbf{PRB-55.} \textbf{Problème.} Un élève raisonne ainsi sur le lancer de deux dés
équilibrés : \og{} la probabilité d'obtenir un six sur un dé vaut $ 1/6 $, donc avec deux dés elle
vaut $ 1/6 + 1/6 = 1/3 $ \fg{}. Déterminez la probabilité d'obtenir au moins un six en la justifiant
sur le tableau des 36 issues, puis énoncez la condition sous laquelle additionner deux probabilités
est légitime et montrez qu'elle n'est pas remplie ici.
\end{problem}

\noindent L'addition naïve des probabilités est l'erreur la plus répandue de tout
l'enseignement de la discipline, et elle a résisté longtemps : les premiers auteurs --- Cardan au
XVIe siècle dans son \textit{Liber de ludo aleae}, Galilée vers 1620 --- ont dû forger la
notion de \og{} cas également possibles \fg{} précisément pour se donner un garde-fou contre ce genre de
glissement. Le point à retenir dépasse les dés : additionner des probabilités revient à supposer
que deux choses ne peuvent pas arriver ensemble, et c'est une hypothèse sur le monde, pas une règle
de calcul.

\begin{refs}
\item Contexte : Événement (probabilités) --- Wikipédia (\url{https://fr.wikipedia.org/wiki/%C3%89v%C3%A9nement_(probabilit%C3%A9s)})
\item Contexte : Jérôme Cardan --- Wikipédia (\url{https://fr.wikipedia.org/wiki/J%C3%A9r%C3%B4me_Cardan})
\end{refs}

\bigskip

\begin{problem}[{La pièce qui \og{} doit \fg{} tomber sur pile --- Collège, short}]
\textbf{PRB-56.} \textbf{Problème.} Une pièce équilibrée vient de tomber cinq fois de suite sur
face, et un joueur en conclut : \og{} pile est en retard, il a maintenant plus d'une chance sur deux de
sortir au prochain lancer \fg{}. Réfutez ce raisonnement, puis expliquez comment il peut être faux
alors même qu'il est vrai que, sur un très grand nombre de lancers, la fréquence de \og{} pile \fg{} se
rapproche de $ 1/2 $.
\end{problem}

\noindent Laplace décrit déjà cette illusion dans son \textit{Essai philosophique sur les
probabilités} (1814) : il y raconte des joueurs de loterie qui misent sur les numéros non sortis
depuis longtemps, persuadés qu'ils sont \og{} dus \fg{}, et il en fait l'exemple type des erreurs de
jugement sur le hasard. Deux siècles plus tard, les sites de loterie publient encore la liste des
numéros \og{} en retard \fg{}, et les joueurs les jouent. La seconde partie de la question est la vraie :
la fréquence se rapproche bien de $ 1/2 $, mais elle y arrive en \textit{diluant} les écarts
passés, jamais en les corrigeant --- la version chiffrée de cette distinction est PRB-05, dans la
bande Lycée.

\begin{refs}
\item Contexte : Erreur du parieur --- Wikipédia (\url{https://fr.wikipedia.org/wiki/Erreur_du_parieur})
\item Contexte : Pierre-Simon de Laplace --- Wikipédia (\url{https://fr.wikipedia.org/wiki/Pierre-Simon_de_Laplace})
\item Contexte : Loi des grands nombres --- Wikipédia (\url{https://fr.wikipedia.org/wiki/Loi_des_grands_nombres})
\end{refs}

\bigskip


\section*{Statistique}

\begin{problem}[{Le salaire moyen chez Duval --- Collège}]
\textbf{STA-41.} \textbf{Problème.} L'entreprise Duval emploie dix personnes : huit ouvriers payés 1 500 € par
mois, un chef d'atelier payé 2 400 € et le patron, qui se verse 21 600 €.
\begin{enumerate}
\item[(a)] Calculez la moyenne et la médiane de ces dix salaires. Justifiez chaque calcul, et en
particulier la façon dont vous déterminez la médiane d'une série de dix valeurs.
\item[(b)] L'entreprise affiche sur son site le salaire moyen de ses employés. Expliquez pourquoi cette
information, bien qu'exacte, donne une image fausse de l'entreprise. Dites quel indicateur vous
publieriez à sa place et défendez votre choix.
\item[(c)] Le patron double son propre salaire, les neuf autres restant inchangés. Dites ce que
deviennent la moyenne, la médiane et l'étendue, en justifiant les trois réponses. Énoncez ensuite,
en une phrase précise, ce que cet exemple révèle sur la sensibilité de chacun de ces trois
indicateurs aux valeurs extrêmes.
\end{enumerate}

\end{problem}

\noindent Le désaccord entre moyenne et médiane n'est pas un artifice de manuel : c'est la
raison pour laquelle l'INSEE publie systématiquement les deux quand il décrit les salaires en
France, le salaire médian y étant nettement inférieur au salaire moyen. L'écart entre les deux
nombres est lui-même une information --- il mesure à quel point la distribution est tirée vers le
haut par une minorité. La moyenne appartient à la tradition des moindres carrés de Legendre et
Gauss ; la médiane fut défendue par Laplace, qui démontra qu'elle minimise l'erreur absolue. La
démonstration de cette propriété est STA-01, dans la bande Lycée.

\begin{refs}
\item Contexte : Moyenne arithmétique --- Wikipédia (\url{https://fr.wikipedia.org/wiki/Moyenne_arithm%C3%A9tique})
\item Contexte : Médiane (statistiques) --- Wikipédia (\url{https://fr.wikipedia.org/wiki/M%C3%A9diane_(statistiques)})
\item Contexte : Statistique descriptive --- Wikipédia (\url{https://fr.wikipedia.org/wiki/Statistique_descriptive})
\end{refs}

\bigskip

\begin{problem}[{L'axe qui ne part pas de zéro --- Collège}]
\textbf{STA-42.} \textbf{Problème.} Un journal local publie un diagramme en barres du nombre de cambriolages
enregistrés dans la ville sur cinq années consécutives : 1 010, 1 025, 1 040, 1 030 et 1 060.
L'axe vertical du graphique commence à 1 000 et s'arrête à 1 070. Le titre de l'article est \og{} La
délinquance flambe \fg{}.
\begin{enumerate}
\item[(a)] Calculez l'augmentation entre la première et la dernière année, en nombre de cambriolages puis
en pourcentage. Justifiez vos deux calculs et dites lequel des deux vous paraît le plus honnête à
publier.
\item[(b)] Tracez le même diagramme avec un axe vertical partant de 0. Décrivez précisément ce qui change
dans l'impression visuelle, et expliquez pourquoi les deux graphiques sont pourtant tous les deux
exacts.
\item[(c)] Un axe tronqué n'est pas toujours malhonnête. Donnez un exemple de série de données pour
laquelle partir de zéro serait au contraire une mauvaise idée, puis rédigez la règle que vous
proposeriez à ce journal pour distinguer les deux situations.
\end{enumerate}

\end{problem}

\noindent Darrell Huff a consacré à ce procédé un chapitre entier de \textit{How to Lie with
Statistics} (1954), l'un des livres de statistique les plus vendus de l'histoire ; il l'appelait
le graphique \og{} gee-whiz \fg{}, du nom de l'exclamation qu'il cherche à provoquer. Le point délicat, et
c'est celui de la question (c), est qu'aucune règle mécanique ne fonctionne : un axe partant de
zéro rendrait illisible une courbe de températures corporelles, et personne ne songerait à le
reprocher. Ce qui distingue l'honnêteté du mensonge n'est donc pas la position de l'axe mais
l'accord entre ce que le graphique fait voir et ce que le texte affirme.

\begin{refs}
\item Contexte : Visualisation de données --- Wikipédia (\url{https://fr.wikipedia.org/wiki/Visualisation_de_donn%C3%A9es})
\item Contexte : Diagramme à barres --- Wikipédia (\url{https://fr.wikipedia.org/wiki/Diagramme_%C3%A0_barres})
\item Contexte : Graphique trompeur --- Wikipédia (en anglais) (\url{https://en.wikipedia.org/wiki/Misleading_graph})
\end{refs}

\bigskip

\begin{problem}[{Le sondage de la cour de récréation --- Collège}]
\textbf{STA-43.} \textbf{Problème.} Camille veut connaître le sport préféré des 600 élèves de son collège. Un
mercredi à midi, elle interroge les 40 élèves présents sur le terrain de basket, dont 28 répondent
\og{} le basket \fg{}, et elle conclut : \og{} 70 \% des élèves du collège préfèrent le basket \fg{}.
\begin{enumerate}
\item[(a)] Énoncez au moins trois raisons distinctes pour lesquelles cette conclusion n'est pas fondée.
Pour chacune, dites exactement quelle étape du raisonnement de Camille elle attaque.
\item[(b)] Camille répond : \og{} j'ai interrogé 40 élèves, c'est beaucoup \fg{}. Expliquez pourquoi interroger
400 élèves au même endroit et au même moment ne corrigerait pas le défaut principal de son
enquête. Votre explication doit porter sur la façon dont les élèves interrogés ont été choisis.
\item[(c)] Concevez un protocole de sondage réellement applicable dans votre collège. Décrivez-le en cinq
lignes, justifiez chacun de vos choix, et terminez en indiquant ce que votre protocole ne garantit
toujours pas.
\end{enumerate}

\end{problem}

\noindent L'histoire des sondages est une suite d'échecs instructifs. En 1948, tous les
grands instituts américains annoncèrent la défaite de Harry Truman face à Thomas Dewey ; le
\textit{Chicago Tribune} imprima même sa une \og{} Dewey Defeats Truman \fg{} avant le dépouillement, et la
photographie de Truman brandissant le journal est restée célèbre. Les instituts avaient cessé de
sonder trop tôt et travaillaient sur des échantillons construits par quotas mal actualisés. Le cas
plus ancien et plus spectaculaire du \textit{Literary Digest}, qui recueillit deux millions de
réponses et se trompa quand même, est traité en STA-32 dans la bande Lycée : la leçon commune est
que la manière de choisir les personnes interrogées compte davantage que leur nombre.

\begin{refs}
\item Contexte : Sondage d'opinion --- Wikipédia (\url{https://fr.wikipedia.org/wiki/Sondage_d%27opinion})
\item Contexte : Biais de sélection --- Wikipédia (\url{https://fr.wikipedia.org/wiki/Biais_de_s%C3%A9lection})
\item Contexte : Dewey Defeats Truman --- Wikipédia (\url{https://fr.wikipedia.org/wiki/Dewey_Defeats_Truman})
\end{refs}

\bigskip

\begin{problem}[{Deux classes, une seule moyenne --- Collège}]
\textbf{STA-44.} \textbf{Problème.} Deux classes de neuf élèves ont passé le même contrôle, noté sur 20. Voici les
notes obtenues.

\noindent Classe A : 9, 10, 10, 11, 11, 11, 12, 12, 13 
Classe B : 3, 4, 5, 11, 13, 15, 16, 16, 16
\begin{enumerate}
\item[(a)] Vérifiez que les deux classes ont exactement la même moyenne, puis calculez pour chacune la
médiane et l'étendue. Justifiez chaque calcul.
\item[(b)] Le professeur doit dire laquelle des deux classes \og{} a mieux réussi \fg{}. Construisez d'abord
l'argument le plus solide en faveur de la classe A, puis l'argument le plus solide en faveur de la
classe B, chaque fois chiffres à l'appui. Concluez sur ce que la question posée a d'incomplet.
\item[(c)] Proposez un indicateur simple, calculable au collège, qui distinguerait ces deux séries mieux
que la moyenne. Justifiez qu'il les distingue effectivement, puis donnez un exemple de deux séries
que votre indicateur ne distingue pas non plus.
\end{enumerate}

\end{problem}

\noindent Une moyenne seule ne décrit jamais une série : c'est la première leçon de la
statistique descriptive, et elle a mis un siècle à se formaliser. On a longtemps mesuré la
dispersion par l'écart moyen ; c'est Karl Pearson qui, en 1893, a introduit le terme d'écart type
et imposé peu à peu la mesure que vous rencontrerez au lycée. Le choix entre plusieurs mesures de
dispersion n'a rien d'évident et fut sérieusement débattu --- la question (c) vous fait refaire ce
débat en miniature, y compris sa partie la plus honnête : reconnaître ce que votre indicateur
laisse encore passer.

\begin{refs}
\item Contexte : Série statistique --- Wikipédia (\url{https://fr.wikipedia.org/wiki/S%C3%A9rie_statistique})
\item Contexte : Médiane (statistiques) --- Wikipédia (\url{https://fr.wikipedia.org/wiki/M%C3%A9diane_(statistiques)})
\item Contexte : Karl Pearson --- Wikipédia (\url{https://fr.wikipedia.org/wiki/Karl_Pearson})
\end{refs}

\bigskip

\begin{problem}[{La moyenne des moyennes --- Collège}]
\textbf{STA-45.} \textbf{Problème.} Au dernier contrôle de mathématiques, la classe de 4e A, qui compte
28 élèves, a obtenu 12 de moyenne ; la classe de 4e B, qui en compte 12, a obtenu 16.
\begin{enumerate}
\item[(a)] Le principal annonce que \og{} la moyenne des deux classes réunies est 14 \fg{}. Montrez que c'est
faux, calculez la valeur correcte et justifiez votre méthode.
\item[(b)] Énoncez précisément la condition sous laquelle la moyenne des moyennes coïncide avec la
moyenne de l'ensemble. Démontrez-la dans le cas de deux groupes d'effectifs quelconques.
\item[(c)] Un élève a obtenu 8 de moyenne au premier trimestre, où il avait trois notes, et 16 au
deuxième, où il n'en avait qu'une. Il annonce 12 de moyenne sur les deux trimestres. Dites ce qui
est vrai et ce qui est faux dans son calcul, et déterminez en justifiant la valeur correcte.
\end{enumerate}

\end{problem}

\noindent Confondre la moyenne des moyennes et la moyenne globale est l'erreur
arithmétique la plus fréquemment publiée, et ce n'est pas une simple étourderie : c'est le
mécanisme même du paradoxe de Simpson, où deux groupes peuvent chacun avoir un meilleur taux de
réussite qu'un concurrent tout en ayant, réunis, un taux inférieur. La démonstration complète est
STA-02, dans la bande Lycée. Le système français des coefficients sur les bulletins scolaires est
d'ailleurs une moyenne pondérée déguisée : la question (c) montre que même sans coefficient
affiché, les effectifs jouent déjà ce rôle.

\begin{refs}
\item Contexte : Moyenne pondérée --- Wikipédia (\url{https://fr.wikipedia.org/wiki/Moyenne_pond%C3%A9r%C3%A9e})
\item Contexte : Paradoxe de Simpson --- Wikipédia (\url{https://fr.wikipedia.org/wiki/Paradoxe_de_Simpson})
\end{refs}

\bigskip

\begin{problem}[{Le diagramme circulaire du foyer --- Collège}]
\textbf{STA-46.} \textbf{Problème.} Le foyer socio-éducatif d'un collège publie un diagramme circulaire des
activités choisies par ses 240 inscrits. Les angles au centre des secteurs sont : théâtre 90$^\circ$,
échecs 60$^\circ$, journal 45$^\circ$, robotique 120$^\circ$, autres activités 45$^\circ$.
\begin{enumerate}
\item[(a)] Retrouvez l'effectif de chaque activité et justifiez la méthode de conversion que vous
employez pour passer d'un angle à un effectif. Vérifiez votre résultat.
\item[(b)] L'année suivante, le foyer compte 300 inscrits et le secteur \og{} robotique \fg{} occupe toujours
120$^\circ$. Le président en conclut que la robotique n'a pas progressé. Réfutez-le en distinguant
soigneusement effectif et fréquence, et dites ce que le diagramme permet malgré tout
d'affirmer.
\item[(c)] Le président souhaite mettre en valeur le journal du collège et propose de dessiner le
diagramme en perspective, le disque étant incliné. Expliquez géométriquement pourquoi cette
présentation fausse la lecture, et déterminez quels secteurs elle avantage.
\end{enumerate}

\end{problem}

\noindent Le diagramme circulaire est une invention de l'Écossais William Playfair, qui en
publia le premier exemple connu dans son \textit{Statistical Breviary} de 1801 ; on lui doit aussi le
diagramme en barres et la courbe chronologique, c'est-à-dire l'essentiel du vocabulaire graphique
que l'on emploie encore. Il a mauvaise réputation chez les spécialistes : Edward Tufte, dans
\textit{The Visual Display of Quantitative Information}, soutient que l'œil compare mal les angles et
que la seule chose pire qu'un diagramme circulaire est plusieurs diagrammes circulaires côte à
côte. La question (c) explique la moitié de cette mauvaise réputation.

\begin{refs}
\item Contexte : Diagramme circulaire --- Wikipédia (\url{https://fr.wikipedia.org/wiki/Diagramme_circulaire})
\item Contexte : William Playfair --- Wikipédia (\url{https://fr.wikipedia.org/wiki/William_Playfair})
\item Contexte : Edward Tufte --- Wikipédia (\url{https://fr.wikipedia.org/wiki/Edward_Tufte})
\end{refs}

\bigskip

\begin{problem}[{Effectif ou fréquence ? --- Collège, short}]
\textbf{STA-47.} \textbf{Problème.} Dans un collège de 800 élèves, 200 sont demi-pensionnaires ;
dans le collège voisin, 300 élèves sur 1 500 le sont. Un journaliste écrit : \og{} le second collège
compte plus de demi-pensionnaires, la demi-pension y est donc plus populaire \fg{} --- dites, en
justifiant par les effectifs puis par les fréquences, ce que cette phrase confond, et rédigez la
comparaison correcte.
\end{problem}

\noindent C'est pour cette raison que les statistiques officielles ne publient presque
jamais un effectif seul : l'INSEE et les agences sanitaires donnent des taux --- pour mille
habitants, pour cent mille --- précisément parce qu'un nombre brut mélange deux informations, la
grandeur du phénomène et la taille de la population où on le mesure. Un nombre de cas qui augmente
dans une population qui augmente ne dit rien tant qu'on n'a pas divisé. Séparer \og{} combien \fg{} de
\og{} quelle part \fg{} est le premier réflexe de lecture statistique, et l'un des plus souvent oubliés.

\begin{refs}
\item Contexte : Fréquence (statistiques) --- Wikipédia (\url{https://fr.wikipedia.org/wiki/Fr%C3%A9quence_(statistiques)})
\item Contexte : Pourcentage --- Wikipédia (\url{https://fr.wikipedia.org/wiki/Pourcentage})
\end{refs}

\bigskip

\begin{problem}[{Un pourcentage d'un pourcentage --- Collège, short}]
\textbf{STA-48.} \textbf{Problème.} Dans un collège, 40 \% des élèves sont en cycle 4 et, parmi
ces derniers, 30 \% apprennent l'allemand ; un délégué en conclut que \og{} 70 \% des élèves du collège
apprennent l'allemand en cycle 4 \fg{}. Débusquez l'erreur, déterminez en le justifiant le pourcentage
correct, et expliquez pourquoi la réponse ne dépend pas de l'effectif total du collège.
\end{problem}

\noindent Prendre un pourcentage d'un pourcentage, c'est multiplier des proportions, et
non les additionner --- une règle que l'arithmétique commerciale maîtrisait déjà bien avant la
statistique. Fibonacci l'applique aux intérêts et aux changes dans le \textit{Liber abaci} (1202), et
les écoles d'abaque de Toscane l'enseignaient aux futurs marchands du XIVe siècle. La
même erreur produit son cousin le plus célèbre : une hausse de 20 \% suivie d'une baisse de 20 \% ne
ramène pas au point de départ, et la question ci-dessus contient déjà la raison de ce
phénomène.

\begin{refs}
\item Contexte : Pourcentage --- Wikipédia (\url{https://fr.wikipedia.org/wiki/Pourcentage})
\item Contexte : Fréquence (statistiques) --- Wikipédia (\url{https://fr.wikipedia.org/wiki/Fr%C3%A9quence_(statistiques)})
\end{refs}

\bigskip

\begin{problem}[{La médiane d'un tableau d'effectifs --- Collège, short}]
\textbf{STA-49.} \textbf{Problème.} Dans une classe de 27 élèves, on a relevé le nombre de frères
et sœurs de chacun : 4 élèves en ont 0, 9 en ont 1, 7 en ont 2, 4 en ont 3, 2 en ont 4 et 1 en a 6.
Déterminez la médiane et la moyenne de cette série en justifiant votre méthode directement sur le
tableau d'effectifs, sans réécrire les 27 valeurs une à une, puis expliquez pourquoi la moyenne
n'est pas un nombre entier alors qu'aucun élève ne peut avoir un nombre non entier de frères et
sœurs.
\end{problem}

\noindent Le tableau d'effectifs est la forme sous laquelle les données arrivent
réellement --- un recensement ne vous remet pas une liste de soixante-sept millions de lignes, il
vous remet des comptages. Savoir lire une médiane et une moyenne directement dans un tableau
d'effectifs cumulés est donc une compétence de praticien avant d'être un exercice. Quant à la
dernière question, elle touche à quelque chose de profond : une moyenne n'est pas un individu de la
série, c'est un point d'équilibre, et rien n'oblige un point d'équilibre à exister parmi les
données.

\begin{refs}
\item Contexte : Médiane (statistiques) --- Wikipédia (\url{https://fr.wikipedia.org/wiki/M%C3%A9diane_(statistiques)})
\item Contexte : Série statistique --- Wikipédia (\url{https://fr.wikipedia.org/wiki/S%C3%A9rie_statistique})
\end{refs}

\bigskip

\begin{problem}[{Le pictogramme qui double --- Collège, short}]
\textbf{STA-50.} \textbf{Problème.} Pour illustrer que le budget du club de sport a doublé, un
graphiste dessine deux ballons dont le second a un diamètre deux fois plus grand que le premier.
Déterminez, calculs à l'appui, par quel facteur l'image exagère la hausse --- en distinguant le cas
où le lecteur compare des aires de disque et celui où il compare des volumes de ballon --- puis
proposez une construction correcte et justifiez-la.
\end{problem}

\noindent Darrell Huff consacre un chapitre de \textit{How to Lie with Statistics} (1954) à
ce procédé, qu'il tenait pour l'un des plus efficaces parce que le graphique reste, au sens strict,
véridique : le second dessin est bien deux fois plus grand que le premier \textit{dans une dimension}.
Le lecteur, lui, ne mesure pas des diamètres, il perçoit des surfaces ou des volumes, et l'écart
entre les deux est précisément ce que la question vous demande de chiffrer. C'est le seul exercice
de cette page où la géométrie et la statistique se rencontrent --- et ce n'est pas un hasard :
tromper avec un graphique, c'est presque toujours exploiter une différence entre ce qui est mesuré
et ce qui est vu.

\begin{refs}
\item Contexte : Pictogramme --- Wikipédia (\url{https://fr.wikipedia.org/wiki/Pictogramme})
\item Contexte : How to Lie with Statistics --- Wikipédia (en anglais) (\url{https://en.wikipedia.org/wiki/How_to_Lie_with_Statistics})
\item Contexte : Graphique trompeur --- Wikipédia (en anglais) (\url{https://en.wikipedia.org/wiki/Misleading_graph})
\end{refs}

\bigskip

\begin{problem}[{Combien d'élèves peuvent être au-dessus de la moyenne ? --- Collège}]
\textbf{STA-51.} \textbf{Problème.} On considère les notes d'un contrôle dans une classe de 25 élèves.
\begin{enumerate}
\item[(a)] Est-il possible que 24 des 25 élèves aient une note strictement supérieure à la moyenne de la
classe ? Construisez un exemple explicite, ou démontrez que c'est impossible.
\item[(b)] Démontrez qu'il est en revanche impossible que les 25 élèves soient tous strictement au-dessus
de la moyenne de la classe. Votre démonstration doit porter sur la somme des notes, et rester
valable pour un nombre quelconque d'élèves.
\item[(c)] Un journal titre : \og{} la moitié des collèges français sont en dessous de la moyenne nationale \fg{}.
Expliquez ce que ce titre a de vide, dites ce qu'il faudrait mesurer à la place pour que
l'information ait un contenu, et distinguez soigneusement le rôle de la moyenne et celui de la
médiane dans votre réponse.
\end{enumerate}

\end{problem}

\noindent L'idée que presque tout le monde peut être au-dessus de la moyenne porte un nom
en anglais, l'effet Lake Wobegon, d'après la ville imaginaire de l'humoriste Garrison Keillor \og{} où
tous les enfants sont au-dessus de la moyenne \fg{}. Elle n'est pas seulement drôle : le psychologue Ola
Svenson a publié en 1981 une enquête devenue classique où une large majorité d'automobilistes se
jugeaient plus prudents et plus habiles que la moyenne des conducteurs. La question (a) montre que
cela n'a rien d'impossible arithmétiquement, la question (b) dit où se trouve la vraie limite, et
la question (c) explique pourquoi confondre moyenne et médiane transforme une banalité
arithmétique en scandale de presse.

\begin{refs}
\item Contexte : Moyenne arithmétique --- Wikipédia (\url{https://fr.wikipedia.org/wiki/Moyenne_arithm%C3%A9tique})
\item Contexte : Supériorité illusoire --- Wikipédia (\url{https://fr.wikipedia.org/wiki/Sup%C3%A9riorit%C3%A9_illusoire})
\item Article : O. Svenson, \og{} Are we all less risky and more skillful than our fellow drivers? \fg{}, \textit{Acta Psychologica} 47 (1981), 143--148
\end{refs}

\bigskip

\begin{problem}[{Trois pourcentages dans la même case --- Collège}]
\textbf{STA-52.} \textbf{Problème.} On a interrogé 400 élèves d'un collège, dont 240 filles et 160 garçons, en
leur demandant s'ils pratiquent un sport en club. Parmi les filles, 156 répondent oui ; parmi les
garçons, 112.
\begin{enumerate}
\item[(a)] Recopiez et complétez le tableau à double entrée croisant le sexe et la pratique d'un sport en
club, marges comprises. Justifiez chaque case que vous ajoutez.
\item[(b)] À partir de la seule case \og{} filles pratiquant un sport en club \fg{}, on peut former les trois
quotients 156 sur 240, 156 sur 268 et 156 sur 400. Énoncez, pour chacun, la question précise à
laquelle il répond. Expliquez ensuite pourquoi confondre ces trois nombres est l'erreur la plus
courante dans la lecture d'un tel tableau.
\item[(c)] Le journal du collège titre : \og{} les filles font moins de sport que les garçons \fg{}. Dites ce que
les données permettent d'affirmer et ce qu'elles ne permettent pas, en distinguant explicitement
effectif et fréquence, puis nommez au moins une information manquante qui pourrait changer votre
jugement.
\end{enumerate}

\end{problem}

\noindent Le tableau à double entrée est le plus vieil instrument de la statistique
appliquée : John Graunt s'en sert déjà dans ses \textit{Natural and Political Observations Made upon the
Bills of Mortality} (1662), le premier travail de statistique démographique, pour croiser causes
de décès et années. Sa simplicité est trompeuse --- la confusion entre les trois quotients de la
question (b) est la racine commune de la plupart des erreurs d'interprétation, depuis les tests
médicaux mal lus jusqu'au paradoxe de Simpson (STA-02, bande Lycée). Apprendre à dire à voix haute
\og{} sur quoi je divise, et pourquoi \fg{} règle une bonne partie de ces malentendus.

\begin{refs}
\item Contexte : Tableau de contingence --- Wikipédia (\url{https://fr.wikipedia.org/wiki/Tableau_de_contingence})
\item Contexte : John Graunt --- Wikipédia (\url{https://fr.wikipedia.org/wiki/John_Graunt})
\item Contexte : Fréquence (statistiques) --- Wikipédia (\url{https://fr.wikipedia.org/wiki/Fr%C3%A9quence_(statistiques)})
\end{refs}

\bigskip


\section*{Mathématiques discrètes et informatique}

\begin{problem}[{Trier un jeu de cartes, et démontrer que ça marche --- Collège}]
\textbf{DCS-51.} \textbf{Problème.} Posez devant vous sept cartes de valeurs toutes différentes, faces
visibles, dans un ordre quelconque. Voici une méthode pour les ranger : parcourez toutes les
cartes non encore rangées, repérez la plus petite, échangez-la avec la première carte non
rangée ; puis recommencez sur ce qui reste.
\begin{enumerate}
\item[(a)] Déroulez la méthode à la main sur les valeurs 5, 3, 9, 1, 7, 2, 8, en écrivant l'état
complet de la rangée après chaque échange.
\item[(b)] Écrivez la phrase qui décrit ce que contient le début de la rangée et qui reste vraie
après chaque échange. Démontrez qu'elle est encore vraie après l'échange suivant, et qu'à la
fin elle dit exactement \og{} la rangée est triée \fg{}.
\item[(c)] Comptez le nombre de comparaisons faites sur sept cartes, puis sur n cartes. Donnez la
formule et démontrez-la.
\item[(d)] Voici une seconde méthode : prenez les cartes une par une et glissez chacune à sa place
parmi celles déjà tenues en main. Déroulez-la sur la même liste, puis comparez les deux
méthodes sur une liste déjà triée, et sur une liste triée à l'envers.
\end{enumerate}

\end{problem}

\noindent La seconde méthode est celle que tout joueur de cartes exécute sans y
penser depuis des siècles ; la première est celle qu'un débutant réinvente au tableau. Elles
portent aujourd'hui les noms de tri par insertion et de tri par sélection, et elles étaient
déjà l'objet des tout premiers programmes écrits : le tri par fusion figure dans un
manuscrit de John von Neumann de 1945 pour l'EDVAC, et Donald Knuth a consacré au seul
problème du tri la moitié d'un volume de \textit{The Art of Computer Programming}. Ce que ce
problème demande n'est pas de trier --- vous savez trier --- mais de \textit{démontrer} que la
méthode trie, et la différence entre les deux est exactement le sujet de cette page. La
question (b) est le premier invariant de boucle de votre vie ; DCS-02 en donne la version
lycée.

\begin{refs}
\item Contexte : Tri par sélection --- Wikipédia (\url{https://fr.wikipedia.org/wiki/Tri_par_s%C3%A9lection})
\item Contexte : Tri par insertion --- Wikipédia (\url{https://fr.wikipedia.org/wiki/Tri_par_insertion})
\item Contexte : Invariant de boucle --- Wikipédia (\url{https://fr.wikipedia.org/wiki/Invariant_de_boucle})
\end{refs}

\bigskip

\begin{problem}[{Devine mon nombre en sept questions --- Collège}]
\textbf{DCS-52.} \textbf{Problème.} Votre adversaire choisit en secret un entier compris entre 1 et 100. Vous
lui posez des questions auxquelles il ne répond que par \og{} oui \fg{} ou par \og{} non \fg{}, et il ne ment
jamais.
\begin{enumerate}
\item[(a)] Décrivez une méthode qui trouve son nombre à coup sûr en sept questions au plus. Elle
doit être assez précise pour qu'une autre personne puisse l'appliquer sans vous.
\item[(b)] Démontrez qu'elle réussit toujours --- pour les cent nombres possibles, et pas seulement
pour ceux que vous avez essayés.
\item[(c)] Démontrez qu'aucune méthode, si maligne soit-elle, ne peut garantir la victoire en six
questions. (Forme suggérée : comptez ce que six réponses peuvent au plus distinguer.)
\item[(d)] Jusqu'à quel nombre pourriez-vous aller avec sept questions ? avec dix ? Démontrez vos
réponses. Puis dites ce qui change si l'adversaire a le droit de mentir une fois.
\end{enumerate}

\end{problem}

\noindent La question (c) est le cœur de l'affaire, et elle est d'une nature
entièrement différente de la (a) : elle ne parle pas de votre méthode, elle parle de toutes
les méthodes possibles, y compris celles que personne n'a encore imaginées. C'est ce qu'on
appelle une borne inférieure, et savoir en démontrer une est ce qui sépare celui qui
programme de celui qui comprend. La variante avec mensonge de la question (d) est un vrai
sujet de recherche, posé indépendamment par Alfréd Rényi (1961) et Stanisław Ulam (1976) ;
c'est le point de départ des codes correcteurs d'erreurs, où l'on cherche justement à
retrouver la vérité malgré des réponses fausses. DCS-53 vous fera reconnaître la même
méthode dans un geste que vous faites depuis toujours.

\begin{refs}
\item Contexte : Recherche dichotomique --- Wikipédia (\url{https://fr.wikipedia.org/wiki/Recherche_dichotomique})
\item Contexte : Puissance de deux --- Wikipédia (\url{https://fr.wikipedia.org/wiki/Puissance_de_deux})
\item Contexte : Le jeu des vingt questions --- Wikipédia (en anglais) (\url{https://en.wikipedia.org/wiki/Twenty_questions})
\end{refs}

\bigskip

\begin{problem}[{Chercher un mot dans le dictionnaire --- Collège, short}]
\textbf{DCS-53.} \textbf{Problème.} Décrivez, assez précisément pour qu'un autre puisse
l'appliquer sans vous, la méthode que vous employez réellement pour trouver un mot dans un
dictionnaire de papier. Démontrez ensuite qu'elle s'arrête toujours --- soit en trouvant le
mot, soit en établissant qu'il n'y figure pas --- et majorez le nombre de fois qu'elle vous
fait ouvrir le volume dans un dictionnaire de 60 000 entrées.
\end{problem}

\noindent Presque personne ne feuillette un dictionnaire page à page, et presque
personne n'a jamais formulé la méthode qu'il emploie à la place : c'est un algorithme que
l'on exécute des années durant sans l'avoir écrit une seule fois. Le mettre en mots est
l'exercice, et la majoration demandée est la surprise --- le nombre est ridiculement petit,
et c'est cette petitesse qui explique pourquoi la même idée fait fonctionner les moteurs de
recherche. DCS-02 en donne la version formelle, avec invariant de boucle et bogue de
débordement à la clé.

\begin{refs}
\item Contexte : Recherche dichotomique --- Wikipédia (\url{https://fr.wikipedia.org/wiki/Recherche_dichotomique})
\item Voir aussi : DCS-02, la correction de la recherche dichotomique (\url{#DCS-02})
\end{refs}

\bigskip

\begin{problem}[{Une consigne qu'un robot peut suivre --- Collège}]
\textbf{DCS-54.} \textbf{Problème.} Un exécutant ne comprend que quatre ordres : \og{} avance d'un pas \fg{},
\og{} tourne d'un quart de tour vers la droite \fg{}, \og{} pose un jeton \fg{}, et \og{} répète N fois : \dots{} \fg{}.
Il n'a aucun bon sens : il fait exactement ce qui est écrit, et rien d'autre.
\begin{enumerate}
\item[(a)] Écrivez la suite d'ordres qui lui fait parcourir le contour d'un carré de côté quatre
pas. Démontrez qu'il revient à sa position de départ \textit{et} à son orientation de
départ.
\item[(b)] Écrivez la suite d'ordres qui lui fait poser un jeton à chaque marche d'un escalier de
cinq marches, chaque marche mesurant deux pas en avant puis deux pas de côté.
\item[(c)] Faites exécuter vos deux textes à la lettre par quelqu'un d'autre, qui a le droit --- le
devoir --- de choisir la lecture la plus bête chaque fois qu'il hésite. Notez chaque endroit
où il a pu choisir, et corrigez le texte jusqu'à ce qu'aucun choix ne subsiste.
\item[(d)] Une consigne peut être parfaitement claire et fausse quand même. Donnez un exemple de
chacun des quatre cas : claire et juste, claire et fausse, ambiguë mais juste par chance,
ambiguë et fausse. Dites lequel des quatre est le plus difficile à repérer.
\end{enumerate}

\end{problem}

\noindent Seymour Papert a bâti le langage Logo autour de cette idée en 1967 : une
tortue qui n'obéit qu'à quelques ordres, et l'enfant qui découvre que \og{} faire un carré \fg{} et
\og{} dire comment faire un carré \fg{} sont deux compétences différentes. Scratch, aujourd'hui, en
est l'héritier direct. Le point sérieux, c'est que l'ordinateur est exactement cet exécutant
sans bon sens : il ne devine rien, ne corrige rien, et n'a jamais l'intention que vous aviez.
La question (c) est la véritable expérience --- presque tout le monde écrit d'abord une
consigne que seul quelqu'un qui sait déjà pourrait suivre, et la découvrir est plus utile
que dix leçons.

\begin{refs}
\item Contexte : Algorithme --- Wikipédia (\url{https://fr.wikipedia.org/wiki/Algorithme})
\item Contexte : Logo (langage) --- Wikipédia (\url{https://fr.wikipedia.org/wiki/Logo_(langage)})
\item Contexte : Seymour Papert --- Wikipédia (\url{https://fr.wikipedia.org/wiki/Seymour_Papert})
\end{refs}

\bigskip

\begin{problem}[{Les cinq cartes qui devinent votre nombre --- Collège}]
\textbf{DCS-55.} \textbf{Problème.} On fabrique cinq cartes portant chacune une liste de nombres entre 1
et 31 : sur la première figurent tous les nombres dont l'écriture comme somme de nombres
choisis parmi 1, 2, 4, 8, 16 utilise le 1 ; sur la deuxième, tous ceux dont l'écriture
utilise le 2 ; et ainsi de suite. Vous pensez à un nombre, on vous montre les cartes une à
une, vous dites seulement sur lesquelles il figure --- et l'on annonce votre nombre.
\begin{enumerate}
\item[(a)] Démontrez que tout entier de 1 à 31 s'écrit d'exactement une façon comme somme de
nombres choisis parmi 1, 2, 4, 8 et 16. Il y a deux choses à établir : qu'une telle écriture
existe, et qu'il n'y en a pas deux.
\item[(b)] Déduisez-en pourquoi le tour fonctionne, et dites précisément quel calcul il faut faire
une fois connue la liste des cartes où le nombre figure.
\item[(c)] Écrivez la première carte en entier. Combien de nombres porte-t-elle ? Démontrez-le sans
les compter un à un.
\item[(d)] Combien de cartes faudrait-il pour aller jusqu'à 1000 ? Démontrez votre réponse. Et si
chaque carte pouvait recevoir trois réponses au lieu de deux, quels nombres écririez-vous à
la place de 1, 2, 4, 8, 16 ?
\end{enumerate}

\end{problem}

\noindent Leibniz a publié en 1703, dans les mémoires de l'Académie royale des
sciences de Paris, une \textit{Explication de l'arithmétique binaire} où il montre que deux
chiffres suffisent à écrire tous les nombres ; il y voyait une élégance presque théologique,
et n'imaginait évidemment pas les machines qui en vivraient trois siècles plus tard. Le tour
de cartes est aujourd'hui l'une des activités les plus répandues de l'informatique
\og{} débranchée \fg{}, qui enseigne les idées de la discipline sans le moindre ordinateur --- et il
est honnête, en ceci qu'il ne cache aucun truc : le tour marche parce que le théorème de la
question (a) est vrai. La variante à trois réponses de la question (d) est un très vieux
problème, celui des poids de Bachet de Méziriac (1612), où il s'agit de peser toutes les
masses entières avec le moins de poids possible.

\begin{refs}
\item Contexte : Système binaire --- Wikipédia (\url{https://fr.wikipedia.org/wiki/Syst%C3%A8me_binaire})
\item Contexte : Claude-Gaspard Bachet de Méziriac --- Wikipédia (\url{https://fr.wikipedia.org/wiki/Claude-Gaspard_Bachet_de_M%C3%A9ziriac})
\item Cours : CS Unplugged --- activités d'informatique sans ordinateur (en anglais) (\url{https://www.csunplugged.org/})
\end{refs}

\bigskip

\begin{problem}[{L'algorithme d'Euclide, déroulé à la main --- Collège}]
\textbf{DCS-56.} \textbf{Problème.} Pour trouver le plus grand diviseur commun de deux entiers strictement
positifs a et b, avec $ a > b $, on remplace le couple $ (a, b) $ par le couple
$ (b, r) $, où r est le reste de la division euclidienne de a par b ; on recommence
jusqu'à obtenir un reste nul.
\begin{enumerate}
\item[(a)] Déroulez-le à la main sur 1071 et 462, puis sur 2024 et 748, en écrivant chaque division
sur une ligne.
\item[(b)] Démontrez le fait sur lequel tout repose : les diviseurs communs de a et de b sont
exactement les diviseurs communs de b et de r. Commencez par le cas plus simple où
$ r = a - b $.
\item[(c)] Démontrez que le procédé s'arrête toujours, quel que soit le couple de départ.
\item[(d)] Rassemblez (b) et (c) en une seule démonstration complète : le dernier reste non nul est
bien le plus grand diviseur commun cherché. C'est ce qu'on appelle démontrer la correction
d'un algorithme.
\item[(e)] Cherchez, parmi les couples d'entiers inférieurs à 100, celui qui fait faire le plus
d'étapes à l'algorithme. Regardez la suite des nombres qui apparaissent au fil du calcul :
la reconnaissez-vous ?
\end{enumerate}

\end{problem}

\noindent C'est le plus vieil algorithme non trivial encore en service : il figure
aux propositions 1 et 2 du livre VII des \textit{Éléments} d'Euclide, vers 300 av. J.-C., sous
forme de soustractions répétées de longueurs, et il tourne aujourd'hui dans chaque
transaction chiffrée que fait votre téléphone. La question (e) touche à un moment
historique : en 1844, Gabriel Lamé a démontré que le pire cas de l'algorithme d'Euclide est
lié à la suite de Fibonacci --- l'un des tout premiers résultats jamais obtenus sur le
\textit{temps de calcul} d'une méthode, et non sur son résultat. DCS-30 propose une variante
qui se passe entièrement de divisions.

\begin{refs}
\item Contexte : Algorithme d'Euclide --- Wikipédia (\url{https://fr.wikipedia.org/wiki/Algorithme_d%27Euclide})
\item Contexte : Plus grand commun diviseur --- Wikipédia (\url{https://fr.wikipedia.org/wiki/Plus_grand_commun_diviseur})
\item Contexte : Gabriel Lamé --- Wikipédia (\url{https://fr.wikipedia.org/wiki/Gabriel_Lam%C3%A9})
\end{refs}

\bigskip

\begin{problem}[{Le bit qui rattrape une erreur --- Collège}]
\textbf{DCS-57.} \textbf{Problème.} Disposez 36 cartes en carré $ 6 \times 6 $, chacune montrant au hasard
sa face noire ou sa face blanche. Ajoutez une septième carte au bout de chaque ligne, de
sorte que chaque ligne de sept cartes contienne un nombre pair de faces noires ; puis, de la
même façon, une septième carte au bas de chaque colonne. La case du coin est alors réclamée
deux fois, une fois par sa ligne et une fois par sa colonne.
\begin{enumerate}
\item[(a)] Démontrez que ces deux exigences ne se contredisent jamais : la couleur que réclame la
ligne du coin est toujours celle que réclame sa colonne. (Comptez de deux façons le nombre
total de faces noires.)
\item[(b)] Pendant que vous avez le dos tourné, quelqu'un retourne exactement une carte. Démontrez
que vous pouvez toujours dire laquelle.
\item[(c)] Il en retourne maintenant exactement deux. Démontrez que vous vous apercevez toujours
que quelque chose a changé, puis donnez un exemple où vous ne pouvez pas dire lesquelles.
\item[(d)] Et s'il en retournait quatre, bien choisies ? Décidez, avec démonstration, si vous êtes
encore certain de vous apercevoir de quelque chose.
\end{enumerate}

\end{problem}

\noindent À la fin des années 1940, Richard Hamming disposait des Bell Labs le
week-end pour faire tourner ses calculs sur une machine à relais ; au moindre défaut de
lecture, la machine s'arrêtait et le week-end était perdu. \og{} Si elle sait qu'il y a une
erreur, pourquoi ne saurait-elle pas où ? \fg{} --- de cette exaspération est sorti le code de
Hamming (1950) et, avec lui, toute la théorie des codes correcteurs. Les cartes de ce
problème en sont la version la plus simple, et elles vous font traverser en quatre questions
tout le sujet : détecter une erreur, la localiser, découvrir la limite du procédé, et voir
exactement quelles configurations d'erreurs passent au travers. Le même principe protège
vos numéros de carte bancaire, les codes ISBN des livres et les données envoyées par les
sondes spatiales.

\begin{refs}
\item Contexte : Somme de contrôle --- Wikipédia (\url{https://fr.wikipedia.org/wiki/Somme_de_contr%C3%B4le})
\item Contexte : Code de Hamming --- Wikipédia (\url{https://fr.wikipedia.org/wiki/Code_de_Hamming})
\item Contexte : Richard Hamming --- Wikipédia (\url{https://fr.wikipedia.org/wiki/Richard_Hamming})
\end{refs}

\bigskip

\begin{problem}[{Deux exécutions de la même consigne --- Collège, short}]
\textbf{DCS-58.} \textbf{Problème.} Appliquez à la liste 4, 7, 1 la consigne \og{} prends le
premier nombre, ajoute 2 et multiplie par 3, écris le résultat, puis recommence avec le
nombre suivant \fg{} de deux façons toutes deux défendables et qui ne donnent pas la même sortie.
Réécrivez ensuite la consigne pour qu'aucune autre lecture ne soit possible, et démontrez que
votre version ne laisse plus le choix.
\end{problem}

\noindent Le 23 septembre 1999, la sonde Mars Climate Orbiter s'est perdue parce que
deux équipes avaient lu la même grandeur dans deux unités différentes, l'une en newtons,
l'autre en livres-force : une mission de plusieurs centaines de millions de dollars détruite
par une consigne que l'on pouvait lire de deux façons. Un ordinateur ne
demande jamais \og{} tu veux dire quoi ? \fg{} --- il choisit une lecture, la vôtre ou l'autre, et
continue. Écrire sans ambiguïté n'est donc pas une élégance de rédaction mais une exigence
technique, et la question difficile de ce problème est la dernière : démontrer qu'un texte
ne peut plus être lu de deux manières est beaucoup plus dur que de le croire clair.

\begin{refs}
\item Contexte : Ambiguïté --- Wikipédia (\url{https://fr.wikipedia.org/wiki/Ambigu%C3%AFt%C3%A9})
\item Contexte : Mars Climate Orbiter --- Wikipédia (\url{https://fr.wikipedia.org/wiki/Mars_Climate_Orbiter})
\item Contexte : Spécification (informatique) --- Wikipédia (\url{https://fr.wikipedia.org/wiki/Sp%C3%A9cification_(informatique)})
\end{refs}

\bigskip

\begin{problem}[{Sortir du labyrinthe --- Collège}]
\textbf{DCS-59.} \textbf{Problème.} Vous êtes dans un labyrinthe, sans plan et sans vue d'ensemble. La
\og{} règle de la main droite \fg{} dit : posez la main droite sur un mur et avancez sans jamais la
décoller.
\begin{enumerate}
\item[(a)] Dessinez un labyrinthe dont tous les murs forment un seul bloc rattaché au bord
extérieur, et suivez-y la règle. Décrivez ce qui se passe, et expliquez pourquoi vous
finissez par ressortir.
\item[(b)] Dessinez maintenant un labyrinthe où la règle échoue --- où l'on tourne indéfiniment sans
jamais atteindre le but. Dites ce que votre labyrinthe a que celui de (a) n'avait pas.
\item[(c)] Voici une méthode qui, elle, ne se laisse pas piéger. En entrant dans un couloir, tracez
un trait à son entrée ; en en sortant, tracez un trait à sa sortie. N'empruntez jamais un
couloir portant déjà deux traits ; à un carrefour où vous n'êtes jamais venu, prenez
n'importe quel couloir vierge ; à un carrefour déjà visité, revenez sur vos pas si le couloir
d'où vous venez ne porte qu'un trait, et sinon prenez un couloir en portant le moins.
Démontrez que chaque couloir est parcouru au plus deux fois, une fois dans chaque sens.
\item[(d)] Expliquez, avec le plus de précision que vous pourrez, pourquoi cette méthode finit par
visiter tout carrefour accessible depuis votre point de départ. Puis désignez l'endroit de
votre explication qui s'appuie sur quelque chose que vous n'avez pas démontré.
\end{enumerate}

\end{problem}

\noindent La méthode de la question (c) est due à Trémaux, ingénieur français du
XIXe siècle, dont Édouard Lucas rapporte le procédé dans le premier tome de ses
\textit{Récréations mathématiques} (1882), où il occupe le chapitre sur les labyrinthes. On la
reconnaît aujourd'hui comme le parcours en profondeur, l'un des deux ou trois algorithmes de
graphes dont tout le reste dépend --- et c'est un fait remarquable qu'un problème posé par un
promeneur perdu et un problème posé par un compilateur reçoivent la même réponse. La question
(d) est délibérément honnête sur son niveau : elle vous demande de repérer le trou de votre
propre démonstration, ce qui est un savoir-faire à part entière et souvent le premier pas
vers la bonne démonstration.

\begin{refs}
\item Contexte : Parcours en profondeur --- Wikipédia (\url{https://fr.wikipedia.org/wiki/Parcours_en_profondeur})
\item Contexte : Algorithmes de résolution de labyrinthe --- Wikipédia (en anglais) (\url{https://en.wikipedia.org/wiki/Maze-solving_algorithm})
\item Lecture : É. Lucas, \textit{Récréations mathématiques}, tome I (1882), le chapitre sur les labyrinthes
\end{refs}

\bigskip

\begin{problem}[{Trois cartes, combien de comparaisons ? --- Collège, short}]
\textbf{DCS-60.} \textbf{Problème.} On vous donne trois cartes de valeurs différentes, face
cachée, et vous n'avez le droit que de comparer deux cartes à la fois. Démontrez que trois
comparaisons suffisent toujours à les ranger dans l'ordre, et qu'aucune méthode n'y parvient
à coup sûr en deux.
\end{problem}

\noindent La première moitié se règle en exhibant une méthode ; la seconde parle de
toutes les méthodes possibles, et se règle en comptant ce que deux réponses peuvent au plus
distinguer. C'est la plus petite borne inférieure honnête qui existe, et c'est le même
argument qui, poussé jusqu'au bout, démontre qu'aucun tri par comparaisons ne peut être
sensiblement plus rapide que ceux qu'on connaît. DCS-29 en donne la version \og{} trouver le plus
grand \fg{} et DCS-51 la version \og{} ranger sept cartes \fg{}.

\begin{refs}
\item Contexte : Tri par comparaison --- Wikipédia (\url{https://fr.wikipedia.org/wiki/Tri_par_comparaison})
\item Contexte : Algorithme de tri --- Wikipédia (\url{https://fr.wikipedia.org/wiki/Algorithme_de_tri})
\end{refs}

\bigskip


\section*{Casse-tête}

\begin{problem}[{Les trois boîtes mal étiquetées --- Collège}]
\textbf{PUZ-101.} \textbf{Problème.} Trois boîtes fermées portent les étiquettes \og{} pommes \fg{}, \og{} oranges \fg{} et
\og{} mélange \fg{}. L'une ne contient que des pommes, une autre que des oranges, la troisième un
mélange des deux --- mais les étiquettes ont été posées n'importe comment et, par malchance,
aucune des trois n'est correcte. Vous pouvez sortir un fruit de la boîte de votre choix, sans
regarder à l'intérieur.
\begin{enumerate}
\item[(a)] Déterminez le contenu des trois boîtes en ne sortant qu'un seul fruit, et démontrez que
votre procédé aboutit quel que soit le fruit que vous obtenez.
\item[(b)] Démontrez qu'après ce tirage la réponse est réellement forcée : il n'existe pas deux
répartitions différentes des contenus qui soient compatibles à la fois avec \og{} aucune
étiquette n'est correcte \fg{} et avec le fruit que vous avez vu.
\item[(c)] On affaiblit l'hypothèse : on vous garantit seulement qu'\textit{au moins} une étiquette est
fausse. Une boîte mixte peut vous donner l'un ou l'autre fruit, et c'est elle qui décide,
pas vous. Déterminez s'il existe encore une procédure qui garantisse la réponse en un nombre
fini de tirages, et démontrez ce que vous affirmez.
\end{enumerate}

\end{problem}

\noindent Ce casse-tête circule depuis des décennies dans les colonnes de jeux et
dans les entretiens d'embauche, et il doit sa longévité à un détail que la plupart des gens
lisent sans le voir : la phrase \og{} aucune étiquette n'est correcte \fg{} n'est pas une décoration
de l'histoire, c'est la donnée qui rend un tirage unique suffisant. Traduisez-la en langage de
permutations --- quelles répartitions des trois contenus sur les trois étiquettes restent
debout ? --- et le casse-tête devient un exercice de comptage plutôt qu'un exercice de chance.
La question (c) est là pour montrer ce que coûte l'affaiblissement d'une hypothèse : un énoncé
soluble en une observation peut devenir insoluble en un million, et la démonstration de cette
impossibilité se mène en imaginant un adversaire qui choisit les fruits au fur et à mesure,
uniquement pour vous garder dans le doute. C'est exactement la technique qu'emploient les
informaticiens pour minorer le nombre de questions nécessaires à un algorithme.

\begin{refs}
\item Contexte : Martin Gardner --- Wikipédia (\url{https://fr.wikipedia.org/wiki/Martin_Gardner})
\item Contexte : Raisonnement par l'absurde --- Wikipédia (\url{https://fr.wikipedia.org/wiki/Raisonnement_par_l%27absurde})
\item Voir aussi : PUZ-01 (\url{#PUZ-01}) ci-dessous, pour la même façon de compter ce qu'une seule observation peut trancher
\end{refs}

\bigskip

\begin{problem}[{L'escargot au fond du puits --- Collège, short}]
\textbf{PUZ-102.} \textbf{Problème.} Un escargot part du fond d'un puits de $ 10 $ mètres :
chaque jour il grimpe de $ 3 $ mètres, chaque nuit il en redescend $ 2 $. Déterminez le
jour où il atteint le bord, et démontrez que votre compte est exact --- c'est-à-dire que ni la
veille ni le lendemain ne conviennent.
\end{problem}

\noindent Les puits, les fosses et les murs rongés par les rats peuplent les recueils
d'arithmétique du Moyen Âge : le \textit{Liber abaci} de Leonardo Fibonacci, en 1202, contient
la version du lion tombé au fond d'une fosse, qui remonte d'une certaine fraction le jour et
reglisse d'une autre la nuit. Ce qui a fait traverser huit siècles à ce problème, ce n'est pas
la soustraction $ 3 - 2 $ mais le fait que le dernier jour ne ressemble pas aux autres : le
régime régulier cesse de s'appliquer au moment précis où l'on en a besoin. C'est la même
erreur d'une unité que l'on commet en comptant les piquets d'une clôture, en numérotant les
pages d'un livre ou en écrivant une boucle informatique, et savoir la traquer vaut mieux que
savoir diviser.

\begin{refs}
\item Contexte : Liber abaci --- Wikipédia (\url{https://fr.wikipedia.org/wiki/Liber_abaci})
\item Contexte : Leonardo Fibonacci --- Wikipédia (\url{https://fr.wikipedia.org/wiki/Leonardo_Fibonacci})
\end{refs}

\bigskip

\begin{problem}[{Le carré magique de Luo Shu --- Collège}]
\textbf{PUZ-103.} \textbf{Problème.} Un carré magique d'ordre $ 3 $ est un tableau de $ 3 \times 3 $ cases
rempli avec les entiers de $ 1 $ à $ 9 $, chacun employé une fois et une seule, de sorte
que les trois lignes, les trois colonnes et les deux diagonales aient toutes la même
somme.
\begin{enumerate}
\item[(a)] Déterminez, avec démonstration, la valeur de cette somme commune, puis celle que porte
nécessairement la case centrale.
\item[(b)] Complétez le carré ci-dessous et démontrez qu'il n'admet qu'une seule complétion :
\[ \begin{array}{ccc} 8 & 1 & ? \\ ? & ? & ? \\ ? & ? & 2 \end{array} \]
\item[(c)] Démontrez qu'il existe exactement huit carrés magiques d'ordre $ 3 $, et qu'ils se
déduisent tous d'un même carré par rotations et symétries.
\item[(d)] On remplace $ 1, \dots, 9 $ par neuf entiers distincts quelconques, la définition étant
inchangée. Démontrez que la case centrale vaut toujours le tiers de la somme magique, et
construisez une infinité de tels carrés.
\end{enumerate}

\end{problem}

\noindent Le carré de Luo Shu est le plus ancien carré magique connu : la tradition
chinoise le fait apparaître sur la carapace d'une tortue sortie de la rivière Luo, et il
figure dans les textes cosmologiques bien avant d'être vu comme un objet mathématique. On le
retrouve ensuite partout --- dans le monde arabe, dans l'Inde médiévale, puis en Europe, où
Albrecht Dürer place un carré magique d'ordre $ 4 $ dans sa gravure \textit{Melencolia I} de
$ 1514 $, daté par les deux cases centrales de sa dernière ligne. La beauté de l'ordre
$ 3 $ est qu'il ne laisse aucune place au hasard : les huit alignements se recoupent
tellement que presque tout est forcé dès les premières lignes de raisonnement, et la question
(c) demande de transformer cette impression en démonstration. C'est un excellent premier
contact avec l'idée de classification --- non pas exhiber un objet, mais prouver qu'on les a
tous.

\begin{refs}
\item Contexte : Carré de Luo Shu --- Wikipédia (\url{https://fr.wikipedia.org/wiki/Carr%C3%A9_de_Luo_Shu})
\item Contexte : Carré magique (mathématiques) --- Wikipédia (\url{https://fr.wikipedia.org/wiki/Carr%C3%A9_magique_(math%C3%A9matiques)})
\item Voir aussi : PUZ-16 (\url{#PUZ-16}) ci-dessous, la variante avec des carrés parfaits, toujours ouverte
\end{refs}

\bigskip

\begin{problem}[{Les quatre poids de Bachet --- Collège}]
\textbf{PUZ-104.} \textbf{Problème.} Une balance à deux plateaux ne fait que comparer : elle dit quel côté
descend, ou que les deux s'équilibrent. Vous cherchez un jeu de quatre poids, de masses
entières en kilogrammes, qui permette de peser n'importe quelle masse entière de $ 1 $ à
$ 40 $ kg : la marchandise est posée sur un plateau, et vous avez le droit de répartir vos
poids sur les \textit{deux} plateaux à votre guise.
\begin{enumerate}
\item[(a)] Trouvez un tel jeu de quatre poids, et démontrez qu'il pèse effectivement toutes les
masses entières de $ 1 $ à $ 40 $.
\item[(b)] Démontrez qu'aucun jeu de quatre poids ne peut faire mieux : quelles que soient les quatre
masses choisies, il est impossible de peser toutes les masses entières de $ 1 $ à
$ 41 $.
\item[(c)] Reprenez tout avec la règle du marchand : les poids ne peuvent aller que sur le plateau
opposé à la marchandise. Déterminez le meilleur jeu de quatre poids et démontrez la borne
correspondante.
\item[(d)] Généralisez les deux régimes à $ n $ poids : donnez dans chaque cas la plus grande masse
$ M(n) $ telle que toutes les masses entières de $ 1 $ à $ M(n) $ soient pesables, et
démontrez vos deux formules.
\end{enumerate}

\end{problem}

\noindent Claude-Gaspard Bachet de Méziriac pose ce problème dans ses \textit{Problèmes
plaisants et délectables qui se font par les nombres}, en $ 1612 $, le premier recueil
imprimé de mathématiques récréatives digne de ce nom --- c'est aussi l'homme qui traduisit
Diophante en latin, dans la marge duquel Fermat écrira sa fameuse note. Des versions du
problème des poids circulaient déjà chez les marchands médiévaux, mais Bachet en fait une
question de principe, et la réponse est l'une des plus jolies rencontres entre le commerce et
la numération : autoriser un poids à changer de plateau, c'est lui donner trois emplois
possibles au lieu de deux, et compter le nombre de résultats accessibles revient à écrire les
nombres dans une base. La question (b) est la vraie question du site : trouver un jeu qui
marche est affaire de tâtonnement, démontrer qu'aucun jeu ne fait mieux demande de raisonner
sur tous les jeux à la fois.

\begin{refs}
\item Contexte : Claude-Gaspard Bachet de Méziriac --- Wikipédia (\url{https://fr.wikipedia.org/wiki/Claude-Gaspard_Bachet_de_M%C3%A9ziriac})
\item Contexte : Système ternaire --- Wikipédia (\url{https://fr.wikipedia.org/wiki/Syst%C3%A8me_ternaire})
\item Contexte : Problème de pesée --- Wikipédia (\url{https://fr.wikipedia.org/wiki/Probl%C3%A8me_de_pes%C3%A9e})
\end{refs}

\bigskip

\begin{problem}[{Casser la tablette de chocolat --- Collège, short}]
\textbf{PUZ-105.} \textbf{Problème.} Une tablette de chocolat est un rectangle de
$ m \times n $ carrés séparés par des rainures, et un coup consiste à prendre un morceau et
à le casser en deux en suivant une rainure, d'un bord à l'autre. Déterminez le nombre de
cassures nécessaires pour réduire la tablette en $ mn $ carrés isolés, et démontrez que ce
nombre ne dépend jamais de l'ordre dans lequel on casse.
\end{problem}

\noindent Voilà le premier exemple d'invariant qu'on donne aux élèves dans les cercles
mathématiques russes, et il est difficile d'en trouver un meilleur : la stratégie ne sert à
rien, l'ingéniosité ne sert à rien, parce qu'une certaine quantité varie de la même façon à
chaque coup quel que soit le coup. Le jour où l'on comprend qu'il faut suivre cette
quantité-là plutôt que la forme des morceaux, on a appris un réflexe qui resservira devant
des problèmes autrement plus sérieux --- un invariant transforme une infinité de scénarios en
une seule ligne de calcul. Le raisonnement se rédige aussi bien par récurrence que par
comptage direct, et comparer les deux rédactions est instructif.

\begin{refs}
\item Contexte : Invariant --- Wikipédia (\url{https://fr.wikipedia.org/wiki/Invariant})
\item Contexte : Raisonnement par récurrence --- Wikipédia (\url{https://fr.wikipedia.org/wiki/Raisonnement_par_r%C3%A9currence})
\item Lecture : Dmitri Fomin, Sergey Genkin \& Ilia Itenberg, \textit{Mathematical Circles (Russian Experience)}, chapitre sur les invariants
\end{refs}

\bigskip

\begin{problem}[{L'héritage des dix-sept chameaux --- Collège}]
\textbf{PUZ-106.} \textbf{Problème.} Un homme laisse dix-sept chameaux à ses trois fils : la moitié à l'aîné,
le tiers au deuxième, le neuvième au cadet. Comme dix-sept n'est divisible ni par deux, ni par
trois, ni par neuf, les fils se querellent. Un voisin prête alors un chameau : le partage se
fait sur dix-huit bêtes, chacun emporte sa part entière, et le voisin repart avec son
chameau.
\begin{enumerate}
\item[(a)] Effectuez le partage sur dix-huit bêtes, et vérifiez que chaque fils reçoit un nombre
entier de chameaux et qu'il en reste exactement un à rendre.
\item[(b)] Chaque fils a-t-il reçu exactement la part que le testament lui promettait ? Comparez et
expliquez précisément d'où vient l'écart.
\item[(c)] Formalisez ce qui fait marcher le tour : on cherche trois fractions $ \frac{1}{a} $,
$ \frac{1}{b} $, $ \frac{1}{c} $ avec $ 2 \leq a \leq b \leq c $ et un troupeau de
$ N $ bêtes tels que l'emprunt d'une seule bête donne à chacun un nombre entier de bêtes et
laisse exactement une bête à rendre. Écrivez cette condition sous forme d'équation.
\item[(d)] Trouvez tous les testaments possibles, et démontrez que votre liste est complète. Traitez
ensuite le cas de deux héritiers, puis de quatre.
\end{enumerate}

\end{problem}

\noindent L'histoire se raconte de l'Iran au Maghreb comme un conte de sagesse, avec
un cadi ou un vieillard à la place du voisin ; la plus ancienne trace imprimée qu'on lui
connaisse en Europe se trouve chez Niccolò Fontana Tartaglia, au milieu du seizième siècle,
dans son \textit{General trattato di numeri et misure}. Le tour n'est pas un tour de magie mais
une identité entre fractions unitaires, du genre de celles que les scribes de l'Égypte
ancienne manipulaient déjà : dire laquelle, et pourquoi elle permet au chameau emprunté de
revenir, c'est exactement le travail des questions (b) et (c). La question (d) est un vrai
petit problème de théorie des nombres --- chercher toutes les décompositions d'un nombre en
somme de fractions unitaires est une activité qui remonte au papyrus Rhind et qui n'a jamais
cessé d'occuper les mathématiciens.

\begin{refs}
\item Contexte : Fraction égyptienne --- Wikipédia (\url{https://fr.wikipedia.org/wiki/Fraction_%C3%A9gyptienne})
\item Contexte : Niccolò Fontana Tartaglia --- Wikipédia (\url{https://fr.wikipedia.org/wiki/Niccol%C3%B2_Fontana_Tartaglia})
\end{refs}

\bigskip

\begin{problem}[{Le verre d'eau et le verre de vin --- Collège, short}]
\textbf{PUZ-107.} \textbf{Problème.} Un verre contient un décilitre d'eau, un autre un
décilitre de vin ; on prend une cuillerée d'eau qu'on verse dans le vin, on remue autant qu'on
veut, puis on reprend une cuillerée du mélange qu'on reverse dans l'eau. Y a-t-il alors plus
de vin dans l'eau que d'eau dans le vin, et pouvez-vous le démontrer sans rien supposer sur
l'homogénéité du mélange ?
\end{problem}

\noindent C'est un classique des recueils de récréations anglais de la fin du dix-neuvième
siècle, et sa réputation vient de ce que les calculs de proportions, qui semblent la voie
naturelle, sont à la fois pénibles et inutiles : la bonne démonstration est une comptabilité,
pas un calcul de concentrations, et elle tient en deux lignes qui ne supposent rien du tout
sur la façon dont on a remué. On peut même laisser la cuillerée entièrement non mélangée, ou
la mélanger parfaitement, ou n'importe quoi entre les deux, sans changer la réponse --- c'est
ce genre de robustesse qui signale qu'une quantité se conserve quelque part. Le même
raisonnement, transposé, sert en chimie et en comptabilité en partie double.

\begin{refs}
\item Contexte : Problème du mélange eau-vin --- Wikipédia (en anglais) (\url{https://en.wikipedia.org/wiki/Wine/water_mixing_problem})
\item Contexte : Conservation de la masse --- Wikipédia (\url{https://fr.wikipedia.org/wiki/Conservation_de_la_masse})
\end{refs}

\bigskip

\begin{problem}[{L'éléphant de Cao Chong --- Collège}]
\textbf{PUZ-108.} \textbf{Problème.} On veut connaître la masse d'un éléphant. On dispose d'une rivière, d'une
barque assez grande pour le porter, d'un tas de pierres, d'un morceau de craie et d'une
balance ordinaire qui ne pèse que de petites charges --- une pierre à la fois.
\begin{enumerate}
\item[(a)] Décrivez une méthode qui donne la masse de l'éléphant, et justifiez chaque étape : énoncez
en particulier, avec précision, le principe physique qui autorise la comparaison que vous
faites.
\item[(b)] Démontrez ce principe dans le cas dont vous vous servez : deux chargements qui font
flotter la même barque exactement à la même ligne d'eau ont la même masse. Vous ne supposerez
que ceci --- la barque flotte à l'équilibre, et l'eau a partout la même masse volumique.
\item[(c)] Version chiffrée. À la ligne de flottaison, la barque a une section horizontale constante
de $ 6 $ m$^2$, et chargée de l'éléphant elle s'enfonce de $ 40 $ cm de plus qu'à vide.
Calculez la masse de l'éléphant, dites où exactement vous avez utilisé l'hypothèse \og{} section
constante \fg{}, et expliquez ce que devient le calcul si la coque s'évase vers le haut.
\item[(d)] On opère maintenant dans un bassin étroit, dont le niveau monte sensiblement quand la
barque s'enfonce. La méthode de la question (a) reste-t-elle exacte ? Démontrez votre
réponse.
\end{enumerate}

\end{problem}

\noindent Cao Chong, fils de Cao Cao, est mort à douze ans en l'an $ 208 $, au début
de la période des Trois Royaumes ; les \textit{Chroniques des Trois Royaumes} de Chen Shou
racontent que, tout enfant, il résolut le problème de la pesée d'un éléphant offert à son père
en le faisant monter dans une barque, en marquant la ligne d'eau, puis en remplaçant l'animal
par des pierres jusqu'à retrouver la même marque. L'histoire est aujourd'hui encore l'une des
plus connues des écoliers chinois, et rien n'indique que la Chine de cette époque connaissait
Archimède, mort quatre siècles plus tôt à Syracuse : la même idée physique a été trouvée deux
fois. Ce qu'il y a de mathématique ici, c'est le procédé de substitution --- on ne mesure pas
l'objet, on lui substitue quelque chose de mesurable qui lui est égal pour la seule propriété
qui compte ---, et c'est le geste de base de toute la métrologie, des étalons du kilogramme aux
balances de laboratoire.

\begin{refs}
\item Contexte : Cao Chong --- Wikipédia (\url{https://fr.wikipedia.org/wiki/Cao_Chong})
\item Contexte : Poussée d'Archimède --- Wikipédia (\url{https://fr.wikipedia.org/wiki/Pouss%C3%A9e_d%27Archim%C3%A8de})
\item Voir aussi : PUZ-59 (\url{#PUZ-59}) ci-dessous, où la même poussée décide du niveau d'un bassin
\end{refs}

\bigskip

\begin{problem}[{La corde autour de la Terre --- Collège, short}]
\textbf{PUZ-109.} \textbf{Problème.} Une corde tendue fait exactement le tour de l'équateur
d'une planète parfaitement sphérique, en épousant le sol. On l'allonge d'un mètre et on la
dispose en cercle à hauteur constante au-dessus du sol : calculez cette hauteur, puis
démontrez qu'elle est la même pour la Terre, pour un ballon de football et pour le Soleil.
\end{problem}

\noindent On fait remonter ce casse-tête à William Whiston, successeur de Newton à
Cambridge, qui le glisse dans un ouvrage du tout début du dix-huitième siècle ; depuis, il
n'a jamais cessé d'être posé, parce qu'il ne se démode pas --- l'écart entre ce que la formule
donne et ce que l'intuition annonce reste violent même quand on connaît la réponse. Le calcul
est à la portée d'un élève de cinquième, et c'est bien là ce qui rend le problème
philosophique : vous devez arbitrer entre une formule que vous savez par cœur et une
conviction qui hurle le contraire, et l'un des deux se trompe. Une variante bien plus
difficile consiste à ne pas soulever la corde partout, mais à la tirer en un seul point pour
la décoller du sol : la hauteur obtenue est alors surprenante en sens inverse, et sa
détermination sort de l'arithmétique élémentaire.

\begin{refs}
\item Contexte : Circonférence --- Wikipédia (\url{https://fr.wikipedia.org/wiki/Circonf%C3%A9rence})
\item Contexte : William Whiston --- Wikipédia (\url{https://fr.wikipedia.org/wiki/William_Whiston})
\item Contexte : Pi --- Wikipédia (\url{https://fr.wikipedia.org/wiki/Pi})
\end{refs}

\bigskip

\begin{problem}[{Le champ de Didon --- Collège}]
\textbf{PUZ-110.} \textbf{Problème.} Vous disposez d'une corde de longueur $ L $ et vous voulez enclore le
champ d'aire la plus grande possible.
\begin{enumerate}
\item[(a)] Parmi tous les rectangles de périmètre $ L $, déterminez celui dont l'aire est maximale,
et démontrez qu'aucun autre rectangle ne fait mieux.
\item[(b)] Le champ borde maintenant une rivière rectiligne, et l'on n'a pas besoin de clôturer ce
côté-là. Déterminez le meilleur rectangle, démontrez-le, et comparez son aire à celle
obtenue en (a).
\item[(c)] Démontrez que le disque de périmètre $ L $ est plus grand que tous les rectangles de
périmètre $ L $ --- la valeur $ \pi \approx 3{,}14 $ suffit. Traitez de même le cas de la
rivière, en comparant le demi-disque au meilleur rectangle de (b).
\item[(d)] Le théorème isopérimétrique affirme que le disque bat toute autre figure de même
périmètre. Énoncez-le proprement, puis expliquez ce qui manquait à la démonstration de Jakob
Steiner en $ 1838 $ et ce qu'il a fallu lui ajouter.
\end{enumerate}

\end{problem}

\noindent Virgile raconte dans l'\textit{Énéide} que Didon, fuyant Tyr, obtint sur la
côte d'Afrique autant de terre qu'une peau de bœuf pouvait en couvrir ; elle fit découper la
peau en lanières très fines, les mit bout à bout, et enferma dans cette corde la colline qui
devint Carthage. Les mathématiciens ont gardé le nom : le \og{} problème de Didon \fg{} est le premier
problème d'optimisation de forme de l'histoire, et Zénodore, vers le deuxième siècle avant
notre ère, démontrait déjà que le cercle l'emporte sur les polygones. La question (d) est une
leçon d'honnêteté mathématique : Steiner produit en $ 1838 $ de superbes arguments de
symétrisation montrant que toute figure autre que le disque peut être améliorée, mais cela ne
démontre pas que l'optimum existe --- Weierstrass devra fournir plus tard l'argument
d'existence. Un raisonnement qui prouve \og{} si un maximum existe, c'est celui-là \fg{} n'est pas une
démonstration tant que la parenthèse n'est pas refermée.

\begin{refs}
\item Contexte : Isopérimétrie --- Wikipédia (\url{https://fr.wikipedia.org/wiki/Isop%C3%A9rim%C3%A9trie})
\item Contexte : Didon --- Wikipédia (\url{https://fr.wikipedia.org/wiki/Didon})
\item Contexte : Jakob Steiner --- Wikipédia (\url{https://fr.wikipedia.org/wiki/Jakob_Steiner})
\end{refs}

\bigskip

\begin{problem}[{Le carré qui gagne un carreau --- Collège}]
\textbf{PUZ-111.} \textbf{Problème.} Sur du papier quadrillé, prenez le carré $ 8 \times 8 $ dont les coins
opposés sont $ (0,0) $ et $ (8,8) $, et découpez-le ainsi : une coupe verticale en
$ x = 5 $ sur toute la hauteur ; dans la bande de gauche, la coupe droite qui joint
$ (0,3) $ à $ (5,5) $ ; dans la bande de droite, la coupe droite qui joint $ (5,0) $ à
$ (8,8) $. Vous obtenez deux trapèzes et deux triangles. Rangez maintenant ces quatre pièces
dans un rectangle $ 13 \times 5 $ de coins $ (0,0) $ et $ (13,5) $ : un triangle en
$ (0,0), (8,0), (8,3) $, un trapèze en $ (0,0), (0,5), (5,5), (5,2) $, le second triangle
en $ (5,2), (5,5), (13,5) $ et le second trapèze en $ (8,0), (13,0), (13,5), (8,3) $.
\begin{enumerate}
\item[(a)] Vérifiez que les quatre pièces découpent exactement le carré --- ni trou, ni recouvrement ---
et calculez leur aire totale. Vérifiez ensuite qu'aucune des quatre ne dépasse du rectangle et
qu'elles ne se recouvrent pas.
\item[(b)] Le rectangle a pour aire $ 65 $. D'où vient le carreau supplémentaire ? Démontrez que
les points $ (0,0) $, $ (8,3) $ et $ (13,5) $ ne sont pas alignés, puis calculez
exactement l'aire de la partie du rectangle que les quatre pièces laissent découverte.
\item[(c)] Recommencez avec le carré $ 13 \times 13 $ et le rectangle $ 8 \times 21 $ : coupe
verticale en $ x = 8 $, coupe de $ (0,5) $ à $ (8,8) $ à gauche, coupe de $ (8,0) $ à
$ (13,13) $ à droite. Cette fois l'écart change de signe : dites lequel, expliquez ce que
cela signifie concrètement pour les pièces, et démontrez que dans les deux cas l'écart vaut
exactement une unité d'aire.
\item[(d)] Démontrez l'identité de Cassini $ F_{n-1}F_{n+1} - F_n^2 = (-1)^n $ pour la suite de
Fibonacci $ 1, 1, 2, 3, 5, 8, 13, 21, \dots $, et déduisez-en que le tour se refait
indéfiniment. Expliquez enfin, en comparant les pentes $ F_{n-2}/F_{n-1} $ et
$ F_{n-1}/F_n $, pourquoi l'illusion devient de plus en plus convaincante quand $ n $
grandit.
\end{enumerate}

\end{problem}

\noindent Le paradoxe est publié par Oskar Schlömilch en $ 1868 $ ; Sam Loyd s'en
attribuera plus tard la paternité, comme de beaucoup d'autres, et Lewis Carroll, qui était
professeur de mathématiques à Oxford, s'y est intéressé de près. Sa version en triangle, due à
Paul Curry vers $ 1953 $, a fait le tour du monde sous le nom de triangle de Curry. Ce qu'il
y a d'instructif ici, c'est que la démonstration ne consiste pas à trouver l'erreur du
découpage --- le découpage du carré est parfaitement exact --- mais à découvrir que l'assemblage
en rectangle est \textit{presque} exact, et que le \og{} presque \fg{} a une aire mesurable : une
lamelle si allongée que l'œil la prend pour un trait de crayon. Derrière l'illusion se cache
une identité algébrique sur la suite de Fibonacci, et c'est pourquoi le même truc marche avec
$ 3, 5, 8 $, avec $ 5, 8, 13 $, avec $ 8, 13, 21 $, toujours à une unité près.

\begin{refs}
\item Contexte : Paradoxe du carré manquant --- Wikipédia (\url{https://fr.wikipedia.org/wiki/Paradoxe_du_carr%C3%A9_manquant})
\item Contexte : Identités de Cassini, de Catalan et de Vajda --- Wikipédia (\url{https://fr.wikipedia.org/wiki/Identit%C3%A9s_de_Cassini,_de_Catalan_et_de_Vajda})
\item Contexte : Suite de Fibonacci --- Wikipédia (\url{https://fr.wikipedia.org/wiki/Suite_de_Fibonacci})
\end{refs}

\bigskip

\begin{problem}[{La mouche entre les deux trains --- Collège}]
\textbf{PUZ-112.} \textbf{Problème.} Deux trains roulent l'un vers l'autre sur la même voie, chacun à
$ 50 $ km/h, séparés au départ par $ 100 $ km. Une mouche quitte au même instant le
pare-brise du premier à $ 75 $ km/h, atteint le second, fait aussitôt demi-tour, revient au
premier, repart, et ainsi de suite jusqu'à être écrasée entre les deux.
\begin{enumerate}
\item[(a)] Calculez la distance totale parcourue par la mouche, en justifiant chaque étape --- à
commencer par l'instant de la collision.
\item[(b)] Calculez la longueur du premier trajet de la mouche, puis celle du deuxième. Démontrez que
ces longueurs forment une suite géométrique dont vous déterminerez la raison, et expliquez
comment une somme d'une infinité de longueurs toutes strictement positives peut être
finie.
\item[(c)] Démontrez que la mouche fait une infinité de demi-tours avant la collision, alors que la
durée totale du vol est finie.
\item[(d)] Renversez le temps. Montrez que, connaissant le lieu et l'instant de la collision ainsi que
la vitesse de la mouche, on ne peut pas reconstituer son point de départ : plusieurs histoires
différentes aboutissent exactement au même écrasement.
\end{enumerate}

\end{problem}

\noindent L'anecdote est la plus célèbre des mathématiques du vingtième siècle : on
pose le problème à John von Neumann, qui donne la réponse presque immédiatement ; son
interlocuteur, un peu déçu, remarque que la plupart des gens essaient de sommer la série au
lieu de voir l'astuce, et von Neumann répond qu'il a justement sommé la série. Les deux voies
sont légitimes et le problème vaut par leur confrontation : l'une regarde le temps, l'autre
regarde la distance, et la question (b) vous demande de vérifier qu'elles donnent bien le même
nombre. La question (c) est un authentique paradoxe de Zénon apprivoisé --- une infinité
d'événements logés dans un intervalle fini ---, et la question (d) montre qu'une équation du
mouvement peut se laisser remonter dans un sens et pas dans l'autre, ce qui est une remarque
plus profonde qu'elle n'en a l'air.

\begin{refs}
\item Contexte : John von Neumann --- Wikipédia (\url{https://fr.wikipedia.org/wiki/John_von_Neumann})
\item Contexte : Série géométrique --- Wikipédia (\url{https://fr.wikipedia.org/wiki/S%C3%A9rie_g%C3%A9om%C3%A9trique})
\item Contexte : Paradoxes de Zénon --- Wikipédia (\url{https://fr.wikipedia.org/wiki/Paradoxes_de_Z%C3%A9non})
\end{refs}

\bigskip

\begin{problem}[{Le tournoi et le deuxième meilleur --- Collège}]
\textbf{PUZ-113.} \textbf{Problème.} Huit joueurs disputent un tournoi à élimination directe : quatre matchs au
premier tour, deux au deuxième, une finale. On suppose qu'il existe un classement strict des
huit joueurs, que le meilleur des deux gagne toujours, et qu'il n'y a pas de match nul.
\begin{enumerate}
\item[(a)] Démontrez que le vainqueur du tournoi est bien le meilleur des huit. Montrez en revanche
que le finaliste malheureux n'est pas nécessairement le deuxième, en exhibant un classement et
un tableau où le vrai deuxième est éliminé dès le premier tour.
\item[(b)] Le tournoi a coûté sept matchs. Organisez des matchs supplémentaires qui désignent à coup
sûr le deuxième meilleur joueur, pour un total de neuf matchs, et démontrez que votre méthode
ne se trompe jamais.
\item[(c)] Démontrez que huit matchs ne suffisent pas : aucune organisation, même adaptative --- chaque
match étant choisi en fonction des résultats précédents ---, ne peut garantir de désigner à la
fois le premier et le deuxième en huit matchs. C'est la moitié difficile du problème.
\item[(d)] Généralisez à $ n $ joueurs : donnez une méthode en $ n - 2 + \lceil \log_2 n \rceil $
matchs, et démontrez qu'on ne peut pas faire mieux.
\end{enumerate}

\end{problem}

\noindent En $ 1883 $, Charles Dodgson --- que le monde connaît sous le nom de Lewis
Carroll --- publie une brochure sur les tournois de tennis pour dénoncer une injustice bien
réelle : la méthode d'attribution du deuxième prix est fausse, et le joueur qui la reçoit n'est
en général pas le deuxième meilleur. Le problème a été repris sous forme mathématique par Hugo
Steinhaus, et le nombre exact de matchs nécessaires a été établi par S. S. Kislitsyn en
$ 1964 $ ; Donald Knuth en donne le traitement de référence dans le troisième volume de
\textit{The Art of Computer Programming}, où le tournoi devient un algorithme de sélection et le
match une comparaison. La question (c) est le cœur du sujet : pour minorer le nombre de matchs,
on ne raisonne pas sur les tournois possibles mais on imagine un adversaire qui, à chaque
match, choisit le vainqueur de façon à retarder le plus longtemps possible le moment où vous
saurez --- technique dite de l'adversaire, aujourd'hui centrale en complexité algorithmique.

\begin{refs}
\item Contexte : Tournoi à élimination directe --- Wikipédia (\url{https://fr.wikipedia.org/wiki/Tournoi_%C3%A0_%C3%A9limination_directe})
\item Contexte : Lewis Carroll --- Wikipédia (\url{https://fr.wikipedia.org/wiki/Lewis_Carroll})
\item Lecture : Donald E. Knuth, \textit{The Art of Computer Programming}, vol. 3, \S{} 5.3.3 (sélection par comparaisons)
\end{refs}

\bigskip

\begin{problem}[{Le treize tombe-t-il souvent un vendredi ? --- Collège}]
\textbf{PUZ-114.} \textbf{Problème.} On travaille dans le calendrier grégorien, celui qui est en usage
aujourd'hui : une année est bissextile si elle est multiple de $ 4 $, sauf si elle est
multiple de $ 100 $, sauf si elle est multiple de $ 400 $.
\begin{enumerate}
\item[(a)] Démontrez qu'un mois contient un vendredi $ 13 $ si et seulement si son premier jour est
un dimanche.
\item[(b)] Démontrez que toute année, commune ou bissextile, contient au moins un vendredi $ 13 $.
Déterminez ensuite le nombre maximal de vendredis $ 13 $ que peut compter une même année, et
démontrez votre réponse.
\item[(c)] Démontrez que $ 400 $ années grégoriennes comptent un nombre entier de semaines.
Déduisez-en que le calendrier se répète à l'identique tous les $ 400 $ ans, et donc que la
question \og{} quel jour de la semaine le $ 13 $ tombe-t-il le plus souvent ? \fg{} a bien un
sens.
\item[(d)] Répondez-y : déterminez, avec démonstration, le jour de la semaine le plus fréquent pour un
$ 13 $, et le moins fréquent. Le comptage est long mais fini --- organisez-le pour qu'il tienne
en une page.
\end{enumerate}

\end{problem}

\noindent La peur du vendredi $ 13 $ est récente : elle ne se répand en Europe et en
Amérique qu'au dix-neuvième siècle, en soudant deux superstitions plus anciennes et
indépendantes, celle du vendredi et celle du nombre treize. Le fait mathématique, lui, est
attesté depuis les années $ 1930 $, quand quelqu'un a pris la peine de compter : dans le
cycle grégorien de $ 400 $ ans, le $ 13 $ tombe un vendredi un peu plus souvent que
n'importe quel autre jour de la semaine. Rien dans la réforme de $ 1582 $ ne visait cet
effet ; il résulte du hasard heureux --- la question (c) --- qui fait que la longueur du cycle
grégorien est un multiple de sept, ce qui n'était nullement obligé et qui rend toute la
question calculable à la main. C'est un bel exemple de problème où l'arithmétique modulaire
sert à remplacer un calcul infini par un tableau fini.

\begin{refs}
\item Contexte : Vendredi treize --- Wikipédia (\url{https://fr.wikipedia.org/wiki/Vendredi_treize})
\item Contexte : Calendrier grégorien --- Wikipédia (\url{https://fr.wikipedia.org/wiki/Calendrier_gr%C3%A9gorien})
\item Contexte : Congruence sur les entiers --- Wikipédia (\url{https://fr.wikipedia.org/wiki/Congruence_sur_les_entiers})
\end{refs}

\bigskip

\begin{problem}[{Les allumettes à déplacer --- Collège}]
\textbf{PUZ-115.} \textbf{Problème.} On écrit des nombres en chiffres romains avec des allumettes : le
$ \mathrm{I} $ coûte une allumette, le $ \mathrm{V} $ et le $ \mathrm{X} $ en coûtent
deux chacun, le signe $ + $ en coûte deux, le signe $ - $ une seule, et le signe $ = $
deux. Sur la table on lit l'égalité fausse
\[ \mathrm{XI} + \mathrm{I} = \mathrm{X}. \]
Un coup consiste à retirer une allumette de la figure et à la reposer ailleurs : le nombre
total d'allumettes ne change donc pas, il est interdit d'en ajouter, d'en retirer ou d'en
laisser une qui ne serve à rien, et ce qui reste sur la table doit être une égalité bien
formée --- un nombre, un signe d'opération, un nombre, le signe $ = $, un nombre.
\begin{enumerate}
\item[(a)] Rendez l'égalité vraie en déplaçant une seule allumette.
\item[(b)] Dressez la liste de toutes les égalités vraies que l'on peut obtenir à partir de cette
figure en un seul déplacement, et démontrez que votre liste est complète. C'est ici que
l'invariant \og{} le nombre d'allumettes ne change pas \fg{} fait le gros du travail, en réduisant une
recherche vague à une vérification finie.
\item[(c)] On n'utilise plus que les chiffres $ \mathrm{I} $, $ \mathrm{V} $, $ \mathrm{X} $,
écrits selon les règles habituelles de la numération romaine, ce qui donne les nombres de
$ 1 $ à $ 39 $. Déterminez le plus grand nombre que l'on puisse écrire avec six
allumettes, et démontrez qu'aucun autre ne fait mieux.
\item[(d)] Même question avec sept, puis huit, puis neuf allumettes. Déterminez à partir de combien
d'allumettes on peut écrire le plus grand nombre disponible, et démontrez-le.
\end{enumerate}

\end{problem}

\noindent Les casse-tête d'allumettes envahissent la presse populaire européenne à la
fin du dix-neuvième siècle, à l'époque où l'allumette devient l'objet le plus banal d'une
poche, et Henry Dudeney comme Sam Loyd en publient par centaines. On les prend pour de simples
jeux d'observation, à tort : la méthode qui les résout proprement est celle-là même qu'on
emploie en combinatoire sérieuse --- d'abord un invariant qui élimine d'un coup l'immense
majorité des figures imaginables, ensuite un examen exhaustif du petit nombre de cas qui
survivent. La question (c) ajoute une optimisation authentique, où l'on doit non seulement
exhiber une écriture économique mais démontrer qu'aucune autre ne fait mieux, en majorant ce
qu'une allumette peut rapporter.

\begin{refs}
\item Contexte : Numération romaine --- Wikipédia (\url{https://fr.wikipedia.org/wiki/Num%C3%A9ration_romaine})
\item Contexte : Casse-tête d'allumettes --- Wikipédia (en anglais) (\url{https://en.wikipedia.org/wiki/Matchstick_puzzle})
\item Contexte : Sam Loyd --- Wikipédia (\url{https://fr.wikipedia.org/wiki/Sam_Loyd})
\end{refs}

\bigskip

\begin{problem}[{La côte et la vitesse moyenne --- Collège, short}]
\textbf{PUZ-116.} \textbf{Problème.} Vous montez une côte de deux kilomètres à la vitesse
moyenne de $ 15 $ km/h. À quelle vitesse devez-vous la redescendre pour que votre vitesse
moyenne sur l'aller-retour atteigne $ 30 $ km/h --- et si aucune vitesse ne convient,
démontrez-le.
\end{problem}

\noindent Le psychologue Max Wertheimer, fondateur de la théorie de la forme, aimait
soumettre des énigmes à ses amis pour observer comment ils pensaient ; il posa celle-ci à
Albert Einstein, et l'anecdote, qu'il rapporte dans \textit{Productive Thinking}, veut
qu'Einstein soit d'abord tombé dans le piège. Le piège est le mot \og{} moyenne \fg{} : une vitesse
moyenne n'est pas la moyenne des vitesses, c'est la distance totale divisée par le temps
total, et ces deux quantités ne se comportent pas du tout de la même façon quand on les
combine. Une fois l'énoncé traduit en temps plutôt qu'en vitesses, tout devient limpide --- y
compris le fait qu'une réponse puisse ne pas exister, ce qui est une conclusion parfaitement
respectable et demande, elle aussi, une démonstration.

\begin{refs}
\item Contexte : Moyenne harmonique --- Wikipédia (\url{https://fr.wikipedia.org/wiki/Moyenne_harmonique})
\item Contexte : Max Wertheimer --- Wikipédia (\url{https://fr.wikipedia.org/wiki/Max_Wertheimer})
\end{refs}

\bigskip


\section*{L'art de la démonstration}

\begin{problem}[{\og{} J'ai essayé avec 3, 5 et 7, donc c'est vrai \fg{} --- Collège}]
\textbf{PRF-64.} \textbf{Problème.} Un camarade affirme : \og{} J'ai essayé avec 3, 5 et 7, ça marche à chaque
fois, donc c'est vrai. \fg{} Ce problème demande ce qui manque exactement à cette phrase.
\begin{enumerate}
\item[(a)] Calculez $ n^2 + n + 11 $ pour n = 0, 1, 2, \dots{}, 9, et constatez que vous obtenez chaque
fois un nombre premier. Décidez ensuite si l'affirmation \og{} $ n^2 + n + 11 $ est un nombre
premier pour tout entier naturel n \fg{} est vraie, et démontrez votre réponse.
\item[(b)] Placez 2 points sur un cercle et tracez la corde qui les joint : le disque est coupé en
2 morceaux. Recommencez avec 3, puis 4, puis 5, puis 6 points, en traçant toutes les cordes
et en plaçant les points de sorte que jamais trois cordes ne se croisent au même endroit ;
comptez les morceaux à chaque fois. Écrivez la suite de nombres obtenue, dites quelle règle
elle vous suggère, puis dites si vos comptages la démontrent.
\item[(c)] Un seul contre-exemple suffit-il à réfuter une affirmation qui commence par \og{} pour
tout \fg{} ? Suffit-il à réfuter une affirmation qui commence par \og{} il existe \fg{} ? Justifiez les
deux réponses.
\item[(d)] Rédigez en trois ou quatre phrases ce qu'il faudrait ajouter à \og{} j'ai essayé avec 3, 5
et 7 \fg{} pour en faire une démonstration --- et pourquoi mille essais de plus n'y suffiraient
toujours pas.
\end{enumerate}

\end{problem}

\noindent Les deux pièges de ce problème sont célèbres et ils ont piégé de vrais
mathématiciens. Euler remarqua en 1772 que $ n^2 + n + 41 $ donne un nombre premier pour
quarante valeurs consécutives de n avant de céder ; 11 appartient à la même petite famille
de nombres, dits nombres chanceux d'Euler, et c'est pour cela qu'il tient si longtemps.
Le compte des morceaux du disque est le piège inverse : la suite commence si franchement
qu'on croit la connaître, et le sixième terme n'est pas celui qu'on annonce --- Leo Moser en
a fait un classique de la mise en garde. La leçon est la même dans les deux cas et elle est
au fondement de tout ce site : vérifier n'est pas démontrer, l'un donne de la confiance,
l'autre donne la certitude, et rien ne transforme la première en la seconde.

\begin{refs}
\item Contexte : Contre-exemple --- Wikipédia (\url{https://fr.wikipedia.org/wiki/Contre-exemple})
\item Contexte : Nombre chanceux d'Euler --- Wikipédia (\url{https://fr.wikipedia.org/wiki/Nombre_chanceux_d%27Euler})
\item Contexte : Partage d'un disque par des cordes --- Wikipédia (en anglais) (\url{https://en.wikipedia.org/wiki/Dividing_a_circle_into_areas})
\end{refs}

\bigskip

\begin{problem}[{Deux impairs font un pair --- démontrez-le --- Collège}]
\textbf{PRF-65.} \textbf{Problème.} Un entier est pair s'il s'écrit $ 2k $ pour un certain entier k, et
impair s'il s'écrit $ 2k+1 $ pour un certain entier k. Tout ce qui suit ne doit se servir
que de ces deux écritures.
\begin{enumerate}
\item[(a)] Vérifiez sur cinq couples de votre choix que la somme de deux entiers impairs est
paire, puis dites en une phrase pourquoi ces cinq vérifications ne démontrent rien.
\item[(b)] Démontrez que la somme de deux entiers impairs est paire. Votre argument doit valoir
pour tous les couples d'entiers impairs d'un seul coup ; montrez du doigt l'endroit précis
où il cesse de parler d'exemples.
\item[(c)] Justifiez le point sur lequel repose (b) et qu'on oublie toujours de justifier : tout
entier impair s'écrit bien $ 2k+1 $ pour un certain entier k, et tout nombre de cette
forme est bien impair.
\item[(d)] Démontrez de la même façon que le produit de deux entiers impairs est impair, que la
somme de trois entiers impairs est impaire, et que la somme de trois entiers consécutifs est
divisible par 3. Décidez enfin si la somme de quatre entiers consécutifs est divisible
par 4.
\end{enumerate}

\end{problem}

\noindent C'est presque toujours la première vraie démonstration qu'un élève écrit,
et elle contient déjà tout le mécanisme : une lettre remplace un nombre, et la ligne de
calcul qui suit parle de tous les cas à la fois. Les Grecs travaillaient ces énoncés bien
avant qu'on sache écrire une lettre pour un nombre --- le livre IX des \textit{Éléments}
d'Euclide démontre en propositions séparées que la somme de nombres pairs est paire, que la
somme d'un nombre pair de nombres impairs est paire, et ainsi de suite, entièrement en
langue et en longueurs. Comparer leur rédaction à la vôtre est instructif : le calcul
littéral n'a pas rendu la démonstration possible, il l'a rendue courte. PRF-74 vous fera
revenir sur le texte que vous aurez écrit ici.

\begin{refs}
\item Contexte : Parité (arithmétique) --- Wikipédia (\url{https://fr.wikipedia.org/wiki/Parit%C3%A9_(arithm%C3%A9tique)})
\item Contexte : Éléments d'Euclide --- Wikipédia (\url{https://fr.wikipedia.org/wiki/%C3%89l%C3%A9ments_d%27Euclide})
\item Contexte : Calcul algébrique --- Wikipédia (\url{https://fr.wikipedia.org/wiki/Calcul_alg%C3%A9brique})
\end{refs}

\bigskip

\begin{problem}[{Un nombre et son double --- Collège, short}]
\textbf{PRF-66.} \textbf{Problème.} Démontrez ou réfutez : un entier et son double n'ont
jamais la même parité. Si l'affirmation est fausse, énoncez à sa place la règle exacte qui
dit quand un entier n et son double $ 2n $ ont la même parité, et démontrez cette
règle.
\end{problem}

\noindent \og{} Démontrez ou réfutez \fg{} est un genre à part : la moitié du travail
consiste à décider de quel côté on se range, et la première chose à faire est d'essayer de
casser l'énoncé plutôt que de le croire. Beaucoup d'affirmations fausses en mathématiques
sont presque vraies --- vraies sur la moitié des cas, ou sur tous sauf un --- et c'est
précisément ce qui les rend dangereuses. PRF-38 pose la même question au lycée, sur les
nombres irrationnels.

\begin{refs}
\item Contexte : Parité (arithmétique) --- Wikipédia (\url{https://fr.wikipedia.org/wiki/Parit%C3%A9_(arithm%C3%A9tique)})
\item Contexte : Contre-exemple --- Wikipédia (\url{https://fr.wikipedia.org/wiki/Contre-exemple})
\end{refs}

\bigskip

\begin{problem}[{Où se cache l'erreur : une \og{} démonstration \fg{} que 1 = 2 --- Collège, short}]
\textbf{PRF-67.} \textbf{Problème.} On part de deux nombres égaux, $ a = b $, et on
enchaîne :
\[ a^2 = ab, \qquad a^2 - b^2 = ab - b^2, \qquad (a-b)(a+b) = b(a-b), \qquad a+b = b, \qquad 2b = b, \qquad 2 = 1. \]
Désignez la ligne exactement fausse --- pas \og{} quelque part par là \fg{}, la ligne --- et écrivez en
une phrase la règle qu'elle enfreint.
\end{problem}

\noindent Toutes les lignes sont correctement calculées, et pourtant la conclusion
est fausse : c'est le meilleur argument qui soit contre l'idée qu'une démonstration serait
un calcul juste. Une démonstration est un enchaînement dont chaque pas est \textit{autorisé},
et une seule autorisation manquante suffit à tout ruiner. Ce faux raisonnement circule
depuis des siècles et il a un cousin géométrique tout aussi instructif, le paradoxe du carré
manquant, où un découpage semble faire apparaître une case de plus : là, ce qui triche n'est
pas un calcul mais un dessin. Retenez la morale des deux --- une figure et une ligne de calcul
sont des indices, jamais des preuves.

\begin{refs}
\item Contexte : Division par zéro --- Wikipédia (\url{https://fr.wikipedia.org/wiki/Division_par_z%C3%A9ro})
\item Contexte : Raisonnement fallacieux --- Wikipédia (\url{https://fr.wikipedia.org/wiki/Raisonnement_fallacieux})
\item Contexte : Paradoxe du carré manquant --- Wikipédia (\url{https://fr.wikipedia.org/wiki/Paradoxe_du_carr%C3%A9_manquant})
\end{refs}

\bigskip

\begin{problem}[{Le principe des tiroirs, raconté puis employé --- Collège}]
\textbf{PRF-68.} \textbf{Problème.} Le principe des tiroirs dit ceci : si l'on range plus d'objets qu'on n'a
de tiroirs, alors un tiroir au moins contient au moins deux objets. Il tient en une ligne et
il a l'air de ne rien dire --- pourtant tout ce qui suit en découle.
\begin{enumerate}
\item[(a)] Démontrez le principe lui-même, en deux phrases, par l'absurde : supposez que chaque
tiroir contienne au plus un objet, et comptez.
\item[(b)] Démontrez que parmi 13 personnes quelconques, deux au moins sont nées le même mois.
\item[(c)] On choisit six nombres différents parmi 1, 2, \dots{}, 10. Démontrez que deux d'entre eux ont
pour somme 11.
\item[(d)] Une tête humaine porte au plus 150 000 cheveux. Démontrez que dans une ville de
200 000 habitants, deux personnes au moins ont exactement le même nombre de cheveux.
\item[(e)] Voici le véritable exercice. Pour chacune des questions (b), (c) et (d), écrivez
explicitement ce qui jouait le rôle des objets, ce qui jouait celui des tiroirs, et combien
il y en avait de chaque --- puis dites, en une phrase par question, lequel des deux était le
plus difficile à voir.
\end{enumerate}

\end{problem}

\noindent Le principe porte le nom de Dirichlet, qui s'en est servi comme d'un outil
en 1834 sous le nom de \textit{Schubfachprinzip}, \og{} principe des tiroirs \fg{}, mais l'exemple des
cheveux est bien plus ancien : un recueil français de 1622, la \textit{Récréation
mathématique} publiée sous le nom de Jean Leurechon, remarque déjà qu'il est nécessaire
que deux hommes aient le même nombre de cheveux. Ce qui frappe, c'est le contraste entre la
banalité de l'énoncé et la force de ses conséquences : il produit des affirmations qu'aucune
observation ne pourrait établir, sur des gens qu'on n'a jamais vus. Toute la difficulté, à
chaque emploi, tient dans la question \og{} quels sont les tiroirs ? \fg{} --- car ils ne sont jamais
donnés. PRF-07 et PRF-55 reprennent la question au lycée, avec des tiroirs qu'il faut
fabriquer soi-même.

\begin{refs}
\item Contexte : Principe des tiroirs --- Wikipédia (\url{https://fr.wikipedia.org/wiki/Principe_des_tiroirs})
\item Contexte : Peter Gustav Lejeune Dirichlet --- Wikipédia (\url{https://fr.wikipedia.org/wiki/Peter_Gustav_Lejeune_Dirichlet})
\item Lecture : A. Engel, \textit{Problem-Solving Strategies}, ch. 4 (le principe des tiroirs)
\end{refs}

\bigskip

\begin{problem}[{Deux personnes ont le même nombre d'amis --- Collège, short}]
\textbf{PRF-69.} \textbf{Problème.} Dans un groupe d'au moins deux personnes, où l'amitié
est toujours réciproque, démontrez qu'il existe deux personnes ayant exactement le même
nombre d'amis à l'intérieur du groupe. Rédigez-le comme un raisonnement par l'absurde :
supposez que ces nombres soient tous différents, et poussez cette supposition jusqu'à une
impossibilité.
\end{problem}

\noindent Le raisonnement par l'absurde est la manœuvre la plus étrange de la
logique --- on affirme le contraire de ce qu'on veut établir, on travaille consciencieusement
sur une hypothèse dont on espère qu'elle est fausse, et l'effondrement final est la
démonstration. Ce problème-ci en est une version où l'on ne calcule rien du tout : il n'y a
ni formule, ni figure, et l'énoncé ne dit même pas combien de personnes il y a. Traduit en
langage de graphes --- un point par personne, un trait par amitié --- c'est l'un des premiers
résultats qu'on rencontre dans la théorie née de l'article d'Euler sur les ponts de
Königsberg (1736).

\begin{refs}
\item Contexte : Raisonnement par l'absurde --- Wikipédia (\url{https://fr.wikipedia.org/wiki/Raisonnement_par_l%27absurde})
\item Contexte : Théorie des graphes --- Wikipédia (\url{https://fr.wikipedia.org/wiki/Th%C3%A9orie_des_graphes})
\item Contexte : Principe des tiroirs --- Wikipédia (\url{https://fr.wikipedia.org/wiki/Principe_des_tiroirs})
\end{refs}

\bigskip

\begin{problem}[{Ce qui ne change jamais --- Collège}]
\textbf{PRF-70.} \textbf{Problème.} Un invariant est une quantité qui ne bouge pas alors que la situation,
elle, bouge. Chacune des questions ci-dessous demande de démontrer qu'une chose est
\textit{impossible} --- et aucun essai raté, si nombreux soient-ils, ne peut établir cela.
\begin{enumerate}
\item[(a)] Sept verres sont posés à l'envers sur une table. Un coup consiste à retourner exactement
deux verres, lesquels que vous voulez. Démontrez qu'aucune suite de coups ne mettra les sept
verres à l'endroit.
\item[(b)] On écrit au tableau les nombres 1, 2, \dots{}, 10. Un coup consiste à effacer deux nombres et
à écrire à la place leur différence positive (ou 0 s'ils sont égaux). Après neuf coups il ne
reste qu'un nombre : démontrez que sa parité est fixée d'avance, quels qu'aient été vos
choix.
\item[(c)] Un jeton est posé sur une case d'un quadrillage colorié comme un damier. À chaque coup
il se déplace en diagonale, d'une case exactement. Démontrez qu'il ne pourra jamais atteindre
une case de l'autre couleur, et dites ce que cela raconte du fou aux échecs.
\item[(d)] Pour chacune des trois questions, nommez l'invariant en une phrase. Expliquez ensuite
pourquoi exhiber un invariant démontre une impossibilité, alors que cent tentatives ratées
ne démontrent rien du tout.
\end{enumerate}

\end{problem}

\noindent L'argument d'invariance est peut-être la plus belle idée élémentaire des
mathématiques : au lieu de suivre un processus qui part dans toutes les directions, on
trouve la seule chose qu'il ne peut pas modifier, et l'on regarde si le but demandé la
respecte. S'il ne la respecte pas, le but est hors d'atteinte --- définitivement, pour toutes
les suites de coups possibles, y compris celles que personne n'essaiera jamais. C'est la
première technique du livre d'Arthur Engel pour les olympiades, et ce n'est pas un hasard :
c'est aussi la manière dont on démontre qu'un programme informatique fait bien ce qu'il
prétend (DCS-01 en donne la version échiquier). Le mot \og{} impossible \fg{} est l'un des rares que
les mathématiques prononcent sans trembler ; l'invariant est ce qui leur en donne le
droit.

\begin{refs}
\item Contexte : Invariant --- Wikipédia (\url{https://fr.wikipedia.org/wiki/Invariant})
\item Lecture : A. Engel, \textit{Problem-Solving Strategies}, ch. 1 (le principe d'invariance)
\item Voir aussi : DCS-01, l'échiquier mutilé (\url{applied-discrete-cs.html#DCS-01})
\end{refs}

\bigskip

\begin{problem}[{Une case en moins, et le coloriage qui décide --- Collège}]
\textbf{PRF-71.} \textbf{Problème.} On colorie un carré $ 5 \times 5 $ comme un damier, les quatre coins
recevant la même couleur. On retire une case, puis on essaie de recouvrir les 24 cases
restantes par 12 dominos, chaque domino couvrant exactement deux cases voisines par un
côté.
\begin{enumerate}
\item[(a)] Comptez les cases de chaque couleur. Démontrez ensuite qu'aucun pavage n'existe dès que
la case retirée est de la couleur la moins nombreuse.
\item[(b)] Cherchez maintenant, parmi les cases de l'autre couleur, celles pour lesquelles un
pavage existe réellement. Chaque réponse demande sa justification : une construction dans un
sens, une démonstration d'impossibilité dans l'autre.
\item[(c)] Le coloriage ne figure nulle part dans l'énoncé du problème : c'est vous qui l'ajoutez.
Expliquez en quelques phrases ce qu'il rend visible et que le quadrillage nu ne montrait
pas.
\item[(d)] Reprenez (a) en supposant qu'un domino ait aussi le droit de couvrir deux cases voisines
\textit{en diagonale}. Dites ce que devient l'argument, et à quelle ligne exactement il
cesse de fonctionner.
\end{enumerate}

\end{problem}

\noindent Colorier une figure pour démontrer qu'un découpage est impossible est un
tour de main que Solomon Golomb a installé au cœur d'un domaine entier lorsqu'il a inventé
le mot \og{} polyomino \fg{} en 1954 et publié ses premiers résultats de pavage. La force du procédé
tient à ce qu'il ajoute au problème une information qui n'y était pas : le damier n'existe
que dans votre tête, et pourtant c'est lui qui tranche. Remarquez aussi ce que la question
(b) demande de plus que la (a) --- une impossibilité se démontre par un argument, mais une
possibilité se démontre en montrant la chose, et les deux moitiés d'un \og{} pour quelles cases
est-ce faisable ? \fg{} n'ont jamais la même allure. Le résultat frère sur l'échiquier
$ 8 \times 8 $ privé de deux cases de couleurs différentes est dû à Ralph Gomory, et son
argument --- d'une élégance rare --- a été popularisé par Ross Honsberger en 1973.

\begin{refs}
\item Contexte : Polyomino --- Wikipédia (\url{https://fr.wikipedia.org/wiki/Polyomino})
\item Contexte : Problème de l'échiquier mutilé --- Wikipédia (\url{https://fr.wikipedia.org/wiki/Probl%C3%A8me_de_l%27%C3%A9chiquier_mutil%C3%A9})
\item Lecture : S. W. Golomb, \textit{Polyominoes} (2ᵉ éd., 1994), ch. 1
\end{refs}

\bigskip

\begin{problem}[{La réciproque n'est pas l'énoncé --- Collège}]
\textbf{PRF-72.} \textbf{Problème.} La réciproque de \og{} si P alors Q \fg{} est \og{} si Q alors P \fg{}. Ce sont deux
affirmations différentes : démontrer l'une ne démontre pas l'autre, et il arrive
constamment que l'une soit vraie et l'autre fausse.
\begin{enumerate}
\item[(a)] Pour chacun des énoncés suivants, écrivez la réciproque, puis décidez si cette
réciproque est vraie --- par une démonstration si oui, par un contre-exemple si non. (i) Si un
entier se termine par 0, alors il est divisible par 5. (ii) Si un triangle est équilatéral,
alors il est isocèle. (iii) Si un entier est divisible par 6, alors il est divisible par 2
et par 3. (iv) S'il pleut, le sol est mouillé.
\item[(b)] Ceux des énoncés dont la réciproque est vraie méritent la forme \og{} si et seulement si \fg{}.
Réécrivez-les ainsi, et rédigez pour l'un d'eux les deux sens comme deux démonstrations
clairement séparées et étiquetées.
\item[(c)] Un critère de divisibilité dit : si la somme des chiffres d'un entier est divisible par
3, alors l'entier l'est aussi. Écrivez sa réciproque, et dites laquelle des deux vous
employez réellement quand vous testez 1 234 567.
\item[(d)] Expliquez en trois phrases pourquoi \og{} la réciproque est évidente \fg{} est l'une des phrases
les plus dangereuses qu'on puisse écrire dans une copie de mathématiques.
\end{enumerate}

\end{problem}

\noindent Confondre un énoncé et sa réciproque est l'erreur de raisonnement la plus
répandue, et pas seulement en mathématiques : \og{} tous les menteurs se contredisent \fg{} ne dit
rien de celui qui vient de se contredire, et un tribunal qui l'oublie condamne un innocent.
Aristote la traitait déjà comme une faute nommée, et la logique médiévale l'appelait
l'affirmation du conséquent. En classe, l'occasion la plus fréquente de rencontrer la
distinction est géométrique : le théorème de Thalès sert à calculer des longueurs, sa
réciproque sert à démontrer que des droites sont parallèles, et ce sont deux outils
différents (PRF-73). PRF-46 reprend la question au lycée, sur des équivalences plus
coriaces.

\begin{refs}
\item Contexte : Implication réciproque --- Wikipédia (\url{https://fr.wikipedia.org/wiki/Implication_r%C3%A9ciproque})
\item Contexte : Implication (logique) --- Wikipédia (\url{https://fr.wikipedia.org/wiki/Implication_(logique)})
\item Contexte : Critères de divisibilité --- Wikipédia (\url{https://fr.wikipedia.org/wiki/Crit%C3%A8res_de_divisibilit%C3%A9})
\end{refs}

\bigskip

\begin{problem}[{Thalès, et la réciproque de Thalès --- Collège}]
\textbf{PRF-73.} \textbf{Problème.} Deux droites se coupent en A. Sur la première on place B, puis M, du même
côté de A ; sur la seconde on place C, puis N, du même côté de A. Le théorème de Thalès
affirme que si les droites $ (BC) $ et $ (MN) $ sont parallèles, alors
\[ \frac{AB}{AM} \;=\; \frac{AC}{AN} \;=\; \frac{BC}{MN}. \]
\begin{enumerate}
\item[(a)] On donne $ AB = 3 $, $ AM = 5 $, $ AC = 6 $, et $ (BC) $ parallèle à
$ (MN) $. Calculez AN en écrivant, à côté de chaque égalité, l'hypothèse qui l'autorise.
Une seule de vos lignes se sert du parallélisme : entourez-la.
\item[(b)] La réciproque du théorème sert à démontrer qu'une chose est parallèle. Énoncez-la
précisément, hypothèses comprises, puis servez-vous-en pour démontrer que deux droites d'une
figure de votre construction sont parallèles.
\item[(c)] La réciproque exige que les points soient placés \og{} dans le même ordre \fg{} sur les deux
droites. Construisez une figure où les rapports de longueurs sont bien égaux mais où cette
condition n'est pas remplie, et décidez, avec justification, si les droites y sont
parallèles.
\item[(d)] Dites en deux phrases ce que (c) enseigne : ce qui arrive à une démonstration quand on
applique un théorème sans vérifier l'une de ses hypothèses.
\end{enumerate}

\end{problem}

\noindent La tradition fait de Thalès de Milet, vers 600 av. J.-C., le premier
à avoir \textit{démontré} plutôt que constaté, et la plus jolie des légendes le montre mesurant
la hauteur de la grande pyramide par la longueur de son ombre --- c'est-à-dire employant très
exactement le théorème qui porte aujourd'hui son nom. Une curiosité mérite d'être connue :
l'appellation \og{} théorème de Thalès \fg{} pour cet énoncé de proportions est une convention
française, quand l'anglais l'appelle le théorème des intercepts et réserve le nom de Thalès
à un tout autre résultat, celui de l'angle inscrit dans un demi-cercle. Ce théorème est
souvent le premier dont un élève français utilise la réciproque comme un outil de
démonstration à part entière, et il vaut la peine de bien voir que les deux sens ne servent
pas au même travail : l'un mesure, l'autre prouve.

\begin{refs}
\item Contexte : Théorème de Thalès --- Wikipédia (\url{https://fr.wikipedia.org/wiki/Th%C3%A9or%C3%A8me_de_Thal%C3%A8s})
\item Contexte : Thalès --- Wikipédia (\url{https://fr.wikipedia.org/wiki/Thal%C3%A8s})
\item Contexte : Éléments d'Euclide --- Wikipédia (\url{https://fr.wikipedia.org/wiki/%C3%89l%C3%A9ments_d%27Euclide})
\end{refs}

\bigskip

\begin{problem}[{Rédiger : ce qu'un lecteur a le droit d'exiger --- Collège}]
\textbf{PRF-74.} \textbf{Problème.} Trouver un argument et l'écrire sont deux travaux distincts, et le second
s'apprend. Reprenez la démonstration que vous avez trouvée en PRF-65 (b) --- la somme de deux
entiers impairs est paire --- et travaillez-la comme on travaille un texte.
\begin{enumerate}
\item[(a)] Rédigez-la pour un camarade absent ce jour-là : chaque fait employé doit être nommé, et
aucune étape ne doit lui demander de deviner. Comptez ensuite vos phrases.
\item[(b)] Soulignez les mots qui portent la logique --- \og{} soit \fg{}, \og{} supposons \fg{}, \og{} donc \fg{}, \og{} or \fg{},
\og{} il existe \fg{}, \og{} pour tout \fg{} --- et vérifiez qu'aucune phrase n'en est privée par accident.
Une phrase sans mot de liaison est souvent une phrase qui ne démontre rien.
\item[(c)] Donnez votre texte à quelqu'un qui n'a pas cherché le problème, en lui demandant de
s'arrêter à la première ligne qu'il ne croit pas. Réécrivez à partir de cette ligne, puis
recommencez l'opération avec un autre lecteur.
\item[(d)] Faites le même travail sur votre démonstration de PRF-68 (c), qui est d'une autre
nature : une phrase y remplace un calcul. Dites enfin, en deux phrases, ce que change le
passage de \og{} je vois pourquoi c'est vrai \fg{} à \og{} j'ai écrit pourquoi c'est vrai \fg{}.
\end{enumerate}

\end{problem}

\noindent La quatrième et dernière étape de la méthode de George Pólya, dans
\textit{How to Solve It} (1945), traduit en français sous le titre \textit{Comment poser et
résoudre un problème}, s'appelle \og{} revenir en arrière \fg{} : le problème n'est pas fini quand
la réponse est trouvée, il est fini quand l'argument tient debout tout seul, sans son auteur
à côté pour l'expliquer. C'est exactement le service que rend, dans la vie d'un
mathématicien, le rapporteur d'un article : quelqu'un qui n'est pas dans votre tête et qui
s'arrête à la première ligne qu'il ne croit pas. La question (c) vous fait fabriquer ce
lecteur, et c'est de loin l'exercice le plus utile de cette page --- parce qu'une démonstration
qu'on est seul à comprendre n'a encore convaincu personne.

\begin{refs}
\item Contexte : Démonstration (logique et mathématiques) --- Wikipédia (\url{https://fr.wikipedia.org/wiki/D%C3%A9monstration_(logique_et_math%C3%A9matiques)})
\item Contexte : George Pólya --- Wikipédia (\url{https://fr.wikipedia.org/wiki/George_P%C3%B3lya})
\item Lecture : G. Pólya, \textit{Comment poser et résoudre un problème} (trad. fr. de \textit{How to Solve It})
\end{refs}

\bigskip

\begin{problem}[{Toujours vrai, ou seulement souvent vrai --- Collège, short}]
\textbf{PRF-75.} \textbf{Problème.} Décidez, avec démonstration, si l'affirmation \og{} la
somme de deux nombres premiers est un nombre pair \fg{} est toujours vraie ; si elle ne l'est
pas, énoncez et démontrez la règle exacte qui dit quand elle l'est. Dites ensuite en une
phrase ce que change, pour un mathématicien, le passage de \og{} je n'ai trouvé aucune
exception \fg{} à \og{} il ne peut pas y en avoir \fg{}.
\end{problem}

\noindent La distinction entre \og{} toujours vrai \fg{} et \og{} vrai dans tous les cas que
j'ai regardés \fg{} est la frontière même des mathématiques, et il existe un exemple magnifique
juste à côté de celui-ci : la conjecture de Goldbach affirme que tout entier pair supérieur
à 2 est la somme de deux nombres premiers. Elle a été vérifiée par ordinateur pour tous les
entiers pairs jusqu'à plus de quatre milliards de milliards, sans la moindre exception, et
personne ne l'a démontrée depuis qu'Euler et Goldbach en ont discuté par lettres en 1742.
Voilà à quoi ressemble une affirmation \og{} souvent vraie \fg{} : indéfiniment vérifiée, toujours
non démontrée, et donc toujours pas un théorème.

\begin{refs}
\item Contexte : Nombre premier --- Wikipédia (\url{https://fr.wikipedia.org/wiki/Nombre_premier})
\item Contexte : Conjecture de Goldbach --- Wikipédia (\url{https://fr.wikipedia.org/wiki/Conjecture_de_Goldbach})
\item Voir aussi : Problèmes ouverts (\url{open-problems.html})
\end{refs}

\bigskip


\vfill
\noindent\emph{Hors de la Caverne} --- des probl\`emes, pas de r\'eponses.
\end{document}
